
\Data\ID{1003112001}
\Updated{2020-07-15 03:07:12}
\Level{A}
\SubArea{Derivative}
\LANGen{
\Question{
Find the derivative of the function \( f(x)=2x^3+3x^2-1 \) at the point \( x_0=-2 \).}
{\( 12 \)}{\( 11 \)}{\( 36 \)}{\( -5 \)}{\( -12 \)}{}}
\LANGpl{
\Question{
Wyznacz pochodną funkcji  \( f(x)=2x^3+3x^2-1 \) w punkcie \( x_0=-2 \).}
{\( 12 \)}{\( 11 \)}{\( 36 \)}{\( -5 \)}{\( -12 \)}{}}
\LANGcs{
\Question{
Vypočítejte derivaci funkce \( f(x)=2x^3+3x^2-1 \) v bodě \( x_0=-2 \).}
{\( 12 \)}{\( 11 \)}{\( 36 \)}{\( -5 \)}{\( -12 \)}{}}
\LANGsk{
\Question{
Vypočítajte deriváciu funkcie \( f(x)=2x^3+3x^2-1 \) v bode \( x_0=-2 \).}
{\( 12 \)}{\( 11 \)}{\( 36 \)}{\( -5 \)}{\( -12 \)}{}}
\LANGes{
\Question{
Halla la derivada de la función \( f(x)=2x^3+3x^2-1 \) en el punto \( x_0=-2 \).}
{\( 12 \)}{\( 11 \)}{\( 36 \)}{\( -5 \)}{\( -12 \)}{}}
\End

\Data\ID{1003112002}
\Updated{2020-07-15 03:07:12}
\Level{A}
\SubArea{Derivative}
\LANGen{
\Question{
Given the function \( f(x)=x^4-3x^2+2x \), find \( f'(-1) \).}
{\( 4 \)}{\( 12 \)}{\( -8 \)}{\( -10 \)}{\( 2 \)}{}}
\LANGpl{
\Question{
Dana jest funkcja \( f(x)=x^4-3x^2+2x \), oblicz \( f'(-1) \).}
{\( 4 \)}{\( 12 \)}{\( -8 \)}{\( -10 \)}{\( 2 \)}{}}
\LANGcs{
\Question{
Je dána funkce \( f(x)=x^4-3x^2+2x \), vypočítejte \( f'(-1) \).}
{\( 4 \)}{\( 12 \)}{\( -8 \)}{\( -10 \)}{\( 2 \)}{}}
\LANGsk{
\Question{
Je daná funkcia \( f(x)=x^4-3x^2+2x \), vypočítajte \( f'(-1) \).}
{\( 4 \)}{\( 12 \)}{\( -8 \)}{\( -10 \)}{\( 2 \)}{}}
\LANGes{
\Question{
Dada la función \( f(x)=x^4-3x^2+2x \), calcula \( f'(-1) \).}
{\( 4 \)}{\( 12 \)}{\( -8 \)}{\( -10 \)}{\( 2 \)}{}}
\End

\Data\ID{1003112003}
\Updated{2020-07-15 03:07:12}
\Level{A}
\SubArea{Derivative}
\LANGen{
\Question{
Given the function \( g(x)=x^3-3x^2-2x+5 \), find \( g'(2) \).}
{\( -2 \)}{\( 0 \)}{\( 3 \)}{\( -3 \)}{\( 2 \)}{}}
\LANGpl{
\Question{
Dana jest funkcja  \( g(x)=x^3-3x^2-2x+5 \), oblicz \( g'(2) \).}
{\( -2 \)}{\( 0 \)}{\( 3 \)}{\( -3 \)}{\( 2 \)}{}}
\LANGcs{
\Question{
Je dána funkce \( g(x)=x^3-3x^2-2x+5 \), vypočítejte \( g'(2) \).}
{\( -2 \)}{\( 0 \)}{\( 3 \)}{\( -3 \)}{\( 2 \)}{}}
\LANGsk{
\Question{
Je daná funkcie \( g(x)=x^3-3x^2-2x+5 \), vypočítajte \( g'(2) \).}
{\( -2 \)}{\( 0 \)}{\( 3 \)}{\( -3 \)}{\( 2 \)}{}}
\LANGes{
\Question{
Dada la función \( g(x)=x^3-3x^2-2x+5 \), calcula \( g'(2) \).}
{\( -2 \)}{\( 0 \)}{\( 3 \)}{\( -3 \)}{\( 2 \)}{}}
\End

\Data\ID{1003112004}
\Updated{2020-07-15 03:07:12}
\Level{A}
\SubArea{Derivative}
\LANGen{
\Question{
Given the function \( f(x)=x^2-6x+5 \), determine all the values of \( x \) (\( x\in\mathbb{R} \)) for which \( f'(x)=0 \).}
{\( x=3 \)}{\( x=1 \)}{\( x=-3 \)}{\( x=5 \)}{\( x_1=1 \), \(x_2=5 \)}{}}
\LANGpl{
\Question{
Dana jest funkcja  \( f(x)=x^2-6x+5 \), wyznacz wszystkie wartości \( x \) (\( x\in\mathbb{R} \)) dla których \( f'(x)=0 \).}
{\( x=3 \)}{\( x=1 \)}{\( x=-3 \)}{\( x=5 \)}{\( x_1=1 \), \(x_2=5 \)}{}}
\LANGcs{
\Question{
Je dána funkce \( f(x)=x^2-6x+5 \). Určete všechna \( x\in\mathbb{R} \), pro která platí \( f'(x)=0 \).}
{\( x=3 \)}{\( x=1 \)}{\( x=-3 \)}{\( x=5 \)}{\( x_1=1 \), \(x_2=5 \)}{}}
\LANGsk{
\Question{
Je daná funkcia \( f(x)=x^2-6x+5 \). Určte všetky \( x\in\mathbb{R} \), pre ktoré platí \( f'(x)=0 \).}
{\( x=3 \)}{\( x=1 \)}{\( x=-3 \)}{\( x=5 \)}{\( x_1=1 \), \(x_2=5 \)}{}}
\LANGes{
\Question{
Dada la función \( f(x)=x^2-6x+5 \), determina todos los valores de \( x \) (\( x\in\mathbb{R} \)) para los cuales \( f'(x)=0 \).}
{\( x=3 \)}{\( x=1 \)}{\( x=-3 \)}{\( x=5 \)}{\( x_1=1 \), \(x_2=5 \)}{}}
\End

\Data\ID{1003112005}
\Updated{2020-07-15 03:07:12}
\Level{A}
\SubArea{Derivative}
\LANGen{
\Question{
Given the function \( f(x)=-x^2-4x \), determine all the values of \( x \) (\( x\in\mathbb{R} \)) for which \( f'(x)=0 \).}
{\( x=-2 \)}{\( x=2 \)}{\( x=-4 \)}{\( x=4 \)}{\( x_1=0 \), \( x_2=4 \)}{}}
\LANGpl{
\Question{
Dana jest funkcja  \( f(x)=-x^2-4x \), wyznacz wszystkie wartości  \( x \) (\( x\in\mathbb{R} \)) dla których \( f'(x)=0 \).}
{\( x=-2 \)}{\( x=2 \)}{\( x=-4 \)}{\( x=4 \)}{\( x_1=0 \), \( x_2=4 \)}{}}
\LANGcs{
\Question{
Je dána funkce \( f(x)=-x^2-4x \). Určete všechna \( x\in\mathbb{R} \), pro která platí \( f'(x)=0 \).}
{\( x=-2 \)}{\( x=2 \)}{\( x=-4 \)}{\( x=4 \)}{\( x_1=0 \), \( x_2=4 \)}{}}
\LANGsk{
\Question{
Je daná funkcia \( f(x)=-x^2-4x \). Určte všetky \( x\in\mathbb{R} \), pre ktoré platí \( f'(x)=0 \).}
{\( x=-2 \)}{\( x=2 \)}{\( x=-4 \)}{\( x=4 \)}{\( x_1=0 \), \( x_2=4 \)}{}}
\LANGes{
\Question{
Dada la función \( f(x)=-x^2-4x \), determina todos los valores de \( x \) (\( x\in\mathbb{R} \)) para los cuales \( f'(x)=0 \).}
{\( x=-2 \)}{\( x=2 \)}{\( x=-4 \)}{\( x=4 \)}{\( x_1=0 \), \( x_2=4 \)}{}}
\End

\Data\ID{1003112006}
\Updated{2020-07-15 03:07:12}
\Level{A}
\SubArea{Derivative}
\LANGen{
\Question{
Given the function \( f(x)=x^2+7 \), determine all the values of \( x \) (\( x\in\mathbb{R} \)) for which \(|f' (x)|=1 \).}
{\( x_1=\frac12 \), \( x_2 =-\frac12 \)}{\( x=\frac12 \)}{\( x_1=1 \), \( x_2 =-1 \)}{\( x=1 \)}{\( x=-\frac12 \)}{}}
\LANGpl{
\Question{
Dana jest funkcja  \( f(x)=x^2+7 \), wyznacz wszystkie wartości  \( x \) (\( x\in\mathbb{R} \)) dla których  \(|f' (x)|=1 \).}
{\( x_1=\frac12 \), \( x_2 =-\frac12 \)}{\( x=\frac12 \)}{\( x_1=1 \), \( x_2 =-1 \)}{\( x=1 \)}{\( x=-\frac12 \)}{}}
\LANGcs{
\Question{
Je dána funkce \( f(x)=x^2+7 \). Určete všechna \( x\in\mathbb{R} \), pro která platí \(|f' (x)|=1 \).}
{\( x_1=\frac12 \), \( x_2 =-\frac12 \)}{\( x=\frac12 \)}{\( x_1=1 \), \( x_2 =-1 \)}{\( x=1 \)}{\( x=-\frac12 \)}{}}
\LANGsk{
\Question{
Je daná funkcia \( f(x)=x^2+7 \). Určte všetky \( x\in\mathbb{R} \), pre ktoré platí \(|f' (x)|=1 \).}
{\( x_1=\frac12 \), \( x_2 =-\frac12 \)}{\( x=\frac12 \)}{\( x_1=1 \), \( x_2 =-1 \)}{\( x=1 \)}{\( x=-\frac12 \)}{}}
\LANGes{
\Question{
Dada la función \( f(x)=x^2+7 \), determina todos los valores de \( x \) (\( x\in\mathbb{R} \)) para los cuales \(|f' (x)|=1 \).}
{\( x_1=\frac12 \), \( x_2 =-\frac12 \)}{\( x=\frac12 \)}{\( x_1=1 \), \( x_2 =-1 \)}{\( x=1 \)}{\( x=-\frac12 \)}{}}
\End

\Data\ID{1003112007}
\Updated{2020-07-15 03:07:12}
\Level{A}
\SubArea{Derivative}
\LANGen{
\Question{
Which of the statements A, B, C, D given bellow are correct?
\[
\begin{array}{l}
\text{A: }\ \left(3x^3-5x^2+7x\right)'=9x^2-10x+7\text{, }x\in\mathbb{R} \\	
\text{B: }\  \left(x^2+2x^{-2}-5\right)'=2x-4x^{-3}\text{, }x\in\mathbb{R}\setminus\{0\} \\
\text{C: }\  \left(5x^4-7x^3+2\pi\right)'=20x^3-21x^2+2\text{, }x\in\mathbb{R} \\
\text{D: }\  \left(x^4-\frac4{x^4}\right)'=4x^3-\frac{16}{x^3}\text{, }x\in\mathbb{R}\setminus\{0\}
\end{array}
\]
The only correct statements are:}
{A, B}{A, C}{A, B, C}{A, D}{A, C, D}{}}
\LANGpl{
\Question{
Które z poniższych twierdzeń  A, B, C, D są prawidłowe?
\[
\begin{array}{l}
\text{A: }\ \left(3x^3-5x^2+7x\right)'=9x^2-10x+7\text{, }x\in\mathbb{R} \\	
\text{B: }\  \left(x^2+2x^{-2}-5\right)'=2x-4x^{-3}\text{, }x\in\mathbb{R}\setminus\{0\} \\
\text{C: }\  \left(5x^4-7x^3+2\pi\right)'=20x^3-21x^2+2\text{, }x\in\mathbb{R} \\
\text{D: }\  \left(x^4-\frac4{x^4}\right)'=4x^3-\frac{16}{x^3}\text{, }x\in\mathbb{R}\setminus\{0\}
\end{array}
\]
Jedyne poprawne twierdzenia to:}
{A, B}{A, C}{A, B, C}{A, D}{A, C, D}{}}
\LANGcs{
\Question{
Která z tvrzení A, B, C, D uvedených níže jsou správná?
\[
\begin{array}{l}
\text{A: }\ \left(3x^3-5x^2+7x\right)'=9x^2-10x+7\text{, }x\in\mathbb{R} \\	
\text{B: }\  \left(x^2+2x^{-2}-5\right)'=2x-4x^{-3}\text{, }x\in\mathbb{R}\setminus\{0\} \\
\text{C: }\  \left(5x^4-7x^3+2\pi\right)'=20x^3-21x^2+2\text{, }x\in\mathbb{R} \\
\text{D: }\  \left(x^4-\frac4{x^4}\right)'=4x^3-\frac{16}{x^3}\text{, }x\in\mathbb{R}\setminus\{0\}
\end{array}
\]
Jedinými správnými tvrzeními jsou:}
{A, B}{A, C}{A, B, C}{A, D}{A, C, D}{}}
\LANGsk{
\Question{
Ktoré z tvrdení A, B, C, D uvedených nižšie sú správne?
\[
\begin{array}{l}
\text{A: }\ \left(3x^3-5x^2+7x\right)'=9x^2-10x+7\text{, }x\in\mathbb{R} \\	
\text{B: }\  \left(x^2+2x^{-2}-5\right)'=2x-4x^{-3}\text{, }x\in\mathbb{R}\setminus\{0\} \\
\text{C: }\  \left(5x^4-7x^3+2\pi\right)'=20x^3-21x^2+2\text{, }x\in\mathbb{R} \\
\text{D: }\  \left(x^4-\frac4{x^4}\right)'=4x^3-\frac{16}{x^3}\text{, }x\in\mathbb{R}\setminus\{0\}
\end{array}
\]
Jedinými správnymi tvrdeniami sú:}
{A, B}{A, C}{A, B, C}{A, D}{A, C, D}{}}
\LANGes{
\Question{
De los enunciados A, B, C, D que se muestran debajo ¿Cuáles son correctos? 
\[
\begin{array}{l}
\text{A: }\ \left(3x^3-5x^2+7x\right)'=9x^2-10x+7\text{, }x\in\mathbb{R} \\	
\text{B: }\  \left(x^2+2x^{-2}-5\right)'=2x-4x^{-3}\text{, }x\in\mathbb{R}\setminus\{0\} \\
\text{C: }\  \left(5x^4-7x^3+2\pi\right)'=20x^3-21x^2+2\text{, }x\in\mathbb{R} \\
\text{D: }\  \left(x^4-\frac4{x^4}\right)'=4x^3-\frac{16}{x^3}\text{, }x\in\mathbb{R}\setminus\{0\}
\end{array}
\]
Los enunciados correctos son solamente:}
{A, B}{A, C}{A, B, C}{A, D}{A, C, D}{}}
\End

\Data\ID{1003112009}
\Updated{2020-07-15 03:07:12}
\Level{A}
\SubArea{Derivative}
\LANGen{
\Question{
Which of the statements A, B, C, D given bellow are incorrect?
\[
\begin{array}{l}
\text{A: }\ \left(\sqrt{x}-\cos x+\mathrm{ln}\, x\right)'=\frac1{2\sqrt x}+\sin x+\frac1x\text{, }x\in\mathbb{R}^+ \\
\text{B: }\ \left(3\sin x-3^x+x^3 \right)'=3\cos x-3^x+3x^2 \text{, }x\in\mathbb{R} \\
\text{C: }\ \left(\mathrm{cotg}\, x-\mathrm{log}_3 5\right)'=-\frac1{\sin^2 x}\text{, }x\in\mathbb{R}\setminus\bigcup\limits_{k\in\mathbb{Z}}\{k\pi\} \\ 
\text{D: }\ \left(3\sqrt[3]{x^2}-2\mathrm{e}^x \right)'=\frac2{\sqrt[3]x}-2\mathrm{e}^x\text{, }x\in\mathbb{R}\setminus\{0\}
\end{array}
\]
The only incorrect statements are:}
{B}{B, C}{A, B, C}{C, D}{B, C, D}{C}}
\LANGpl{
\Question{
Które z poniższych twierdzeń  A, B, C, D  są nieprawidłowe?
\[
\begin{array}{l}
\text{A: }\ \left(\sqrt{x}-\cos x+\mathrm{ln}\, x\right)'=\frac1{2\sqrt x}+\sin x+\frac1x\text{, }x\in\mathbb{R}^+ \\
\text{B: }\ \left(3\sin x-3^x+x^3 \right)'=3\cos x-3^x+3x^2 \text{, }x\in\mathbb{R} \\
\text{C: }\ \left(\mathrm{cotg}\, x-\mathrm{log}_3 5\right)'=-\frac1{\sin^2 x}\text{, }x\in\mathbb{R}\setminus\bigcup\limits_{k\in\mathbb{Z}}\{k\pi\} \\ 
\text{D: }\ \left(3\sqrt[3]{x^2}-2\mathrm{e}^x \right)'=\frac2{\sqrt[3]x}-2\mathrm{e}^x\text{, }x\in\mathbb{R}\setminus\{0\}
\end{array}
\]
Nieprawidłowe twierdzenia to:}
{B}{B, C}{A, B, C}{C, D}{B, C, D}{C}}
\LANGcs{
\Question{
Která z níže uvedených tvrzení A, B, C, D nejsou pravdivá?
\[
\begin{array}{l}
\text{A: }\ \left(\sqrt{x}-\cos x+\mathrm{ln}\, x\right)'=\frac1{2\sqrt x}+\sin x+\frac1x\text{, }x\in\mathbb{R}^+ \\
\text{B: }\ \left(3\sin x-3^x+x^3 \right)'=3\cos x-3^x+3x^2 \text{, }x\in\mathbb{R} \\
\text{C: }\ \left(\mathrm{cotg}\, x-\mathrm{log}_3 5\right)'=-\frac1{\sin^2 x}\text{, }x\in\mathbb{R}\setminus\bigcup\limits_{k\in\mathbb{Z}}\{k\pi\} \\ 
\text{D: }\ \left(3\sqrt[3]{x^2}-2\mathrm{e}^x \right)'=\frac2{\sqrt[3]x}-2\mathrm{e}^x\text{, }x\in\mathbb{R}\setminus\{0\}
\end{array}
\]
Jedinými nesprávnými tvrzeními jsou:}
{B}{B, C}{A, B, C}{C, D}{B, C, D}{C}}
\LANGsk{
\Question{
Ktoré z tvrdení A, B, C, D uvedených nižšie nie sú správne?
\[
\begin{array}{l}
\text{A: }\ \left(\sqrt{x}-\cos x+\mathrm{ln}\, x\right)'=\frac1{2\sqrt x}+\sin x+\frac1x\text{, }x\in\mathbb{R}^+ \\
\text{B: }\ \left(3\sin x-3^x+x^3 \right)'=3\cos x-3^x+3x^2 \text{, }x\in\mathbb{R} \\
\text{C: }\ \left(\mathrm{cotg}\, x-\mathrm{log}_3 5\right)'=-\frac1{\sin^2 x}\text{, }x\in\mathbb{R}\setminus\bigcup\limits_{k\in\mathbb{Z}}\{k\pi\} \\ 
\text{D: }\ \left(3\sqrt[3]{x^2}-2\mathrm{e}^x \right)'=\frac2{\sqrt[3]x}-2\mathrm{e}^x\text{, }x\in\mathbb{R}\setminus\{0\}
\end{array}
\]
Jedinými nesprávnymi tvrdeniami sú:}
{B}{B, C}{A, B, C}{C, D}{B, C, D}{C}}
\LANGes{
\Question{
De los enunciados A, B, C, D ¿Cuáles  son incorrectos?
\[
\begin{array}{l}
\text{A: }\ \left(\sqrt{x}-\cos x+\mathrm{ln}\, x\right)'=\frac1{2\sqrt x}+\sin x+\frac1x\text{, }x\in\mathbb{R}^+ \\
\text{B: }\ \left(3\sin x-3^x+x^3 \right)'=3\cos x-3^x+3x^2 \text{, }x\in\mathbb{R} \\
\text{C: }\ \left(\mathrm{cotg}\, x-\mathrm{log}_3 5\right)'=-\frac1{\sin^2 x}\text{, }x\in\mathbb{R}\setminus\bigcup\limits_{k\in\mathbb{Z}}\{k\pi\} \\ 
\text{D: }\ \left(3\sqrt[3]{x^2}-2\mathrm{e}^x \right)'=\frac2{\sqrt[3]x}-2\mathrm{e}^x\text{, }x\in\mathbb{R}\setminus\{0\}
\end{array}
\]
Los enunciados incorrectos son:}
{B}{B, C}{A, B, C}{C, D}{B, C, D}{C}}
\End

\Data\ID{1103023901}
\Updated{2020-07-15 07:07:23}
\Level{A}
\SubArea{Derivative}
\IMG{1103023901.pdf}
\LANGen{
\Question{
The graph of \( f' \) is given in figure. Function \( f' \) is the derivative of a function \( f \). Choose which of the following  graphs is the graph of \( f \).}
{}{}{}{}{}{}}
\LANGpl{
\Question{
Rysunek   \( f' \) przedstawia wykres  pochodnej  funkcji \( f \). Wskaż, który z poniższych wykresów jest wykresem funkcji \( f \).}
{}{}{}{}{}{}}
\LANGcs{
\Question{
Na obrázku je graf derivace jedné z níže zobrazených funkcí. Vyberte které.}
{}{}{}{}{}{}}
\LANGsk{
\Question{
Na obrázku je graf derivácie jednej z nižšie zobrazených funkcií. Vyberte ktoré.}
{}{}{}{}{}{}}
\LANGes{
\Question{
NA}
{}{}{}{}{}{}}

\IMGanswer{1103023901a.pdf}
\IMGanswer{1103023901b.pdf}
\IMGanswer{1103023901c.pdf}
\IMGanswer{1103023901d.pdf}\End

\Data\ID{1103023902}
\Updated{2020-09-23 01:09:32}
\Level{A}
\SubArea{Derivative}
\IMG{1103023902.pdf}
\LANGen{
\Question{
The figure shows the derivative of one of the functions displayed below. Choose the function.}
{}{}{}{}{}{}}
\LANGpl{
\Question{
Rysunek \( f' \) przedstawia wykres  pochodnej  funkcji \( f \). Wskaż, który z poniższych wykresów jest wykresem funkcji \( f \).}
{}{}{}{}{}{}}
\LANGcs{
\Question{
Na obrázku je graf derivace jedné z níže zobrazených funkcí. Vyberte které.}
{}{}{}{}{}{}}
\LANGsk{
\Question{
Na obrázku je graf derivácie jednej z nižšie zobrazených funkcií. Vyberte ktoré.}
{}{}{}{}{}{}}
\LANGes{
\Question{
NA}
{}{}{}{}{}{}}

\IMGanswer{1103023902a.pdf}
\IMGanswer{1103023902b.pdf}
\IMGanswer{1103023902c.pdf}
\IMGanswer{1103023902d.pdf}\End

\Data\ID{1103023903}
\Updated{2020-07-15 03:07:10}
\Level{A}
\SubArea{Derivative}
\IMG{1103023903.pdf}
\LANGen{
\Question{
The graph of \( f' \) is given in figure. Function \( f' \) is the derivative of a function \( f \). Choose which of the following  graphs is the graph of \( f \).}
{}{}{}{}{}{}}
\LANGpl{
\Question{
Rysunek  \( f' \) przedstawia wykres  pochodnej  funkcji \( f \). Wskaż, który z poniższych wykresów jest wykresem funkcji \( f \).}
{}{}{}{}{}{}}
\LANGcs{
\Question{
Na obrázku je graf derivace jedné z níže zobrazených funkcí. Vyberte které.}
{}{}{}{}{}{}}
\LANGsk{
\Question{
Na obrázku je graf derivácie jednej z nižšie zobrazených funkcií. Vyberte ktoré.}
{}{}{}{}{}{}}
\LANGes{
\Question{
NA}
{}{}{}{}{}{}}

\IMGanswer{1103023903a.pdf}
\IMGanswer{1103023903b.pdf}
\IMGanswer{1103023903c.pdf}
\IMGanswer{1103023903d.pdf}\End

\Data\ID{1103023904}
\Updated{2020-07-15 03:07:10}
\Level{A}
\SubArea{Derivative}
\IMG{1103023904.pdf}
\LANGen{
\Question{
The graph of \( f' \) is given in figure. Function \( f' \) is the derivative of a function \( f \). Choose which of the following  graphs is the graph of \( f \).}
{}{}{}{}{}{}}
\LANGpl{
\Question{
Rysunek \( f' \) przedstawia wykres  pochodnej  funkcji  \( f \). Wskaż, który z poniższych wykresów jest wykresem funkcji \( f \).}
{}{}{}{}{}{}}
\LANGcs{
\Question{
Na obrázku je graf derivace jedné z níže zobrazených funkcí. Vyberte které.}
{}{}{}{}{}{}}
\LANGsk{
\Question{
Na obrázku je graf derivácie jednej z nižšie zobrazených funkcií. Vyberte ktoré.}
{}{}{}{}{}{}}
\LANGes{
\Question{
NA}
{}{}{}{}{}{}}

\IMGanswer{1103023904a.pdf}
\IMGanswer{1103023904b.pdf}
\IMGanswer{1103023904c.pdf}
\IMGanswer{1103023904d.pdf}\End

\Data\ID{1103023905}
\Updated{2020-07-15 03:07:10}
\Level{A}
\SubArea{Derivative}
\IMG{1103023905.pdf}
\LANGen{
\Question{
The graph of \( f' \) is given in figure. Function \( f' \) is the derivative of a function \( f \). Choose which of the following  graphs is the graph of \( f \).}
{}{}{}{}{}{}}
\LANGpl{
\Question{
Rysunek \( f' \) przedstawia wykres  pochodnej  funkcji \( f \). Wskaż, który z poniższych wykresów jest wykresem funkcji \( f \).}
{}{}{}{}{}{}}
\LANGcs{
\Question{
Na obrázku je graf derivace jedné z níže zobrazených funkcí. Vyberte které.}
{}{}{}{}{}{}}
\LANGsk{
\Question{
Na obrázku je graf derivácie jednej z nižšie zobrazených funkcií. Vyberte ktoré.}
{}{}{}{}{}{}}
\LANGes{
\Question{
NA}
{}{}{}{}{}{}}

\IMGanswer{1103023905a.pdf}
\IMGanswer{1103023905b.pdf}
\IMGanswer{1103023905c.pdf}
\IMGanswer{1103023905d_0.pdf}\End

\Data\ID{1103024401}
\Updated{2020-07-15 03:07:10}
\Level{A}
\SubArea{Derivative}
\IMG{1103024401_0.pdf}
\LANGen{
\Question{
The graph of \( f \) is given in figure. Choose which of the following graphs is the graph of \( f' \)
(function \( f' \) is the derivative of a function \( f \)).}
{}{}{}{}{}{}}
\LANGpl{
\Question{
Rysunek przedstawia wykres funkcji  \( f \). Wybierz, który z poniższych wykresów jest wykresem funkcji \( f' \)
(funkcja  \( f' \) jest pochodną funkcji \( f \)).}
{}{}{}{}{}{}}
\LANGcs{
\Question{
Na obrázku je dán graf funkce \( f \). Vyberte, který z následujících grafů je grafem funkce \( f' \) (funkce  \( f' \) je derivace funkce \( f \)).}
{}{}{}{}{}{}}
\LANGsk{
\Question{
Na obrázku je daný graf funkcie \( f \). Vyberte, ktorý z nasledujúcich grafov je grafom funkcie \( f' \) (funkcia  \( f' \) je derivácia funkcie \( f \)).}
{}{}{}{}{}{}}
\LANGes{
\Question{
NA}
{}{}{}{}{}{}}

\IMGanswer{1103024401a_0.pdf}
\IMGanswer{1103024401b_0.pdf}
\IMGanswer{1103024401c_0.pdf}
\IMGanswer{1103024401d_0.pdf}\End

\Data\ID{1103024402}
\Updated{2020-07-15 03:07:10}
\Level{A}
\SubArea{Derivative}
\IMG{1103024402.pdf}
\LANGen{
\Question{
The graph of \( f \) is given in figure. Choose which of the following graphs is the graph of \( f' \)
(function \( f' \) is the derivative of a function \( f \)).}
{}{}{}{}{}{}}
\LANGpl{
\Question{
Rysunek przedstawia wykres funkcji  \( f \). Wybierz, który z poniższych wykresów jest wykresem funkcji  \( f' \)
(funkcja \( f' \) jest pochodną funkcji \( f \)).}
{}{}{}{}{}{}}
\LANGcs{
\Question{
Na obrázku je dán graf funkce \( f \). Vyberte, který z následujících grafů je grafem funkce \( f' \) (funkce  \( f' \) je derivace funkce \( f \)).}
{}{}{}{}{}{}}
\LANGsk{
\Question{
Na obrázku je daný graf funkcie \( f \). Vyberte, ktorý z nasledujúcich grafov je grafom funkcie \( f' \) (funkcia  \( f' \) je derivácia funkcie \( f \)).}
{}{}{}{}{}{}}
\LANGes{
\Question{
NA}
{}{}{}{}{}{}}

\IMGanswer{1103024402a.pdf}
\IMGanswer{1103024402b.pdf}
\IMGanswer{1103024402c.pdf}
\IMGanswer{1103024402d.pdf}\End

\Data\ID{1103024403}
\Updated{2020-09-23 01:09:56}
\Level{A}
\SubArea{Derivative}
\IMG{1103024403.pdf}
\LANGen{
\Question{
The graph of \( f \) is given in the figure. Choose which of the following graphs is the graph of \( f' \).
(Function \( f' \) is the derivative of a function \( f \)).}
{}{}{}{}{}{}}
\LANGpl{
\Question{
Rysunek przedstawia wykres funkcji \( f \). Wybierz, który z poniższych wykresów jest wykresem funkcji \( f' \)
(funkcja \( f' \) jest pochodną funkcji \( f \)).}
{}{}{}{}{}{}}
\LANGcs{
\Question{
Na obrázku je dán graf funkce \( f \). Vyberte, který z následujících grafů je grafem funkce \( f' \) (funkce  \( f' \) je derivace funkce \( f \)).}
{}{}{}{}{}{}}
\LANGsk{
\Question{
Na obrázku je daný graf funkcie \( f \). Vyberte, ktorý z nasledujúcich grafov je grafom funkcie \( f' \) (funkcia  \( f' \) je derivácia funkcie \( f \)).}
{}{}{}{}{}{}}
\LANGes{
\Question{
NA}
{}{}{}{}{}{}}

\IMGanswer{1103024403a.pdf}
\IMGanswer{1103024403b.pdf}
\IMGanswer{1103024403c.pdf}
\IMGanswer{1103024403d.pdf}\End

\Data\ID{1103024404}
\Updated{2020-09-23 01:09:32}
\Level{A}
\SubArea{Derivative}
\IMG{1103024404.pdf}
\LANGen{
\Question{
The graph of \( f \) is given in the figure. Choose which of the following graphs is the graph of \( f' \).
(Function \( f' \) is the derivative of a function \( f \).)}
{}{}{}{}{}{}}
\LANGpl{
\Question{
Rysunek przedstawia wykres funkcji \( f \). Wybierz, który z poniższych wykresów jest wykresem funkcji \( f' \)
(funkcja \( f' \) jest pochodną funkcji \( f \)).}
{}{}{}{}{}{}}
\LANGcs{
\Question{
Na obrázku je dán graf funkce \( f \). Vyberte, který z následujících grafů je grafem funkce \( f' \) (funkce  \( f' \) je derivace funkce \( f \)).}
{}{}{}{}{}{}}
\LANGsk{
\Question{
Na obrázku je daný graf funkcie \( f \). Vyberte, ktorý z nasledujúcich grafov je grafom funkcie \( f' \) (funkcia  \( f' \) je deriváciou funkcie \( f \)).}
{}{}{}{}{}{}}
\LANGes{
\Question{
NA}
{}{}{}{}{}{}}

\IMGanswer{1103024404a.pdf}
\IMGanswer{1103024404b.pdf}
\IMGanswer{1103024404c.pdf}
\IMGanswer{1103024404d.pdf}\End

\Data\ID{1103024405}
\Updated{2020-07-15 03:07:10}
\Level{A}
\SubArea{Derivative}
\IMG{1103024405_0.pdf}
\LANGen{
\Question{
The graph of \( f \) is given in figure. Choose which of the following graphs is the graph of \( f' \)
(function \( f' \) is the derivative of a function \( f \)).}
{}{}{}{}{}{}}
\LANGpl{
\Question{
Rysunek przedstawia wykres funkcji \( f \).  Wybierz, który z poniższych wykresów jest wykresem funkcji \( f' \)
(funkcja \( f' \) jest pochodną funkcji \( f \)).}
{}{}{}{}{}{}}
\LANGcs{
\Question{
Na obrázku je dán graf funkce \( f \). Vyberte, který z následujících grafů je grafem funkce \( f' \) (funkce  \( f' \) je derivace funkce \( f \)).}
{}{}{}{}{}{}}
\LANGsk{
\Question{
Na obrázku je daný graf funkcie \( f \). Vyberte, ktorý z nasledujúcich grafov je grafom funkcie \( f' \) (funkcia  \( f' \) je deriváciou funkcie \( f \)).}
{}{}{}{}{}{}}
\LANGes{
\Question{
NA}
{}{}{}{}{}{}}

\IMGanswer{1103024405a_0.pdf}
\IMGanswer{1103024405b_0.pdf}
\IMGanswer{1103024405c_0.pdf}
\IMGanswer{1103024405d_0.pdf}\End

\Data\ID{1103024406}
\Updated{2020-07-15 03:07:10}
\Level{A}
\SubArea{Derivative}
\IMG{1103024406.pdf}
\LANGen{
\Question{
The graph of \( f \) is given in figure. Choose which of the following graphs is the graph of \( f' \)
(function \( f' \) is the derivative of a function \( f \)).}
{}{}{}{}{}{}}
\LANGpl{
\Question{
Rysunek przedstawia wykres funkcji  \( f \). Wybierz, który z poniższych wykresów jest wykresem funkcji \( f' \)
(funkcja \( f' \) jest pochodna funkcji \( f \)).}
{}{}{}{}{}{}}
\LANGcs{
\Question{
Na obrázku je dán graf funkce \( f \). Vyberte, který z následujících grafů je grafem funkce \( f' \) (funkce  \( f' \) je derivace funkce \( f \)).}
{}{}{}{}{}{}}
\LANGsk{
\Question{
Na obrázku je daný graf funkcie \( f \). Vyberte, ktorý z nasledujúcich grafov je grafom funkcie \( f' \) (funkcia  \( f' \) je deriváciou funkcie \( f \)).}
{}{}{}{}{}{}}
\LANGes{
\Question{
NA}
{}{}{}{}{}{}}

\IMGanswer{1103024406a.pdf}
\IMGanswer{1103024406b.pdf}
\IMGanswer{1103024406c.pdf}
\IMGanswer{1103024406d.pdf}\End

\Data\ID{1103024407}
\Updated{2020-07-15 03:07:10}
\Level{A}
\SubArea{Derivative}
\IMG{1103024407_0.pdf}
\LANGen{
\Question{
The graph of \( f \) is given in figure. Choose which of the following graphs is the graph of \( f' \)
(function \( f' \) is the derivative of a function \( f \)).}
{}{}{}{}{}{}}
\LANGpl{
\Question{
Rysunek przedstawia funkcję  \( f \). Wybierz, który z poniższych wykresów jest wykresem funkcji \( f' \)
(funkcja \( f' \) jest pochodną funkcji \( f \)).}
{}{}{}{}{}{}}
\LANGcs{
\Question{
Na obrázku je dán graf funkce \( f \). Vyberte, který z následujících grafů je grafem funkce \( f' \) (funkce  \( f' \) je derivace funkce \( f \)).}
{}{}{}{}{}{}}
\LANGsk{
\Question{
Na obrázku je daný graf funkcie \( f \). Vyberte, ktorý z nasledujúcich grafov je grafom funkcie \( f' \) (funkcia  \( f' \) je deriváciou funkcie \( f \)).}
{}{}{}{}{}{}}
\LANGes{
\Question{
NA}
{}{}{}{}{}{}}

\IMGanswer{1103024407a_0.pdf}
\IMGanswer{1103024407b_0.pdf}
\IMGanswer{1103024407c_0.pdf}
\IMGanswer{1103024407d_0.pdf}\End

\Data\ID{1103024408}
\Updated{2020-07-15 03:07:10}
\Level{A}
\SubArea{Derivative}
\IMG{1103024408.pdf}
\LANGen{
\Question{
The graph of \( f \) is given in figure. Choose which of the following graphs is the graph of \( f' \)
(function \( f' \) is the derivative of a function \( f \)).}
{}{}{}{}{}{}}
\LANGpl{
\Question{
Rysunek przedstawia funkcję \( f \). Wybierz, który z poniższych wykresów jest wykresem funkcji \( f' \)
(funkcja \( f' \) jest pochodną funkcji \( f \)).}
{}{}{}{}{}{}}
\LANGcs{
\Question{
Na obrázku je dán graf funkce \( f \). Vyberte, který z následujících grafů je grafem funkce \( f' \) (funkce  \( f' \) je derivace funkce \( f \)).}
{}{}{}{}{}{}}
\LANGsk{
\Question{
Na obrázku je daný graf funkcie \( f \). Vyberte, ktorý z nasledujúcich grafov je grafom funkcie \( f' \) (funkcia  \( f' \) je deriváciou funkcie \( f \)).}
{}{}{}{}{}{}}
\LANGes{
\Question{
NA}
{}{}{}{}{}{}}

\IMGanswer{1103024408a.pdf}
\IMGanswer{1103024408b.pdf}
\IMGanswer{1103024408c.pdf}
\IMGanswer{1103024408d.pdf}\End

\Data\ID{1103024409}
\Updated{2020-07-15 03:07:10}
\Level{A}
\SubArea{Derivative}
\IMG{1103024409.pdf}
\LANGen{
\Question{
The graph of \( f \) is given in figure. Choose which of the following graphs is the graph of \( f' \)
(function \( f' \) is the derivative of a function \( f \)).}
{}{}{}{}{}{}}
\LANGpl{
\Question{
Rysunek przedstawia wykres funkcji  \( f \). Wybierz, który z poniższych wykresów jest wykresem funkcji \( f' \)
(funkcja \( f' \) jest pochodną funkcji  \( f \)).}
{}{}{}{}{}{}}
\LANGcs{
\Question{
Na obrázku je dán graf funkce \( f \). Vyberte, který z následujících grafů je grafem funkce \( f' \) (funkce  \( f' \) je derivace funkce \( f \)).}
{}{}{}{}{}{}}
\LANGsk{
\Question{
Na obrázku je daný graf funkcie \( f \). Vyberte, ktorý z nasledujúcich grafov je grafom funkcie \( f' \) (funkcia  \( f' \) je derivácia funkcie \( f \)).}
{}{}{}{}{}{}}
\LANGes{
\Question{
NA}
{}{}{}{}{}{}}

\IMGanswer{1103024409a.pdf}
\IMGanswer{1103024409b.pdf}
\IMGanswer{1103024409c.pdf}
\IMGanswer{1103024409d.pdf}\End

\Data\ID{1103143701}
\Updated{2020-07-15 03:07:10}
\Level{A}
\SubArea{Derivative}
\LANGen{
\Question{
Choose the graph of a function \( f \) for which \( f(0)=0 \), \( f'(0)=3 \), \( f'(1)=0 \), \( f'(2)=-1 \) (\( f' \) is the derivative of the function \( f \)).}
{}{}{}{}{}{}}
\LANGpl{
\Question{
Wskaż wykres funkcji \( f \) dla której \( f(0)=0 \), \( f'(0)=3 \), \( f'(1)=0 \), \( f'(2)=-1 \) (\( f' \) jest pochodną funkcji \( f \)).}
{}{}{}{}{}{}}
\LANGcs{
\Question{
Vyberte graf funkce \( f \), pro kterou platí \( f(0)=0 \), \( f'(0)=3 \), \( f'(1)=0 \), \( f'(2)=-1 \) (\( f' \) je derivace funkce \( f \)).}
{}{}{}{}{}{}}
\LANGsk{
\Question{
Vyberte graf funkcie \( f \), pre ktorú platí \( f(0)=0 \), \( f'(0)=3 \), \( f'(1)=0 \), \( f'(2)=-1 \) (\( f' \) je derivácia funkcie \( f \)).}
{}{}{}{}{}{}}
\LANGes{
\Question{
NA}
{}{}{}{}{}{}}

\IMGanswer{1103143701a.pdf}
\IMGanswer{1103143701b.pdf}
\IMGanswer{1103143701c.pdf}
\IMGanswer{1103143701d.pdf}\End

\Data\ID{1103143702}
\Updated{2020-07-15 03:07:10}
\Level{A}
\SubArea{Derivative}
\LANGen{
\Question{
Choose the graph of a function \( f \) for which \( f(0)=0 \), \( f'(1)=-1 \), \( f'(2)=0 \), \( f'(3)=3 \) (\( f' \) is the derivative of the function \( f \)).}
{}{}{}{}{}{}}
\LANGpl{
\Question{
Wskaż wykres funkcji \( f \) dla której \( f(0)=0 \), \( f'(1)=-1 \), \( f'(2)=0 \), \( f'(3)=3 \) (\( f' \) jest pochodną funkcji \( f \)).}
{}{}{}{}{}{}}
\LANGcs{
\Question{
Vyberte graf funkce \( f \), pro kterou platí \( f(0)=0 \), \( f'(1)=-1 \), \( f'(2)=0 \), \( f'(3)=3 \) (\( f' \) je derivace funkce \( f \)).}
{}{}{}{}{}{}}
\LANGsk{
\Question{
Vyberte graf funkcie \( f \), pre ktorú platí \( f(0)=0 \), \( f'(1)=-1 \), \( f'(2)=0 \), \( f'(3)=3 \) (\( f' \) je derivácia funkcie \( f \)).}
{}{}{}{}{}{}}
\LANGes{
\Question{
NA}
{}{}{}{}{}{}}

\IMGanswer{1103143702a.pdf}
\IMGanswer{1103143702b.pdf}
\IMGanswer{1103143702c.pdf}
\IMGanswer{1103143702d.pdf}\End

\Data\ID{1103143703}
\Updated{2020-07-15 03:07:10}
\Level{A}
\SubArea{Derivative}
\LANGen{
\Question{
Choose the graph of a function \( f \) for which \( f'(-2)=0 \), \( f'(-1)=-2 \), \( f'(2)=4 \) (\( f' \) is the derivative of the function \( f \)).}
{}{}{}{}{}{}}
\LANGpl{
\Question{
Wskaż wykres funkcji \( f \) dla której \( f'(-2)=0 \), \( f'(-1)=-2 \), \( f'(2)=4 \) (\( f' \) jest pochodną funkcji \( f \)).}
{}{}{}{}{}{}}
\LANGcs{
\Question{
Vyberte graf funkce \( f \), pro kterou platí \( f'(-2)=0 \), \( f'(-1)=-2 \), \( f'(2)=4 \) (\( f' \) je derivace funkce \( f \)).}
{}{}{}{}{}{}}
\LANGsk{
\Question{
Vyberte graf funkcie \( f \), pre ktorú platí \( f'(-2)=0 \), \( f'(-1)=-2 \), \( f'(2)=4 \) (\( f' \) je derivácia funkcie \( f \)).}
{}{}{}{}{}{}}
\LANGes{
\Question{
NA}
{}{}{}{}{}{}}

\IMGanswer{1103143703a_0.pdf}
\IMGanswer{1103143703b_0.pdf}
\IMGanswer{1103143703c_0.pdf}
\IMGanswer{1103143703d_0.pdf}\End

\Data\ID{1103143704}
\Updated{2020-07-15 03:07:10}
\Level{A}
\SubArea{Derivative}
\LANGen{
\Question{
Choose the graph of a function \( f \) for which \( f(-2)=0 \), \( f'(-1)=0 \), \( f'(0)=1 \), \( f'(2)=-3 \) (\( f' \) is the derivative of the function \( f \)).}
{}{}{}{}{}{}}
\LANGpl{
\Question{
Wskaż wykres funkcji \( f \) dla której \( f(-2)=0 \), \( f'(-1)=0 \), \( f'(0)=1 \), \( f'(2)=-3 \) (\( f' \) jest pochodną funkcji \( f \)).}
{}{}{}{}{}{}}
\LANGcs{
\Question{
Vyberte graf funkce \( f \), pro kterou platí \( f(-2)=0 \), \( f'(-1)=0 \), \( f'(0)=1 \), \( f'(2)=-3 \) (\( f' \) je derivace funkce \( f \)).}
{}{}{}{}{}{}}
\LANGsk{
\Question{
Vyberte graf funkcie \( f \), pre ktorú platí \( f(-2)=0 \), \( f'(-1)=0 \), \( f'(0)=1 \), \( f'(2)=-3 \) (\( f' \) je derivácia funkcie \( f \)).}
{}{}{}{}{}{}}
\LANGes{
\Question{
NA}
{}{}{}{}{}{}}

\IMGanswer{1103143704a_0.pdf}
\IMGanswer{1103143704b_0.pdf}
\IMGanswer{1103143704c_0.pdf}
\IMGanswer{1103143704d_0.pdf}\End

\Data\ID{1103143705}
\Updated{2021-02-22 04:02:13}
\Level{A}
\SubArea{Derivative}
\LANGen{
\Question{
Choose the graph of a function \( f \) for which \( f'(-1) > 0 \), \( f'(0)=0 \), \( f'(1) < 0 \), \( f'(3) > 0 \) (\( f' \) is the derivative of the function \( f \)).}
{}{}{}{}{}{}}
\LANGpl{
\Question{
Wskaż wykres funkcji \( f \) dla której \( f'(-1) > 0 \), \( f'(0)=0 \), \( f'(1) < 0 \), \( f'(3) > 0 \) (\( f' \) jest pochodną funkcji \( f \)).}
{}{}{}{}{}{}}
\LANGcs{
\Question{
Vyberte graf funkce \( f \), pro kterou platí \( f'(-1) > 0 \), \( f'(0)=0 \), \( f'(1) < 0 \), \( f'(3) > 0 \) (\( f' \) je derivace funkce \( f \)).}
{}{}{}{}{}{}}
\LANGsk{
\Question{
Vyberte graf funkcie \( f \), pre ktorú platí \( f'(-1) > 0 \), \( f'(0)=0 \), \( f'(1) < 0 \), \( f'(3) > 0 \) (\( f' \) je derivácia funkcie \( f \)).}
{}{}{}{}{}{}}
\LANGes{
\Question{
NA}
{}{}{}{}{}{}}

\IMGanswer{1103143705a_0.pdf}
\IMGanswer{1103143705b_1.pdf}
\IMGanswer{1103143705c.pdf}
\IMGanswer{1103143705d.pdf}\End

\Data\ID{1103163601}
\Updated{2020-07-15 03:07:10}
\Level{A}
\SubArea{Derivative}
\IMG{1103163601.pdf}
\LANGen{
\Question{
The graph of \( f \) is given in the figure. Choose which of the following graphs is the graph of \( f' \). (The function \( f' \) is the derivative of the function \( f \).)}
{}{}{}{}{}{}}
\LANGpl{
\Question{
Dany jest wykres funkcji \( f \).  Wskaż, który z poniższych wykresów jest wykresem funkcji\( f' \). (Funkcja  \( f' \) jest pochodną funkcji \( f \).)}
{}{}{}{}{}{}}
\LANGcs{
\Question{
Na obrázku je dán graf funkce \( f \). Vyberte, který z následujících grafů je grafem funkce \( f' \). (Funkce  \( f' \) je derivace funkce \( f \).)}
{}{}{}{}{}{}}
\LANGsk{
\Question{
Na obrázku je daný graf funkcie \( f \). Vyberte, ktorý z nasledujúcich grafov je grafom funkcie \( f' \). (Funkcia  \( f' \) je derivácia funkcie \( f \).)}
{}{}{}{}{}{}}
\LANGes{
\Question{
NA}
{}{}{}{}{}{}}

\IMGanswer{1103163601a.pdf}
\IMGanswer{1103163601b.pdf}
\IMGanswer{1103163601c.pdf}
\IMGanswer{1103163601d.pdf}\End

\Data\ID{1103164701}
\Updated{2020-07-15 03:07:12}
\Level{A}
\SubArea{Derivative}
\IMG{1103164701.pdf}
\LANGen{
\Question{
The graph of \( f \) is given in the figure, where \( A \), \( B \) and \( C \) are points on the graph. If \( x_A \), \( x_B \) and \( x_C \) denote the \( x \)-coordinates of the points the  \( A \), \( B \) and \( C \), and if \( f' \) is the derivative of \( f \), then:}
{\( f'( x_A ) > f'( x_B ) > f'( x_C ) \)}{\( f'( x_A ) < f'( x_B ) = f' ( x_C ) \)}{\( f'( x_A ) > f'( x_B ) = f'(x_C ) \)}{\( f'(x_A ) < f'( x_B ) < f'( x_C ) \)}{\( f'( x_A ) = f'( x_B ) > f'( x_C ) \)}{}}
\LANGpl{
\Question{
Dany jest wykres funkcji \( f \), gdzie \( A \), \( B \) i \( C \) to punkty na wykresie. Jeśli \( x_A \), \( x_B \) i \( x_C \) określają \( x \)-współrzędne punktów  \( A \), \( B \) i \( C \), oraz  \( f' \) jest pochodną  \( f \), to:}
{\( f'( x_A ) > f'( x_B ) > f'( x_C ) \)}{\( f'( x_A ) < f'( x_B ) = f' ( x_C ) \)}{\( f'( x_A ) > f'( x_B ) = f'(x_C ) \)}{\( f'(x_A ) < f'( x_B ) < f'( x_C ) \)}{\( f'( x_A ) = f'( x_B ) > f'( x_C ) \)}{}}
\LANGcs{
\Question{
Na obrázku je dán graf funkce \( f \), přičemž  \( A \), \( B \) a \( C \) jsou body grafu této funkce. Jestliže \( x_A \), \( x_B \) a \( x_C \) jsou \( x \)-ové souřadnice bodů  \( A \), \( B \) a \( C \), a jestliže \( f' \) je derivace funkce \( f \), pak:}
{\( f'( x_A ) > f'( x_B ) > f'( x_C ) \)}{\( f'( x_A ) < f'( x_B ) = f' ( x_C ) \)}{\( f'( x_A ) > f'( x_B ) = f'(x_C ) \)}{\( f'(x_A ) < f'( x_B ) < f'( x_C ) \)}{\( f'( x_A ) = f'( x_B ) > f'( x_C ) \)}{}}
\LANGsk{
\Question{
Na obrázku je daný graf funkcie \( f \), pričom  \( A \), \( B \) a \( C \) sú body grafu tejto funkcie. Ak \( x_A \), \( x_B \) a \( x_C \) sú \( x \)-ové súradnice bodov  \( A \), \( B \) a \( C \), a ak \( f' \) je derivácia funkcie \( f \), potom:}
{\( f'( x_A ) > f'( x_B ) > f'( x_C ) \)}{\( f'( x_A ) < f'( x_B ) = f' ( x_C ) \)}{\( f'( x_A ) > f'( x_B ) = f'(x_C ) \)}{\( f'(x_A ) < f'( x_B ) < f'( x_C ) \)}{\( f'( x_A ) = f'( x_B ) > f'( x_C ) \)}{}}
\LANGes{
\Question{
Dada la gráfica de la función \( f \) donde \( A \), \( B \) y \( C \) son puntos en la gráfica. Si \( x_A \), \( x_B \) y \( x_C \) denotan las coordenadas del eje \( x \) de los puntos  \( A \), \( B \) y \( C \), y si \( f' \) es la derivada de \( f \), entonces:}
{\( f'( x_A ) > f'( x_B ) > f'( x_C ) \)}{\( f'( x_A ) < f'( x_B ) = f' ( x_C ) \)}{\( f'( x_A ) > f'( x_B ) = f'(x_C ) \)}{\( f'(x_A ) < f'( x_B ) < f'( x_C ) \)}{\( f'( x_A ) = f'( x_B ) > f'( x_C ) \)}{}}
\End

\Data\ID{1103164702}
\Updated{2020-07-15 03:07:12}
\Level{A}
\SubArea{Derivative}
\IMG{1103164702.pdf}
\LANGen{
\Question{
The graph of \( f \) is given in the figure, where \( A \), \( B \) and \( C \) are points on the graph and \( y \)-coordinate of the point \( B \) is the maximum value of the function \( f \). If \( x_A \), \( x_B \) and \( x_C \) denote the \( x \)-coordinates of the points \( A \), \( B \) and \( C \), and if \( f' \) is the derivative of \( f \), then:}
{\( f'( x_A ) > 0 \), \( f'( x_B ) = 0 \), \( f'( x_C ) < 0 \)}{\( f'( x_A ) > 0 \), \( f'( x_B ) > 0 \), \( f'( x_C ) < 0 \)}{\( f'( x_A ) < 0 \), \( f'( x_B ) = 0 \), \( f'( x_C ) < 0 \)}{\( f'( x_A ) < 0 \), \( f'( x_B ) = 0 \), \( f'( x_C ) > 0 \)}{}{}}
\LANGpl{
\Question{
Dany jest wykres funkcji  \( f \), gdzie \( A \), \( B \) i \( C \) to punkty na wykresie oraz  \( y \)-współrzędna punktu \( B \) jest maksymalną wartością funkcji \( f \). Jeśli \( x_A \), \( x_B \) i \( x_C \) oznaczają \( x \)-współrzędne punktów \( A \), \( B \) i \( C \), a  \( f' \) jest pochodną \( f \), to:}
{\( f'( x_A ) > 0 \), \( f'( x_B ) = 0 \), \( f'( x_C ) < 0 \)}{\( f'( x_A ) > 0 \), \( f'( x_B ) > 0 \), \( f'( x_C ) < 0 \)}{\( f'( x_A ) < 0 \), \( f'( x_B ) = 0 \), \( f'( x_C ) < 0 \)}{\( f'( x_A ) < 0 \), \( f'( x_B ) = 0 \), \( f'( x_C ) > 0 \)}{}{}}
\LANGcs{
\Question{
Na obrázku je dán graf funkce \( f \), přičemž  \( A \), \( B \) a \( C \) jsou body grafu této funkce a \( y \)-ová souřadnice bodu \( B \) je největší hodnota funkce \( f \). Jestliže \( x_A \), \( x_B \) a \( x_C \) jsou \( x \)-ové souřadnice bodů  \( A \), \( B \) a \( C \), a jestliže \( f' \) je derivace funkce \( f \), pak:}
{\( f'( x_A ) > 0 \), \( f'( x_B ) = 0 \), \( f'( x_C ) < 0 \)}{\( f'( x_A ) > 0 \), \( f'( x_B ) > 0 \), \( f'( x_C ) < 0 \)}{\( f'( x_A ) < 0 \), \( f'( x_B ) = 0 \), \( f'( x_C ) < 0 \)}{\( f'( x_A ) < 0 \), \( f'( x_B ) = 0 \), \( f'( x_C ) > 0 \)}{}{}}
\LANGsk{
\Question{
Na obrázku je daný graf funkcie \( f \), pričom  \( A \), \( B \) a \( C \) sú body grafu tejto funkcie a \( y \)-ová súradnice bodu \( B \) je najsvätejšia hodnota funkcie \( f \). Ak \( x_A \), \( x_B \) a \( x_C \) sú \( x \)-ové súradnice bodov \( A \), \( B \) a \( C \), a ak \( f' \) je derivácia funkcie \( f \), potom:}
{\( f'( x_A ) > 0 \), \( f'( x_B ) = 0 \), \( f'( x_C ) < 0 \)}{\( f'( x_A ) > 0 \), \( f'( x_B ) > 0 \), \( f'( x_C ) < 0 \)}{\( f'( x_A ) < 0 \), \( f'( x_B ) = 0 \), \( f'( x_C ) < 0 \)}{\( f'( x_A ) < 0 \), \( f'( x_B ) = 0 \), \( f'( x_C ) > 0 \)}{}{}}
\LANGes{
\Question{
Dada la gráfica de \( f \), donde \( A \), \( B \) y \( C \) son puntos de la gráfica y  la coordenada del eje \( y \) del punto \( B \) es el máximo valor de la función \( f \). Si \( x_A \), \( x_B \) y \( x_C \) denotan las coordenadas del eje \( x \) de los puntos \( A \), \( B \) y \( C \), y si \( f' \) es the derivada de \( f \), entonces:}
{\( f'( x_A ) > 0 \), \( f'( x_B ) = 0 \), \( f'( x_C ) < 0 \)}{\( f'( x_A ) > 0 \), \( f'( x_B ) > 0 \), \( f'( x_C ) < 0 \)}{\( f'( x_A ) < 0 \), \( f'( x_B ) = 0 \), \( f'( x_C ) < 0 \)}{\( f'( x_A ) < 0 \), \( f'( x_B ) = 0 \), \( f'( x_C ) > 0 \)}{}{}}
\End

\Data\ID{1103164703}
\Updated{2020-07-15 03:07:12}
\Level{A}
\SubArea{Derivative}
\IMG{1103164703.pdf}
\LANGen{
\Question{
The graph of \( f \) is given in the figure, where \( A \), \( B \) and \( C \) are points on the graph. If \( x_A \), \( x_B \) and \( x_C \) denote the \( x \)-coordinates of the points \( A \), \( B \) and \( C \) and if \( f' \) is the derivative of \( f \), then:}
{\( f'( x_A ) < f'( x_B ) < f'( x_C ) \)}{\( f'( x_A ) < f'( x_B ) = f'( x_C ) \)}{\( f'( x_A ) > f'( x_B ) = f'( x_C ) \)}{\( f'( x_A ) > f'( x_B ) > f'( x_C ) \)}{\( f'( x_A ) = f'( x_B ) < f'( x_C ) \)}{}}
\LANGpl{
\Question{
Dany jest wykres funkcji  \( f \), gdzie \( A \), \( B \) i \( C \) to punkty na wykresie. Jeśli  \( x_A \), \( x_B \) i \( x_C \) oznaczają \( x \)-współrzędne punktów \( A \), \( B \) i \( C \) oraz \( f' \) jest pochodną  \( f \), to:}
{\( f'( x_A ) < f'( x_B ) < f'( x_C ) \)}{\( f'( x_A ) < f'( x_B ) = f'( x_C ) \)}{\( f'( x_A ) > f'( x_B ) = f'( x_C ) \)}{\( f'( x_A ) > f'( x_B ) > f'( x_C ) \)}{\( f'( x_A ) = f'( x_B ) < f'( x_C ) \)}{}}
\LANGcs{
\Question{
Na obrázku je dán graf funkce \( f \), přičemž  \( A \), \( B \) a \( C \) jsou body grafu této funkce. Jestliže \( x_A \), \( x_B \) a \( x_C \) jsou \( x \)-ové souřadnice bodů  \( A \), \( B \) a \( C \), a jestliže \( f' \) je derivace funkce \( f \), pak:}
{\( f'( x_A ) < f'( x_B ) < f'( x_C ) \)}{\( f'( x_A ) < f'( x_B ) = f'( x_C ) \)}{\( f'( x_A ) > f'( x_B ) = f'( x_C ) \)}{\( f'( x_A ) > f'( x_B ) > f'( x_C ) \)}{\( f'( x_A ) = f'( x_B ) < f'( x_C ) \)}{}}
\LANGsk{
\Question{
Na obrázku je daný graf funkcie \( f \), pričom  \( A \), \( B \) a \( C \) sú body grafu tejto funkcie. Ak \( x_A \), \( x_B \) a \( x_C \) sú \( x \)-ové súradnice bodov  \( A \), \( B \) a \( C \), a ak \( f' \) je derivácia funkcie \( f \), potom:}
{\( f'( x_A ) < f'( x_B ) < f'( x_C ) \)}{\( f'( x_A ) < f'( x_B ) = f'( x_C ) \)}{\( f'( x_A ) > f'( x_B ) = f'( x_C ) \)}{\( f'( x_A ) > f'( x_B ) > f'( x_C ) \)}{\( f'( x_A ) = f'( x_B ) < f'( x_C ) \)}{}}
\LANGes{
\Question{
Dada la gráfica de la función \( f \), donde \( A \), \( B \) y \( C \) son puntos de la gráfica. Si \( x_A \), \( x_B \) yd \( x_C \) denotan las coordenadas del eje \( x \) de los puntos \( A \), \( B \) y \( C \) y si \( f' \) es derivada de \( f \), entonces:}
{\( f'( x_A ) < f'( x_B ) < f'( x_C ) \)}{\( f'( x_A ) < f'( x_B ) = f'( x_C ) \)}{\( f'( x_A ) > f'( x_B ) = f'( x_C ) \)}{\( f'( x_A ) > f'( x_B ) > f'( x_C ) \)}{\( f'( x_A ) = f'( x_B ) < f'( x_C ) \)}{}}
\End

\Data\ID{1103164704}
\Updated{2020-07-15 03:07:12}
\Level{A}
\SubArea{Derivative}
\IMG{1103164704.pdf}
\LANGen{
\Question{
The graph of \( f \) is given in the figure, where \( A \), \( B \) and \( C \) are points of the graph, and \( y \)-coordinate of the point \( B \) is the minimum value of the function \( f \). If \( x_A \), \( x_B \) and \( x_C \) denote the \( x \)-coordinates of the points \( A \), \( B \) and \( C \), and if \( f' \) is the derivative of \( f \), then:}
{\( f'( x_A ) < 0 \), \( f'( x_B ) = 0 \), \( f'( x_C ) > 0 \)}{\( f'( x_A ) < 0 \), \( f'( x_B ) = 0 \), \( f'( x_C ) < 0 \)}{\( f'( x_A ) > 0 \), \( f'( x_B ) < 0 \), \( f'( x_C ) < 0 \)}{\( f'( x_A ) > 0 \), \( f'( x_B ) = 0 \), \( f'( x_C ) > 0 \)}{}{}}
\LANGpl{
\Question{
Dany jest wykres funkcji  \( f \), gdzie \( A \), \( B \) i \( C \) to punkty na wykresie, a \( y \)-współrzędna punktu \( B \) to minimalna wartość funkcji \( f \). Jeśli \( x_A \), \( x_B \) i \( x_C \) oznaczają \( x \)-współrzędne  punktów \( A \), \( B \) i \( C \) oraz  \( f' \) jest pochodną  \( f \), to:}
{\( f'( x_A ) < 0 \), \( f'( x_B ) = 0 \), \( f'( x_C ) > 0 \)}{\( f'( x_A ) < 0 \), \( f'( x_B ) = 0 \), \( f'( x_C ) < 0 \)}{\( f'( x_A ) > 0 \), \( f'( x_B ) < 0 \), \( f'( x_C ) < 0 \)}{\( f'( x_A ) > 0 \), \( f'( x_B ) = 0 \), \( f'( x_C ) > 0 \)}{}{}}
\LANGcs{
\Question{
Na obrázku je dán graf funkce \( f \), přičemž  \( A \), \( B \) a \( C \) jsou body grafu této funkce a \( y \)-ová souřadnice bodu \( B \) je nejmenší hodnota funkce \( f \). Jestliže \( x_A \), \( x_B \) a \( x_C \) jsou \( x \)-ové souřadnice bodů  \( A \), \( B \) a \( C \), a jestliže \( f' \) je derivace funkce \( f \), pak:}
{\( f'( x_A ) < 0 \), \( f'( x_B ) = 0 \), \( f'( x_C ) > 0 \)}{\( f'( x_A ) < 0 \), \( f'( x_B ) = 0 \), \( f'( x_C ) < 0 \)}{\( f'( x_A ) > 0 \), \( f'( x_B ) < 0 \), \( f'( x_C ) < 0 \)}{\( f'( x_A ) > 0 \), \( f'( x_B ) = 0 \), \( f'( x_C ) > 0 \)}{}{}}
\LANGsk{
\Question{
Na obrázku je daný graf funkcie \( f \), pričom  \( A \), \( B \) a \( C \) sú body grafu tejto funkcie a \( y \)-ová súradnice bodu \( B \) je najmenšia hodnota funkcie \( f \). Ak \( x_A \), \( x_B \) a \( x_C \) sú \( x \)-ové súradnice bodov  \( A \), \( B \) a \( C \), a ak \( f' \) je derivácia funkcie \( f \), potom:}
{\( f'( x_A ) < 0 \), \( f'( x_B ) = 0 \), \( f'( x_C ) > 0 \)}{\( f'( x_A ) < 0 \), \( f'( x_B ) = 0 \), \( f'( x_C ) < 0 \)}{\( f'( x_A ) > 0 \), \( f'( x_B ) < 0 \), \( f'( x_C ) < 0 \)}{\( f'( x_A ) > 0 \), \( f'( x_B ) = 0 \), \( f'( x_C ) > 0 \)}{}{}}
\LANGes{
\Question{
Dada la gráfica de la función \( f \), donde \( A \), \( B \) y \( C \) son puntos de la gráfica, y la coordenada del eje \( y \) del punto \( B \) es el mínimo valor de la función \( f \). Si \( x_A \), \( x_B \) y \( x_C \) denotan las coordenadas del eje \( x \) de los puntos \( A \), \( B \) y \( C \), y si \( f' \) es derivada de \( f \), entonces:}
{\( f'( x_A ) < 0 \), \( f'( x_B ) = 0 \), \( f'( x_C ) > 0 \)}{\( f'( x_A ) < 0 \), \( f'( x_B ) = 0 \), \( f'( x_C ) < 0 \)}{\( f'( x_A ) > 0 \), \( f'( x_B ) < 0 \), \( f'( x_C ) < 0 \)}{\( f'( x_A ) > 0 \), \( f'( x_B ) = 0 \), \( f'( x_C ) > 0 \)}{}{}}
\End

\Data\ID{1103164705}
\Updated{2020-07-15 03:07:12}
\Level{A}
\SubArea{Derivative}
\IMG{1103164705.pdf}
\LANGen{
\Question{
The graph of \( f \) is given in the figure. Which of the following holds? (\( f' \) is the derivative of the function \( f \).)}
{\( f'(-1)=1 \), \( f'(2)=0 \), \( f'(4)=-2 \)}{\( f'(-1)=1 \), \( f'(2)=2 \), \( f'(4)=0 \)}{\( f'(-1)=0 \), \( f'(2)=0 \), \( f'(4)=-2 \)}{\( f'(-1)=0 \), \( f'(2)=2 \), \( f'(4)=0 \)}{\( f'(-1)=1 \), \( f'(2)=0 \), \( f'(4)=0 \)}{}}
\LANGpl{
\Question{
Dany jest wykres funkcji \( f \). Które z podanych wyrażeń jest poprawne? (\( f' \) jest pochodną funkcji \( f \).)}
{\( f'(-1)=1 \), \( f'(2)=0 \), \( f'(4)=-2 \)}{\( f'(-1)=1 \), \( f'(2)=2 \), \( f'(4)=0 \)}{\( f'(-1)=0 \), \( f'(2)=0 \), \( f'(4)=-2 \)}{\( f'(-1)=0 \), \( f'(2)=2 \), \( f'(4)=0 \)}{\( f'(-1)=1 \), \( f'(2)=0 \), \( f'(4)=0 \)}{}}
\LANGcs{
\Question{
Na obrázku je dán graf funkce \( f \). Které z následujících tvrzení platí? (\( f' \) je derivace funkce \( f \).)}
{\( f'(-1)=1 \), \( f'(2)=0 \), \( f'(4)=-2 \)}{\( f'(-1)=1 \), \( f'(2)=2 \), \( f'(4)=0 \)}{\( f'(-1)=0 \), \( f'(2)=0 \), \( f'(4)=-2 \)}{\( f'(-1)=0 \), \( f'(2)=2 \), \( f'(4)=0 \)}{\( f'(-1)=1 \), \( f'(2)=0 \), \( f'(4)=0 \)}{}}
\LANGsk{
\Question{
Na obrázku je daný graf funkcie \( f \). Ktoré z nasledujúcich tvrdení platí? (\( f' \) je derivácia funkcie \( f \).)}
{\( f'(-1)=1 \), \( f'(2)=0 \), \( f'(4)=-2 \)}{\( f'(-1)=1 \), \( f'(2)=2 \), \( f'(4)=0 \)}{\( f'(-1)=0 \), \( f'(2)=0 \), \( f'(4)=-2 \)}{\( f'(-1)=0 \), \( f'(2)=2 \), \( f'(4)=0 \)}{\( f'(-1)=1 \), \( f'(2)=0 \), \( f'(4)=0 \)}{}}
\LANGes{
\Question{
Dada la gráfica de la función \( f \). ¿Cuál de los enunciados es correcto? (\( f' \) es derivada de la función \( f \).)}
{\( f'(-1)=1 \), \( f'(2)=0 \), \( f'(4)=-2 \)}{\( f'(-1)=1 \), \( f'(2)=2 \), \( f'(4)=0 \)}{\( f'(-1)=0 \), \( f'(2)=0 \), \( f'(4)=-2 \)}{\( f'(-1)=0 \), \( f'(2)=2 \), \( f'(4)=0 \)}{\( f'(-1)=1 \), \( f'(2)=0 \), \( f'(4)=0 \)}{}}
\End

\Data\ID{1103164706}
\Updated{2020-07-15 03:07:12}
\Level{A}
\SubArea{Derivative}
\IMG{1103164706.pdf}
\LANGen{
\Question{
The graph of \( f \) is given in the figure. Which of the following holds? (\( f' \) is the derivative of the function \( f \).)}
{\( f'(0)=1 \), \( f'(1) \) does not exist, \( f'(4)=-2 \)}{\( f'(0)=1 \), \( f'(1)=0 \), \( f'(4)=-2 \)}{\( f'(-1)=0 \), \( f'(2)=0 \), \( f'(3) \)  does not exist}{\( f'(-1)=1 \), \( f'(2)=0 \), \( f'(3)=0 \)}{}{}}
\LANGpl{
\Question{
Dany jest wykres funkcji \( f \). Które z podanych stwierdzeń jest poprawne? (\( f' \) jest pochodną funkcji \( f \).)}
{\( f'(0)=1 \), \( f'(1) \) nie istnieje, \( f'(4)=-2 \)}{\( f'(0)=1 \), \( f'(1)=0 \), \( f'(4)=-2 \)}{\( f'(-1)=0 \), \( f'(2)=0 \), \( f'(3) \)  nie istnieje}{\( f'(-1)=1 \), \( f'(2)=0 \), \( f'(3)=0 \)}{}{}}
\LANGcs{
\Question{
Na obrázku je dán graf funkce \( f \). Které z následujících tvrzení platí? (\( f' \) je derivace funkce \( f \).)}
{\( f'(0)=1 \), \( f'(1) \) neexistuje, \( f'(4)=-2 \)}{\( f'(0)=1 \), \( f'(1)=0 \), \( f'(4)=-2 \)}{\( f'(-1)=0 \), \( f'(2)=0 \), \( f'(3) \)  neexistuje}{\( f'(-1)=1 \), \( f'(2)=0 \), \( f'(3)=0 \)}{}{}}
\LANGsk{
\Question{
Na obrázku je daný graf funkcie \( f \). Ktoré z nasledujúcich tvrdení platí? (\( f' \) je derivácia funkcie \( f \).)}
{\( f'(0)=1 \), \( f'(1) \) neexistuje, \( f'(4)=-2 \)}{\( f'(0)=1 \), \( f'(1)=0 \), \( f'(4)=-2 \)}{\( f'(-1)=0 \), \( f'(2)=0 \), \( f'(3) \)  neexistuje}{\( f'(-1)=1 \), \( f'(2)=0 \), \( f'(3)=0 \)}{}{}}
\LANGes{
\Question{
Dada la gráfica de la función \( f \). ¿Cuál de las gráficas es la de \( f' \)? (\( f' \) es la derivada de la función \( f \).)}
{\( f'(0)=1 \), \( f'(1) \) no existe, \( f'(4)=-2 \)}{\( f'(0)=1 \), \( f'(1)=0 \), \( f'(4)=-2 \)}{\( f'(-1)=0 \), \( f'(2)=0 \), \( f'(3) \)  no existe}{\( f'(-1)=1 \), \( f'(2)=0 \), \( f'(3)=0 \)}{}{}}
\End

\Data\ID{1103164707}
\Updated{2020-07-15 03:07:12}
\Level{A}
\SubArea{Derivative}
\IMG{1103164707.pdf}
\LANGen{
\Question{
The graph of \( g \) is given in the figure. Which of the following holds? (\( g' \) is the derivative of the function \( g \).)}
{\( g'(-1)=0 \), \( g'(1)=2 \), \( g'(5)=-1 \)}{\( g'(-1)=0 \), \( g'(1)=1 \), \( g'(5)=-1 \)}{\( g'(-1)=0 \), \( g'(1)=2 \), \( g'(5)=2 \)}{\( g'(-1)=-2 \), \( g'(1)=0 \), \( g'(5)=2 \)}{\( g'(-1)=-2 \), \( g'(1)=2 \), \( g'(5)=2 \)}{}}
\LANGpl{
\Question{
Dany jest wykres funkcji \( g \). Które z podanych wyrażeń jest poprawne?  (\( g' \) jest pochodną funkcji \( g \).)}
{\( g'(-1)=0 \), \( g'(1)=2 \), \( g'(5)=-1 \)}{\( g'(-1)=0 \), \( g'(1)=1 \), \( g'(5)=-1 \)}{\( g'(-1)=0 \), \( g'(1)=2 \), \( g'(5)=2 \)}{\( g'(-1)=-2 \), \( g'(1)=0 \), \( g'(5)=2 \)}{\( g'(-1)=-2 \), \( g'(1)=2 \), \( g'(5)=2 \)}{}}
\LANGcs{
\Question{
Na obrázku je dán graf funkce \( g \). Které z následujících tvrzení platí? (\( g' \) je derivace funkce \( g \).)}
{\( g'(-1)=0 \), \( g'(1)=2 \), \( g'(5)=-1 \)}{\( g'(-1)=0 \), \( g'(1)=1 \), \( g'(5)=-1 \)}{\( g'(-1)=0 \), \( g'(1)=2 \), \( g'(5)=2 \)}{\( g'(-1)=-2 \), \( g'(1)=0 \), \( g'(5)=2 \)}{\( g'(-1)=-2 \), \( g'(1)=2 \), \( g'(5)=2 \)}{}}
\LANGsk{
\Question{
Na obrázku je daný graf funkcie \( g \). Ktoré z nasledujúcich tvrdení platí? (\( g' \) je derivácie funkcie \( g \).)}
{\( g'(-1)=0 \), \( g'(1)=2 \), \( g'(5)=-1 \)}{\( g'(-1)=0 \), \( g'(1)=1 \), \( g'(5)=-1 \)}{\( g'(-1)=0 \), \( g'(1)=2 \), \( g'(5)=2 \)}{\( g'(-1)=-2 \), \( g'(1)=0 \), \( g'(5)=2 \)}{\( g'(-1)=-2 \), \( g'(1)=2 \), \( g'(5)=2 \)}{}}
\LANGes{
\Question{
Dada la gráfica de la función \( g \). ¿Cuál de los enunciados es correcto? (\( g' \) es derivada de la función \( g \).)}
{\( g'(-1)=0 \), \( g'(1)=2 \), \( g'(5)=-1 \)}{\( g'(-1)=0 \), \( g'(1)=1 \), \( g'(5)=-1 \)}{\( g'(-1)=0 \), \( g'(1)=2 \), \( g'(5)=2 \)}{\( g'(-1)=-2 \), \( g'(1)=0 \), \( g'(5)=2 \)}{\( g'(-1)=-2 \), \( g'(1)=2 \), \( g'(5)=2 \)}{}}
\End

\Data\ID{1103164708}
\Updated{2020-07-15 03:07:12}
\Level{A}
\SubArea{Derivative}
\IMG{1103164708.pdf}
\LANGen{
\Question{
The graph of \( g \) is given in the figure. Which of the following holds? (\( g' \) is the derivative of the function \( g \).)}
{\( g'(1)=2 \), \( g'(2)=2 \), \( g'(4)=-1 \)}{\( g'(-1)=0 \), \( g'(3) \) does not exist, \( g'(4)=3 \)}{\( g'(-1) = -2 \), \( g'(3) \) does not exist, \( g'(4)=-1 \)}{\( g'(1)=0 \), \( g'(2)=2 \), \( g'(4)=3 \)}{}{}}
\LANGpl{
\Question{
Dany jest wykres funkcji \( g \). Które z podanych stwierdzeń jest poprawne? (\( g' \) jest pochodną funkcji  \( g \).)}
{\( g'(1)=2 \), \( g'(2)=2 \), \( g'(4)=-1 \)}{\( g'(-1)=0 \), \( g'(3) \) nie istnieje, \( g'(4)=3 \)}{\( g'(-1) = -2 \), \( g'(3) \) nie istnieje, \( g'(4)=-1 \)}{\( g'(1)=0 \), \( g'(2)=2 \), \( g'(4)=3 \)}{}{}}
\LANGcs{
\Question{
Na obrázku je dán graf funkce \( g \). Které z následujících tvrzení platí? (\( g' \) je derivace funkce \( g \).)}
{\( g'(1)=2 \), \( g'(2)=2 \), \( g'(4)=-1 \)}{\( g'(-1)=0 \), \( g'(3) \) neexistuje, \( g'(4)=3 \)}{\( g'(-1) = -2 \), \( g'(3) \) neexistuje, \( g'(4)=-1 \)}{\( g'(1)=0 \), \( g'(2)=2 \), \( g'(4)=3 \)}{}{}}
\LANGsk{
\Question{
Na obrázku je daný graf funkcie \( g \). Ktoré z nasledujúcich tvrdení platí? (\( g' \) je derivácie funkcie \( g \).)}
{\( g'(1)=2 \), \( g'(2)=2 \), \( g'(4)=-1 \)}{\( g'(-1)=0 \), \( g'(3) \) neexistuje, \( g'(4)=3 \)}{\( g'(-1) = -2 \), \( g'(3) \) neexistuje, \( g'(4)=-1 \)}{\( g'(1)=0 \), \( g'(2)=2 \), \( g'(4)=3 \)}{}{}}
\LANGes{
\Question{
Dada la gráfica de la función \( g \). Elige cuál de los siguientes enunciados es válido. (la función \( g' \) es la derivada de \( g \)).}
{\( g'(1)=2 \), \( g'(2)=2 \), \( g'(4)=-1 \)}{\( g'(-1)=0 \), \( g'(3) \) does not exist, \( g'(4)=3 \)}{\( g'(-1) = -2 \), \( g'(3) \) does not exist, \( g'(4)=-1 \)}{\( g'(1)=0 \), \( g'(2)=2 \), \( g'(4)=3 \)}{}{}}
\End

\Data\ID{1103164709}
\Updated{2020-07-15 03:07:12}
\Level{A}
\SubArea{Derivative}
\IMG{1103164709.pdf}
\LANGen{
\Question{
The graph of \( f \) is given in the figure. Which of the following holds? (\( f' \) is the derivative of the function \( f \).)}
{\( f'(-1) \) does not exist, \( f'(1)=-\frac12 \), \( f'(4)=3 \)}{\( f'(-2)=0 \), \( f'(1)=-2 \), \( f'(3) \) does not exist}{\( f'(-2)=0 \), \( f'(2)=-\frac32 \), \( f'(4)=3 \)}{\( f'(-1)=2 \), \( f'(1)=-\frac12 \), \( f'(3) \)  does not exist}{}{}}
\LANGpl{
\Question{
Dany jest wykres funkcji \( f \). Które z podanych stwierdzeń jest poprawne? (\( f' \) jest pochodną funkcji \( f \).)}
{\( f'(-1) \) nie istnieje, \( f'(1)=-\frac12 \), \( f'(4)=3 \)}{\( f'(-2)=0 \), \( f'(1)=-2 \), \( f'(3) \) nie istnieje}{\( f'(-2)=0 \), \( f'(2)=-\frac32 \), \( f'(4)=3 \)}{\( f'(-1)=2 \), \( f'(1)=-\frac12 \), \( f'(3) \)  nie istnieje}{}{}}
\LANGcs{
\Question{
Na obrázku je dán graf funkce \( f \). Které z následujících tvrzení platí? (\( f' \) je derivace funkce \( f \).)}
{\( f'(-1) \) neexistuje, \( f'(1)=-\frac12 \), \( f'(4)=3 \)}{\( f'(-2)=0 \), \( f'(1)=-2 \), \( f'(3) \) neexistuje}{\( f'(-2)=0 \), \( f'(2)=-\frac32 \), \( f'(4)=3 \)}{\( f'(-1)=2 \), \( f'(1)=-\frac12 \), \( f'(3) \)  neexistuje}{}{}}
\LANGsk{
\Question{
Na obrázku je daný graf funkcie \( f \). Ktoré z nasledujúcich tvrdení platí? (\( f' \) je derivácia funkcie \( f \).)}
{\( f'(-1) \) neexistuje, \( f'(1)=-\frac12 \), \( f'(4)=3 \)}{\( f'(-2)=0 \), \( f'(1)=-2 \), \( f'(3) \) neexistuje}{\( f'(-2)=0 \), \( f'(2)=-\frac32 \), \( f'(4)=3 \)}{\( f'(-1)=2 \), \( f'(1)=-\frac12 \), \( f'(3) \)  neexistuje}{}{}}
\LANGes{
\Question{
Dada la gráfica de la función \( f \). ¿Cuál de los enunciados es correcto? (\( f' \) es derivada de la función \( f \).)}
{\( f'(-1) \) no existe, \( f'(1)=-\frac12 \), \( f'(4)=3 \)}{\( f'(-2)=0 \), \( f'(1)=-2 \), \( f'(3) \) no existe}{\( f'(-2)=0 \), \( f'(2)=-\frac32 \), \( f'(4)=3 \)}{\( f'(-1)=2 \), \( f'(1)=-\frac12 \), \( f'(3) \)  no existe}{}{}}
\End

\Data\ID{1103164710}
\Updated{2020-07-15 03:07:12}
\Level{A}
\SubArea{Derivative}
\IMG{1103164710.pdf}
\LANGen{
\Question{
The graph of \( f \) is given in the figure. Which of the following holds? (\( f' \) is the derivative of the function \( f \).)}
{\( f'(-1)=0 \), \( f'(2)=-2 \), \( f'(4)=0 \)}{\( f'(-1)=0 \), \( f'(2)=-\frac12 \), \( f'(4)=-3 \)}{\( f'(-1)=0 \), \( f'(0)=0 \), \( f'(3)=0 \)}{\( f'(0)=0 \), \( f'(2)=-1 \), \( f'(4)=0 \)}{\( f'(0)=1 \), \( f'(2)=-2 \), \( f'(4)=-3 \)}{}}
\LANGpl{
\Question{
Dany jest wykres funkcji \( f \). Które z podanych wyrażeń jest poprawne?  (\( f' \) jest pochodną funkcji \( f \).)}
{\( f'(-1)=0 \), \( f'(2)=-2 \), \( f'(4)=0 \)}{\( f'(-1)=0 \), \( f'(2)=-\frac12 \), \( f'(4)=-3 \)}{\( f'(-1)=0 \), \( f'(0)=0 \), \( f'(3)=0 \)}{\( f'(0)=0 \), \( f'(2)=-1 \), \( f'(4)=0 \)}{\( f'(0)=1 \), \( f'(2)=-2 \), \( f'(4)=-3 \)}{}}
\LANGcs{
\Question{
Na obrázku je dán graf funkce \( f \). Které z následujících tvrzení platí? (\( f' \) je derivace funkce \( f \).)}
{\( f'(-1)=0 \), \( f'(2)=-2 \), \( f'(4)=0 \)}{\( f'(-1)=0 \), \( f'(2)=-\frac12 \), \( f'(4)=-3 \)}{\( f'(-1)=0 \), \( f'(0)=0 \), \( f'(3)=0 \)}{\( f'(0)=0 \), \( f'(2)=-1 \), \( f'(4)=0 \)}{\( f'(0)=1 \), \( f'(2)=-2 \), \( f'(4)=-3 \)}{}}
\LANGsk{
\Question{
Na obrázku je daný graf funkcie \( f \). Ktoré z nasledujúcich tvrdení platí? (\( f' \) je derivácia funkcie \( f \).)}
{\( f'(-1)=0 \), \( f'(2)=-2 \), \( f'(4)=0 \)}{\( f'(-1)=0 \), \( f'(2)=-\frac12 \), \( f'(4)=-3 \)}{\( f'(-1)=0 \), \( f'(0)=0 \), \( f'(3)=0 \)}{\( f'(0)=0 \), \( f'(2)=-1 \), \( f'(4)=0 \)}{\( f'(0)=1 \), \( f'(2)=-2 \), \( f'(4)=-3 \)}{}}
\LANGes{
\Question{
Dada la gráfica de la función \( f \). ¿Cuál de los enunciados es correcto? (\( f' \) es derivada de la función \( f \).)}
{\( f'(-1)=0 \), \( f'(2)=-2 \), \( f'(4)=0 \)}{\( f'(-1)=0 \), \( f'(2)=-\frac12 \), \( f'(4)=-3 \)}{\( f'(-1)=0 \), \( f'(0)=0 \), \( f'(3)=0 \)}{\( f'(0)=0 \), \( f'(2)=-1 \), \( f'(4)=0 \)}{\( f'(0)=1 \), \( f'(2)=-2 \), \( f'(4)=-3 \)}{}}
\End

\Data\ID{9000062401}
\Updated{2020-07-15 03:07:37}
\Level{A}
\SubArea{Derivative}
\LANGen{
\Question{
Find the derivative of \(f(x) = 3x^{4} - 2x^{3} - 3x^{2}\) 
at the point \(x_{0} = -1\).}
{\(- 12\)}{\(0\)}{\(12\)}{\(24\)}{}{}}
\LANGpl{
\Question{
Wyznacz pochodną funkcji  \(f\colon y = 3x^{4} - 2x^{3} - 3x^{2}\) 
w punkcie  \(x_{0} = -1\).}
{\(- 12\)}{\(0\)}{\(12\)}{\(24\)}{}{}}
\LANGcs{
\Question{
Určete první derivaci funkce \(f\colon y = 3x^{4} - 2x^{3} - 3x^{2}\) 
v bodě \(x_{0} = -1\).}
{\(- 12\)}{\(0\)}{\(12\)}{\(24\)}{}{}}
\LANGsk{
\Question{
Určte prvú deriváciu funkcie \(f\colon y = 3x^{4} - 2x^{3} - 3x^{2}\) 
v bode \(x_{0} = -1\).}
{\(- 12\)}{\(0\)}{\(12\)}{\(24\)}{}{}}
\LANGes{
\Question{
Halla la derivada de \(f(x) = 3x^{4} - 2x^{3} - 3x^{2}\) 
en el punto \(x_{0} = -1\).}
{\(- 12\)}{\(0\)}{\(12\)}{\(24\)}{}{}}
\End

\Data\ID{9000062402}
\Updated{2020-07-15 03:07:37}
\Level{A}
\SubArea{Derivative}
\LANGen{
\Question{
Find the second derivative of \(f(x) = x^{4} - 3x^{2}\) 
at the point \(x_{0} = 1\).}
{\(6\)}{\(- 2\)}{\(- 6\)}{\(1\)}{}{}}
\LANGpl{
\Question{
Wyznacz drugą pochodną funkcji \(f\colon y = x^{4} - 3x^{2}\) 
w punkcie \(x_{0} = 1\).}
{\(6\)}{\(- 2\)}{\(- 6\)}{\(1\)}{}{}}
\LANGcs{
\Question{
Určete druhou derivaci funkce \(f\colon y = x^{4} - 3x^{2}\) 
v bodě \(x_{0} = 1\).}
{\(6\)}{\(- 2\)}{\(- 6\)}{\(1\)}{}{}}
\LANGsk{
\Question{
Určte druhú deriváciu funkcie \(f\colon y = x^{4} - 3x^{2}\) 
v bode \(x_{0} = 1\).}
{\(6\)}{\(- 2\)}{\(- 6\)}{\(1\)}{}{}}
\LANGes{
\Question{
Calcula la segunda derivada de \(f(x) = x^{4} - 3x^{2}\) 
en el punto \(x_{0} = 1\).}
{\(6\)}{\(- 2\)}{\(- 6\)}{\(1\)}{}{}}
\End

\Data\ID{9000070802}
\Updated{2020-07-15 03:07:37}
\Level{A}
\SubArea{Derivative}
\LANGen{
\Question{
Differentiate the following function. 
\[ 
 f(x) = 3 - 2\cos x
\]}
{\(f'(x) = 2\sin x;\ x\in \mathbb{R}\)}{\(f'(x) = 3 + 2\sin x;\ x\in \mathbb{R}\)}{\(f'(x) = 3 - 2\sin x;\ x\in \mathbb{R}\)}{\(f'(x) = 2\cos x;\ x\in \mathbb{R}\)}{}{}}
\LANGpl{
\Question{
Wyznacz pochodną funkcji. 
\[ 
 f\colon y = 3 - 2\cos x
\]}
{\(f'(x) = 2\sin x;\ x\in \mathbb{R}\)}{\(f'(x) = 3 + 2\sin x;\ x\in \mathbb{R}\)}{\(f'(x) = 3 - 2\sin x;\ x\in \mathbb{R}\)}{\(f'(x) = 2\cos x;\ x\in \mathbb{R}\)}{}{}}
\LANGcs{
\Question{
Určete první derivaci funkce \(f\colon y = 3 - 2\cos x\).}
{\(f'(x) = 2\sin x;\ x\in \mathbb{R}\)}{\(f'(x) = 3 + 2\sin x;\ x\in \mathbb{R}\)}{\(f'(x) = 3 - 2\sin x;\ x\in \mathbb{R}\)}{\(f'(x) = 2\cos x;\ x\in \mathbb{R}\)}{}{}}
\LANGsk{
\Question{
Určte prvú deriváciu funkcie \(f\colon y = 3 - 2\cos x\).}
{\(f'(x) = 2\sin x;\ x\in \mathbb{R}\)}{\(f'(x) = 3 + 2\sin x;\ x\in \mathbb{R}\)}{\(f'(x) = 3 - 2\sin x;\ x\in \mathbb{R}\)}{\(f'(x) = 2\cos x;\ x\in \mathbb{R}\)}{}{}}
\LANGes{
\Question{
Deriva la siguiente función.
\[ 
 f(x) = 3 - 2\cos x
\]}
{\(f'(x) = 2\sin x;\ x\in \mathbb{R}\)}{\(f'(x) = 3 + 2\sin x;\ x\in \mathbb{R}\)}{\(f'(x) = 3 - 2\sin x;\ x\in \mathbb{R}\)}{\(f'(x) = 2\cos x;\ x\in \mathbb{R}\)}{}{}}
\End

\Data\ID{9000070803}
\Updated{2020-07-15 03:07:37}
\Level{A}
\SubArea{Derivative}
\LANGen{
\Question{
Differentiate the following function. 
\[ 
 f(x) = 3x^{3} + 2x +\mathrm{e} ^{x}
\]}
{\(f'(x) = 9x^{2} + 2 +\mathrm{e} ^{x};\ x\in \mathbb{R}\)}{\(f'(x) = 6x^{2} + 2x;\ x\in \mathbb{R}\)}{\(f'(x) = 6x^{2} + 2x +\mathrm{e} ^{x};\ x\in \mathbb{R}\)}{\(f'(x) = 9x^{2} + 2;\ x\in \mathbb{R}\)}{}{}}
\LANGpl{
\Question{
Wyznacz pochodną funkcji.
\[ 
 f\colon y = 3x^{3} + 2x +\mathrm{e} ^{x}
\]}
{\(f'(x) = 9x^{2} + 2 +\mathrm{e} ^{x};\ x\in \mathbb{R}\)}{\(f'(x) = 6x^{2} + 2x;\ x\in \mathbb{R}\)}{\(f'(x) = 6x^{2} + 2x +\mathrm{e} ^{x};\ x\in \mathbb{R}\)}{\(f'(x) = 9x^{2} + 2;\ x\in \mathbb{R}\)}{}{}}
\LANGcs{
\Question{
Určete první derivaci funkce \(f\colon y = 3x^{3} + 2x +\mathrm{e} ^{x}\).}
{\(f'(x) = 9x^{2} + 2 +\mathrm{e} ^{x};\ x\in \mathbb{R}\)}{\(f'(x) = 6x^{2} + 2x;\ x\in \mathbb{R}\)}{\(f'(x) = 6x^{2} + 2x +\mathrm{e} ^{x};\ x\in \mathbb{R}\)}{\(f'(x) = 9x^{2} + 2;\ x\in \mathbb{R}\)}{}{}}
\LANGsk{
\Question{
Určte prvú deriváciu funkcie \(f\colon y = 3x^{3} + 2x +\mathrm{e} ^{x}\).}
{\(f'(x) = 9x^{2} + 2 +\mathrm{e} ^{x};\ x\in \mathbb{R}\)}{\(f'(x) = 6x^{2} + 2x;\ x\in \mathbb{R}\)}{\(f'(x) = 6x^{2} + 2x +\mathrm{e} ^{x};\ x\in \mathbb{R}\)}{\(f'(x) = 9x^{2} + 2;\ x\in \mathbb{R}\)}{}{}}
\LANGes{
\Question{
Deriva la siguiente función.
\[ 
 f(x) = 3x^{3} + 2x +\mathrm{e} ^{x}
\]}
{\(f'(x) = 9x^{2} + 2 +\mathrm{e} ^{x};\ x\in \mathbb{R}\)}{\(f'(x) = 6x^{2} + 2x;\ x\in \mathbb{R}\)}{\(f'(x) = 6x^{2} + 2x +\mathrm{e} ^{x};\ x\in \mathbb{R}\)}{\(f'(x) = 9x^{2} + 2;\ x\in \mathbb{R}\)}{}{}}
\End

\Data\ID{9000070804}
\Updated{2020-07-15 03:07:37}
\Level{A}
\SubArea{Derivative}
\LANGen{
\Question{
Differentiate the following function. 
\[ 
 f(x) = 2x^{9} - x^{2} + 7
\]}
{\(f'(x) = 18x^{8} - 2x;\ x\in \mathbb{R}\)}{\(f'(x) = 9x^{8} - 2x + 7;\ x\in \mathbb{R}\)}{\(f'(x) = 18x^{8} - 2x + 7;\ x\in \mathbb{R}\)}{\(f'(x) = 18x^{8} + 2x;\ x\in \mathbb{R}\)}{}{}}
\LANGpl{
\Question{
Wyznacz pochodną funkcji.
\[ 
 f\colon y = 2x^{9} - x^{2} + 7
\]}
{\(f'(x) = 18x^{8} - 2x;\ x\in \mathbb{R}\)}{\(f'(x) = 9x^{8} - 2x + 7;\ x\in \mathbb{R}\)}{\(f'(x) = 18x^{8} - 2x + 7;\ x\in \mathbb{R}\)}{\(f'(x) = 18x^{8} + 2x;\ x\in \mathbb{R}\)}{}{}}
\LANGcs{
\Question{
Určete první derivaci funkce \(f\colon y = 2x^{9} - x^{2} + 7\).}
{\(f'(x) = 18x^{8} - 2x;\ x\in \mathbb{R}\)}{\(f'(x) = 9x^{8} - 2x + 7;\ x\in \mathbb{R}\)}{\(f'(x) = 18x^{8} - 2x + 7;\ x\in \mathbb{R}\)}{\(f'(x) = 18x^{8} + 2x;\ x\in \mathbb{R}\)}{}{}}
\LANGsk{
\Question{
Určte prvú deriváciu funkcie \(f\colon y = 2x^{9} - x^{2} + 7\).}
{\(f'(x) = 18x^{8} - 2x;\ x\in \mathbb{R}\)}{\(f'(x) = 9x^{8} - 2x + 7;\ x\in \mathbb{R}\)}{\(f'(x) = 18x^{8} - 2x + 7;\ x\in \mathbb{R}\)}{\(f'(x) = 18x^{8} + 2x;\ x\in \mathbb{R}\)}{}{}}
\LANGes{
\Question{
Deriva la siguiente función.
\[ 
 f(x) = 2x^{9} - x^{2} + 7
\]}
{\(f'(x) = 18x^{8} - 2x;\ x\in \mathbb{R}\)}{\(f'(x) = 9x^{8} - 2x + 7;\ x\in \mathbb{R}\)}{\(f'(x) = 18x^{8} - 2x + 7;\ x\in \mathbb{R}\)}{\(f'(x) = 18x^{8} + 2x;\ x\in \mathbb{R}\)}{}{}}
\End

\Data\ID{9000070805}
\Updated{2020-07-15 03:07:37}
\Level{A}
\SubArea{Derivative}
\LANGen{
\Question{
Differentiate the following function. 
\[ 
 f(x) = -3x^{3} - x^{2} + 9x
\]}
{\(f'(x) = -9x^{2} - 2x + 9;\ x\in \mathbb{R}\)}{\(f'(x) = 9x^{2} - 2x + 9;\ x\in \mathbb{R}\)}{\(f'(x) = 27x^{2} - 2x;\ x\in \mathbb{R}\)}{\(f'(x) = -9x^{2} - 2x;\ x\in \mathbb{R}\)}{}{}}
\LANGpl{
\Question{
Wyznacz pochodną funkcji.
\[ 
 f\colon y = -3x^{3} - x^{2} + 9x
\]}
{\(f'(x) = -9x^{2} - 2x + 9;\ x\in \mathbb{R}\)}{\(f'(x) = 9x^{2} - 2x + 9;\ x\in \mathbb{R}\)}{\(f'(x) = 27x^{2} - 2x;\ x\in \mathbb{R}\)}{\(f'(x) = -9x^{2} - 2x;\ x\in \mathbb{R}\)}{}{}}
\LANGcs{
\Question{
Určete první derivaci funkce \(f\colon y = -3x^{3} - x^{2} + 9x\).}
{\(f'(x) = -9x^{2} - 2x + 9;\ x\in \mathbb{R}\)}{\(f'(x) = 9x^{2} - 2x + 9;\ x\in \mathbb{R}\)}{\(f'(x) = 27x^{2} - 2x;\ x\in \mathbb{R}\)}{\(f'(x) = -9x^{2} - 2x;\ x\in \mathbb{R}\)}{}{}}
\LANGsk{
\Question{
Určte prvú deriváciu funkcie \(f\colon y = -3x^{3} - x^{2} + 9x\).}
{\(f'(x) = -9x^{2} - 2x + 9;\ x\in \mathbb{R}\)}{\(f'(x) = 9x^{2} - 2x + 9;\ x\in \mathbb{R}\)}{\(f'(x) = 27x^{2} - 2x;\ x\in \mathbb{R}\)}{\(f'(x) = -9x^{2} - 2x;\ x\in \mathbb{R}\)}{}{}}
\LANGes{
\Question{
Deriva la siguiente función.
\[ 
 f(x) = -3x^{3} - x^{2} + 9x
\]}
{\(f'(x) = -9x^{2} - 2x + 9;\ x\in \mathbb{R}\)}{\(f'(x) = 9x^{2} - 2x + 9;\ x\in \mathbb{R}\)}{\(f'(x) = 27x^{2} - 2x;\ x\in \mathbb{R}\)}{\(f'(x) = -9x^{2} - 2x;\ x\in \mathbb{R}\)}{}{}}
\End

\Data\ID{9000070806}
\Updated{2021-01-31 11:01:46}
\Level{A}
\SubArea{Derivative}
\LANGen{
\Question{
Differentiate the following function. 
\[ 
 f(x) = \frac{\pi } 
{x} +\ln 2
\]}
{\(f'(x) = - \frac{\pi }{x^{2}} ;\ x\in \mathbb{R}\setminus \{0\}\)}{\(f'(x) = 0;\ x\in \mathbb{R}\setminus \{0\}\)}{\(f'(x) =\pi ;\ x\in \mathbb{R}\setminus \{0\}\)}{\(f'(x) = \frac{\pi } {x^{2}} ;\ x\in \mathbb{R}\setminus \{0\}\)}{}{}}
\LANGpl{
\Question{
Wyznacz pochodną funkcji.
\[ 
 f\colon y = \frac{\pi } 
{x} +\ln 2
\]}
{\(f'(x) = - \frac{\pi }{x^{2}} ;\ x\in \mathbb{R}\setminus \{0\}\)}{\(f'(x) = 0;\ x\in \mathbb{R}\setminus \{0\}\)}{\(f'(x) =\pi ;\ x\in \mathbb{R}\setminus \{0\}\)}{\(f'(x) = \frac{\pi } {x^{2}} ;\ x\in \mathbb{R}\setminus \{0\}\)}{}{}}
\LANGcs{
\Question{
Určete první derivaci funkce \(f\colon y = \frac{\pi } {x} +\ln 2\).}
{\(f'(x) = - \frac{\pi }{x^{2}} ;\ x\in \mathbb{R}\setminus \{0\}\)}{\(f'(x) = 0;\ x\in \mathbb{R}\setminus \{0\}\)}{\(f'(x) =\pi ;\ x\in \mathbb{R}\setminus \{0\}\)}{\(f'(x) = \frac{\pi } {x^{2}} ;\ x\in \mathbb{R}\setminus \{0\}\)}{}{}}
\LANGsk{
\Question{
Určte prvú deriváciu funkcie \(f\colon y = \frac{\pi } {x} +\ln 2\).}
{\(f'(x) = - \frac{\pi }{x^{2}} ;\ x\in \mathbb{R}\setminus \{0\}\)}{\(f'(x) = 0;\ x\in \mathbb{R}\setminus \{0\}\)}{\(f'(x) =\pi ;\ x\in \mathbb{R}\setminus \{0\}\)}{\(f'(x) = \frac{\pi } {x^{2}} ;\ x\in \mathbb{R}\setminus \{0\}\)}{}{}}
\LANGes{
\Question{
Deriva la siguiente función. \[ 
 f(x) = \frac{\pi } 
{x} +\ln 2
\]}
{\(f'(x) = - \frac{\pi }{x^{2}} ;\ x\in \mathbb{R}\setminus \{0\}\)}{\(f'(x) = 0;\ x\in \mathbb{R}\setminus \{0\}\)}{\(f'(x) =\pi ;\ x\in \mathbb{R}\setminus \{0\}\)}{\(f'(x) = \frac{\pi } {x^{2}} ;\ x\in \mathbb{R}\setminus \{0\}\)}{}{}}
\End

\Data\ID{9000070810}
\Updated{2020-07-15 03:07:37}
\Level{A}
\SubArea{Derivative}
\LANGen{
\Question{
Differentiate the following function. 
\[ 
 f(x)=\log _{5}12
\]}
{\(f'(x) = 0;\ x\in \mathbb{R}\)}{\(f'(x) = \frac{1} 
{\ln 12};\ x\in \mathbb{R}\)}{\(f'(x) = \frac{1} 
{12\ln 5};\ x\in \mathbb{R}\)}{\(f'(x) = 1;\ x\in \mathbb{R}\)}{}{}}
\LANGpl{
\Question{
Wyznacz pochodną funkcji.
\[ 
 f\colon y =\log _{5}12
\]}
{\(f'(x) = 0;\ x\in \mathbb{R}\)}{\(f'(x) = \frac{1} 
{\ln 12};\ x\in \mathbb{R}\)}{\(f'(x) = \frac{1} 
{12\ln 5};\ x\in \mathbb{R}\)}{\(f'(x) = 1;\ x\in \mathbb{R}\)}{}{}}
\LANGcs{
\Question{
Určete první derivaci funkce \(f\colon y =\log _{5}12\).}
{\(f'(x) = 0;\ x\in \mathbb{R}\)}{\(f'(x) = \frac{1} 
{\ln 12};\ x\in \mathbb{R}\)}{\(f'(x) = \frac{1} 
{12\ln 5};\ x\in \mathbb{R}\)}{\(f'(x) = 1;\ x\in \mathbb{R}\)}{}{}}
\LANGsk{
\Question{
Určte prvú deriváciu funkcie \(f\colon y =\log _{5}12\).}
{\(f'(x) = 0;\ x\in \mathbb{R}\)}{\(f'(x) = \frac{1} 
{\ln 12};\ x\in \mathbb{R}\)}{\(f'(x) = \frac{1} 
{12\ln 5};\ x\in \mathbb{R}\)}{\(f'(x) = 1;\ x\in \mathbb{R}\)}{}{}}
\LANGes{
\Question{
Deriva la siguiente función.
\[ 
 f(x)=\log _{5}12
\]}
{\(f'(x) = 0;\ x\in \mathbb{R}\)}{\(f'(x) = \frac{1} 
{\ln 12};\ x\in \mathbb{R}\)}{\(f'(x) = \frac{1} 
{12\ln 5};\ x\in \mathbb{R}\)}{\(f'(x) = 1;\ x\in \mathbb{R}\)}{}{}}
\End

\Data\ID{1003112008}
\Updated{2020-07-15 03:07:12}
\Level{B}
\SubArea{Derivative}
\LANGen{
\Question{
Which of the statements A, B, C, D given bellow are correct?
\[
\begin{array}{l}
\text{A: }\ \left(3x^{-3}-\frac5{x^2} +7\right)'=-\frac9{x^4}-\frac{10}{x^3}\text{, }x\in\mathbb{R}\setminus\{0\} \\
\text{B: }\ \left(\frac{x^3-4}{3x}\right)'=2x+\frac4{3x^2}\text{, }x\in\mathbb{R}\setminus\{0\} \\
\text{C: }\ \left(\frac{x^4-x+1}{x}\right)'=3x^2-\frac{1}{x^2}\text{, }x\in\mathbb{R}\setminus\{0\} \\
\text{D: }\ \left(\frac2{x^2}-\frac3{x^3} \right)'=-\frac4{x^3}+\frac9{x^4}\text{, }x\in\mathbb{R}\setminus\{0\}
\end{array}
\]
The only correct statements are:}
{C, D}{B, C}{A, B}{A, C, D}{B, C, D}{A, D}}
\LANGpl{
\Question{
Które z poniższych twierdzeń  A, B, C, D są prawidłowe?
\[
\begin{array}{l}
\text{A: }\ \left(3x^{-3}-\frac5{x^2} +7\right)'=-\frac9{x^4}-\frac{10}{x^3}\text{, }x\in\mathbb{R}\setminus\{0\} \\
\text{B: }\ \left(\frac{x^3-4}{3x}\right)'=2x+\frac4{3x^2}\text{, }x\in\mathbb{R}\setminus\{0\} \\
\text{C: }\ \left(\frac{x^4-x+1}{x}\right)'=3x^2-\frac{1}{x^2}\text{, }x\in\mathbb{R}\setminus\{0\} \\
\text{D: }\ \left(\frac2{x^2}-\frac3{x^3} \right)'=-\frac4{x^3}+\frac9{x^4}\text{, }x\in\mathbb{R}\setminus\{0\}
\end{array}
\]
Jedyne poprawne twierdzenia to:}
{C, D}{B, C}{A, B}{A, C, D}{B, C, D}{A, D}}
\LANGcs{
\Question{
Která z tvrzení A, B, C, D uvedených níže jsou správná?
\[
\begin{array}{l}
\text{A: }\ \left(3x^{-3}-\frac5{x^2} +7\right)'=-\frac9{x^4}-\frac{10}{x^3}\text{, }x\in\mathbb{R}\setminus\{0\} \\
\text{B: }\ \left(\frac{x^3-4}{3x}\right)'=2x+\frac4{3x^2}\text{, }x\in\mathbb{R}\setminus\{0\} \\
\text{C: }\ \left(\frac{x^4-x+1}{x}\right)'=3x^2-\frac{1}{x^2}\text{, }x\in\mathbb{R}\setminus\{0\} \\
\text{D: }\ \left(\frac2{x^2}-\frac3{x^3} \right)'=-\frac4{x^3}+\frac9{x^4}\text{, }x\in\mathbb{R}\setminus\{0\}
\end{array}
\]
Jedinými správnými tvrzeními jsou:}
{C, D}{B, C}{A, B}{A, C, D}{B, C, D}{A, D}}
\LANGsk{
\Question{
Ktoré z tvrdení A, B, C, D uvedených nižšie sú správne?
\[
\begin{array}{l}
\text{A: }\ \left(3x^{-3}-\frac5{x^2} +7\right)'=-\frac9{x^4}-\frac{10}{x^3}\text{, }x\in\mathbb{R}\setminus\{0\} \\
\text{B: }\ \left(\frac{x^3-4}{3x}\right)'=2x+\frac4{3x^2}\text{, }x\in\mathbb{R}\setminus\{0\} \\
\text{C: }\ \left(\frac{x^4-x+1}{x}\right)'=3x^2-\frac{1}{x^2}\text{, }x\in\mathbb{R}\setminus\{0\} \\
\text{D: }\ \left(\frac2{x^2}-\frac3{x^3} \right)'=-\frac4{x^3}+\frac9{x^4}\text{, }x\in\mathbb{R}\setminus\{0\}
\end{array}
\]
Jedinými správnymi tvrdeniami sú:}
{C, D}{B, C}{A, B}{A, C, D}{B, C, D}{A, D}}
\LANGes{
\Question{
De los enunciados A, B, C, D ¿Cuáles son correctos? 
\[
\begin{array}{l}
\text{A: }\ \left(3x^{-3}-\frac5{x^2} +7\right)'=-\frac9{x^4}-\frac{10}{x^3}\text{, }x\in\mathbb{R}\setminus\{0\} \\
\text{B: }\ \left(\frac{x^3-4}{3x}\right)'=2x+\frac4{3x^2}\text{, }x\in\mathbb{R}\setminus\{0\} \\
\text{C: }\ \left(\frac{x^4-x+1}{x}\right)'=3x^2-\frac{1}{x^2}\text{, }x\in\mathbb{R}\setminus\{0\} \\
\text{D: }\ \left(\frac2{x^2}-\frac3{x^3} \right)'=-\frac4{x^3}+\frac9{x^4}\text{, }x\in\mathbb{R}\setminus\{0\}
\end{array}
\]
Los enunciados correctos son solamente:}
{C, D}{B, C}{A, B}{A, C, D}{B, C, D}{A, D}}
\End

\Data\ID{1003112010}
\Updated{2020-07-15 03:07:12}
\Level{B}
\SubArea{Derivative}
\LANGen{
\Question{
Differentiate the following function:
\[ f(x)=\frac{\sqrt{x\sqrt x}}x \]}
{\( f'(x)=-\frac{\sqrt[4]{x^3}}{4x^2}\text{, }x\in\mathbb{R}^+  \)}{\( f'(x)=-\frac{\sqrt[4]{x^3}}{x^2}\text{, }x\in\mathbb{R}^+	 \)}{\( f'(x)=-\frac{\sqrt[4]x}{4x}\text{, }x\in\mathbb{R}^+	 \)}{\( f'(x)=\frac{\sqrt[4]x}{4x}\text{, }x\in\mathbb{R}^+ \)}{}{}}
\LANGpl{
\Question{
Wyznacz pochodną funkcji.
\[ f(x)=\frac{\sqrt{x\sqrt x}}x \]}
{\( f'(x)=-\frac{\sqrt[4]{x^3}}{4x^2}\text{, }x\in\mathbb{R}^+  \)}{\( f'(x)=-\frac{\sqrt[4]{x^3}}{x^2}\text{, }x\in\mathbb{R}^+	 \)}{\( f'(x)=-\frac{\sqrt[4]x}{4x}\text{, }x\in\mathbb{R}^+	 \)}{\( f'(x)=\frac{\sqrt[4]x}{4x}\text{, }x\in\mathbb{R}^+ \)}{}{}}
\LANGcs{
\Question{
Derivujte následující funkci:
\[ f(x)=\frac{\sqrt{x\sqrt x}}x \]}
{\( f'(x)=-\frac{\sqrt[4]{x^3}}{4x^2}\text{, }x\in\mathbb{R}^+  \)}{\( f'(x)=-\frac{\sqrt[4]{x^3}}{x^2}\text{, }x\in\mathbb{R}^+	 \)}{\( f'(x)=-\frac{\sqrt[4]x}{4x}\text{, }x\in\mathbb{R}^+	 \)}{\( f'(x)=\frac{\sqrt[4]x}{4x}\text{, }x\in\mathbb{R}^+ \)}{}{}}
\LANGsk{
\Question{
Derivujte následujúcu funkciu:
\[ f(x)=\frac{\sqrt{x\sqrt x}}x \]}
{\( f'(x)=-\frac{\sqrt[4]{x^3}}{4x^2}\text{, }x\in\mathbb{R}^+  \)}{\( f'(x)=-\frac{\sqrt[4]{x^3}}{x^2}\text{, }x\in\mathbb{R}^+	 \)}{\( f'(x)=-\frac{\sqrt[4]x}{4x}\text{, }x\in\mathbb{R}^+	 \)}{\( f'(x)=\frac{\sqrt[4]x}{4x}\text{, }x\in\mathbb{R}^+ \)}{}{}}
\LANGes{
\Question{
Deriva la siguiente función: 
\[ f(x)=\frac{\sqrt{x\sqrt x}}x \]}
{\( f'(x)=-\frac{\sqrt[4]{x^3}}{4x^2}\text{, }x\in\mathbb{R}^+  \)}{\( f'(x)=-\frac{\sqrt[4]{x^3}}{x^2}\text{, }x\in\mathbb{R}^+	 \)}{\( f'(x)=-\frac{\sqrt[4]x}{4x}\text{, }x\in\mathbb{R}^+	 \)}{\( f'(x)=\frac{\sqrt[4]x}{4x}\text{, }x\in\mathbb{R}^+ \)}{}{}}
\End

\Data\ID{1003112011}
\Updated{2020-07-15 03:07:12}
\Level{B}
\SubArea{Derivative}
\LANGen{
\Question{
Differentiate the following function:
\[ f(x)=\frac{x-5\sqrt[3]{x^2}}{\sqrt[3]x} \]}
{\( f'(x)=\frac{2\sqrt[3]{x^2}-5\sqrt[3] x}{3x}\text{, }x\in\mathbb{R}\setminus\{0\} \)}{\( f'(x)=\frac{2\sqrt[3]{x^2}-5\sqrt[3]x }{3}\text{, }x\in\mathbb{R}\setminus\{0\} \)}{\(  f'(x)=\frac{2\sqrt[3]x-5\sqrt[3]{x^2}}{3x}\text{, }x\in\mathbb{R}\setminus\{0\} \)}{\(   f'(x)=\frac{\sqrt[3]x}{3x}-5\text{, }x\in\mathbb{R}\setminus\{0\} \)}{}{}}
\LANGpl{
\Question{
Wyznacz pochodną funkcji.
\[ f(x)=\frac{x-5\sqrt[3]{x^2}}{\sqrt[3]x} \]}
{\( f'(x)=\frac{2\sqrt[3]{x^2}-5\sqrt[3] x}{3x}\text{, }x\in\mathbb{R}\setminus\{0\} \)}{\( f'(x)=\frac{2\sqrt[3]{x^2}-5\sqrt[3]x }{3}\text{, }x\in\mathbb{R}\setminus\{0\} \)}{\(  f'(x)=\frac{2\sqrt[3]x-5\sqrt[3]{x^2}}{3x}\text{, }x\in\mathbb{R}\setminus\{0\} \)}{\(   f'(x)=\frac{\sqrt[3]x}{3x}-5\text{, }x\in\mathbb{R}\setminus\{0\} \)}{}{}}
\LANGcs{
\Question{
Derivujte následující funkci:
\[ f(x)=\frac{x-5\sqrt[3]{x^2}}{\sqrt[3]x} \]}
{\( f'(x)=\frac{2\sqrt[3]{x^2}-5\sqrt[3] x}{3x}\text{, }x\in\mathbb{R}\setminus\{0\} \)}{\( f'(x)=\frac{2\sqrt[3]{x^2}-5\sqrt[3]x }{3}\text{, }x\in\mathbb{R}\setminus\{0\} \)}{\(  f'(x)=\frac{2\sqrt[3]x-5\sqrt[3]{x^2}}{3x}\text{, }x\in\mathbb{R}\setminus\{0\} \)}{\(   f'(x)=\frac{\sqrt[3]x}{3x}-5\text{, }x\in\mathbb{R}\setminus\{0\} \)}{}{}}
\LANGsk{
\Question{
Derivujte následujúcu funkciu:
\[ f(x)=\frac{x-5\sqrt[3]{x^2}}{\sqrt[3]x} \]}
{\( f'(x)=\frac{2\sqrt[3]{x^2}-5\sqrt[3] x}{3x}\text{, }x\in\mathbb{R}\setminus\{0\} \)}{\( f'(x)=\frac{2\sqrt[3]{x^2}-5\sqrt[3]x }{3}\text{, }x\in\mathbb{R}\setminus\{0\} \)}{\(  f'(x)=\frac{2\sqrt[3]x-5\sqrt[3]{x^2}}{3x}\text{, }x\in\mathbb{R}\setminus\{0\} \)}{\(   f'(x)=\frac{\sqrt[3]x}{3x}-5\text{, }x\in\mathbb{R}\setminus\{0\} \)}{}{}}
\LANGes{
\Question{
Deriva la siguiente función: 
\[ f(x)=\frac{x-5\sqrt[3]{x^2}}{\sqrt[3]x} \]}
{\( f'(x)=\frac{2\sqrt[3]{x^2}-5\sqrt[3] x}{3x}\text{, }x\in\mathbb{R}\setminus\{0\} \)}{\( f'(x)=\frac{2\sqrt[3]{x^2}-5\sqrt[3]x }{3}\text{, }x\in\mathbb{R}\setminus\{0\} \)}{\(  f'(x)=\frac{2\sqrt[3]x-5\sqrt[3]{x^2}}{3x}\text{, }x\in\mathbb{R}\setminus\{0\} \)}{\(   f'(x)=\frac{\sqrt[3]x}{3x}-5\text{, }x\in\mathbb{R}\setminus\{0\} \)}{}{}}
\End

\Data\ID{1003230201}
\Updated{2020-07-15 03:07:21}
\Level{B}
\SubArea{Derivative}
\LANGen{
\Question{
Given the function \( f(x)=x\cdot\,\mathrm{e}^x \), find the set of all \( x \), \( x\in\mathbb{R} \), for which \( f'(x)=0 \).}
{\( \{-1\} \)}{\( \{-1;0\} \)}{\( \{0;1\} \)}{\( \emptyset \)}{\( \{0\} \)}{\( \{1\} \)}}
\LANGpl{
\Question{
Dana je funkcja \( f(x)=x\cdot\,\mathrm{e}^x \), wskaż zbiór wszystkich \( x \), \( x\in\mathbb{R} \), dla których \( f'(x)=0 \).}
{\( \{-1\} \)}{\( \{-1;0\} \)}{\( \{0;1\} \)}{\( \emptyset \)}{\( \{0\} \)}{\( \{1\} \)}}
\LANGcs{
\Question{
Je dána funkce \( f(x)=x\cdot\,\mathrm{e}^x \). Najděte množinu všech \( x \), \( x\in\mathbb{R} \), pro která platí \( f'(x)=0 \).}
{\( \{-1\} \)}{\( \{-1;0\} \)}{\( \{0;1\} \)}{\( \emptyset \)}{\( \{0\} \)}{\( \{1\} \)}}
\LANGsk{
\Question{
Je daná funkcia \( f(x)=x\cdot\,\mathrm{e}^x \). Nájdite množinu všetkých \( x \), \( x\in\mathbb{R} \), pre ktoré platí \( f'(x)=0 \).}
{\( \{-1\} \)}{\( \{-1;0\} \)}{\( \{0;1\} \)}{\( \emptyset \)}{\( \{0\} \)}{\( \{1\} \)}}
\LANGes{
\Question{
Dada la función\( f(x)=x\cdot\,\mathrm{e}^x \), halla el conjunto de todos los \( x \), \( x\in\mathbb{R} \), para los cuales \( f'(x)=0 \).}
{\( \{-1\} \)}{\( \{-1;0\} \)}{\( \{0;1\} \)}{\( \emptyset \)}{\( \{0\} \)}{\( \{1\} \)}}
\End

\Data\ID{1003230202}
\Updated{2020-07-15 03:07:21}
\Level{B}
\SubArea{Derivative}
\LANGen{
\Question{
Given the function \( f(x)=\frac x{4-x^2} \), find the set of all \( x \), \( x\in\mathbb{R} \), for which \( f'(x)=0 \).}
{\( \emptyset \)}{\( \{ 2 \} \)}{\( \{ -2;2 \} \)}{\( \{ -2 \} \)}{\( \{ 0 \} \)}{\( \{ 4 \} \)}}
\LANGpl{
\Question{
Dana jest funkcja \( f(x)=\frac x{4-x^2} \), wskaż zbiór wszystkich \( x \), \( x\in\mathbb{R} \), dla których \( f'(x)=0 \).}
{\( \emptyset \)}{\( \{ 2 \} \)}{\( \{ -2;2 \} \)}{\( \{ -2 \} \)}{\( \{ 0 \} \)}{\( \{ 4 \} \)}}
\LANGcs{
\Question{
Je dána funkce \( f(x)=\frac x{4-x^2} \). Najděte množinu všech \( x \), \( x\in\mathbb{R} \), pro která platí \( f'(x)=0 \).}
{\( \emptyset \)}{\( \{ 2 \} \)}{\( \{ -2;2 \} \)}{\( \{ -2 \} \)}{\( \{ 0 \} \)}{\( \{ 4 \} \)}}
\LANGsk{
\Question{
Je daná funkcia \( f(x)=\frac x{4-x^2} \). Nájdite množinu všetkých \( x \), \( x\in\mathbb{R} \), pre ktoré platí \( f'(x)=0 \).}
{\( \emptyset \)}{\( \{ 2 \} \)}{\( \{ -2;2 \} \)}{\( \{ -2 \} \)}{\( \{ 0 \} \)}{\( \{ 4 \} \)}}
\LANGes{
\Question{
Dada la función \( f(x)=\frac x{4-x^2} \), halla el conjunto de todos los \( x \), \( x\in\mathbb{R} \), para los cuales \( f'(x)=0 \).}
{\( \emptyset \)}{\( \{ 2 \} \)}{\( \{ -2;2 \} \)}{\( \{ -2 \} \)}{\( \{ 0 \} \)}{\( \{ 4 \} \)}}
\End

\Data\ID{1003230203}
\Updated{2020-07-15 03:07:21}
\Level{B}
\SubArea{Derivative}
\LANGen{
\Question{
Given the function \( f(x)=\frac{\sqrt x}{\ln x} \), find the set of all \( x \), \( x\in\mathbb{R} \), for which \( f'(x)=0 \).}
{\( \left\{ \mathrm{e}^2 \right\} \)}{\( \{ \mathrm{e} \} \)}{\( \left\{ \sqrt{\mathrm{e}} \right\} \)}{\( \left\{ \frac1{\mathrm{e}};\mathrm{e} \right\} \)}{\( \{ 2 \} \)}{\( \left\{ 1;\mathrm{e}^2 \right\} \)}}
\LANGpl{
\Question{
Dana jest funkcja \( f(x)=\frac{\sqrt x}{\ln x} \), wskaż zbiór wszystkich \( x \), \( x\in\mathbb{R} \), dla których \( f'(x)=0 \).}
{\( \left\{ \mathrm{e}^2 \right\} \)}{\( \{ \mathrm{e} \} \)}{\( \left\{ \sqrt{\mathrm{e}} \right\} \)}{\( \left\{ \frac1{\mathrm{e}};\mathrm{e} \right\} \)}{\( \{ 2 \} \)}{\( \left\{ 1;\mathrm{e}^2 \right\} \)}}
\LANGcs{
\Question{
Je dána funkce \( f(x)=\frac{\sqrt x}{\ln x} \). Najděte množinu všech \( x \), \( x\in\mathbb{R} \), pro která platí \( f'(x)=0 \).}
{\( \left\{ \mathrm{e}^2 \right\} \)}{\( \{ \mathrm{e} \} \)}{\( \left\{ \sqrt{\mathrm{e}} \right\} \)}{\( \left\{ \frac1{\mathrm{e}};\mathrm{e} \right\} \)}{\( \{ 2 \} \)}{\( \left\{ 1;\mathrm{e}^2 \right\} \)}}
\LANGsk{
\Question{
Je daná funkcia \( f(x)=\frac{\sqrt x}{\ln x} \). Nájdite množinu všetkých \( x \), \( x\in\mathbb{R} \), pre ktoré platí \( f'(x)=0 \).}
{\( \left\{ \mathrm{e}^2 \right\} \)}{\( \{ \mathrm{e} \} \)}{\( \left\{ \sqrt{\mathrm{e}} \right\} \)}{\( \left\{ \frac1{\mathrm{e}};\mathrm{e} \right\} \)}{\( \{ 2 \} \)}{\( \left\{ 1;\mathrm{e}^2 \right\} \)}}
\LANGes{
\Question{
Dada la función \( f(x)=\frac{\sqrt x}{\ln x} \), halla el conjunto de todos los \( x \), \( x\in\mathbb{R} \), para los cuales \( f'(x)=0 \).}
{\( \left\{ \mathrm{e}^2 \right\} \)}{\( \{ \mathrm{e} \} \)}{\( \left\{ \sqrt{\mathrm{e}} \right\} \)}{\( \left\{ \frac1{\mathrm{e}};\mathrm{e} \right\} \)}{\( \{ 2 \} \)}{\( \left\{ 1;\mathrm{e}^2 \right\} \)}}
\End

\Data\ID{1003230204}
\Updated{2020-07-15 03:07:21}
\Level{B}
\SubArea{Derivative}
\LANGen{
\Question{
Which of the statements A, B, C given bellow are correct?
\[
\begin{array}{l}
\text{A: }\left(\frac1{x^3}\cdot\cos x\right)' =-\frac{\cos x+\sin x}{x^4},\ x\in\mathbb{R}\setminus\{0\} \\
\text{B: }\bigl(\left(1-x^3\right)\cdot\ln x \bigr)'=-3x^2\ln x+\frac1x - x^2,\ x\in\mathbb{R}^+ \\
\text{C: } \left(5^x\cdot\sqrt[5]x\right)'=5^{x-1}\sqrt[5]x\left(5\ln5+\frac1x\right),\ x\in\mathbb{R}\setminus\{0\} 
\end{array}
\]
The only correct statements are:}
{B, C}{A, C}{A, B}{B}{A}{C}}
\LANGpl{
\Question{
Które z poniższych wyrażeń  A, B, C jest poprawne?
\[
\begin{array}{l}
\text{A: }\left(\frac1{x^3}\cdot\cos x\right)' =-\frac{\cos x+\sin x}{x^4},\ x\in\mathbb{R}\setminus\{0\} \\
\text{B: }\bigl(\left(1-x^3\right)\cdot\ln x \bigr)'=-3x^2\ln x+\frac1x - x^2,\ x\in\mathbb{R}^+ \\
\text{C: } \left(5^x\cdot\sqrt[5]x\right)'=5^{x-1}\sqrt[5]x\left(5\ln5+\frac1x\right),\ x\in\mathbb{R}\setminus\{0\} 
\end{array}
\]
Jedyne poprawne wyrażenia to:}
{B, C}{A, C}{A, B}{B}{A}{C}}
\LANGcs{
\Question{
Která z tvrzení A, B, C uvedených níže jsou správná?
\[
\begin{array}{l}
\text{A: }\left(\frac1{x^3}\cdot\cos x\right)' =-\frac{\cos x+\sin x}{x^4},\ x\in\mathbb{R}\setminus\{0\} \\
\text{B: }\bigl(\left(1-x^3\right)\cdot\ln x \bigr)'=-3x^2\ln x+\frac1x - x^2,\ x\in\mathbb{R}^+ \\
\text{C: } \left(5^x\cdot\sqrt[5]x\right)'=5^{x-1}\sqrt[5]x\left(5\ln5+\frac1x\right),\ x\in\mathbb{R}\setminus\{0\} 
\end{array}
\]
Jedinými správnými tvrzeními jsou:}
{B, C}{A, C}{A, B}{B}{A}{C}}
\LANGsk{
\Question{
Ktoré z tvrdení A, B, C uvedených nižšie sú správne?
\[
\begin{array}{l}
\text{A: }\left(\frac1{x^3}\cdot\cos x\right)' =-\frac{\cos x+\sin x}{x^4},\ x\in\mathbb{R}\setminus\{0\} \\
\text{B: }\bigl(\left(1-x^3\right)\cdot\ln x \bigr)'=-3x^2\ln x+\frac1x - x^2,\ x\in\mathbb{R}^+ \\
\text{C: } \left(5^x\cdot\sqrt[5]x\right)'=5^{x-1}\sqrt[5]x\left(5\ln5+\frac1x\right),\ x\in\mathbb{R}\setminus\{0\} 
\end{array}
\]
Jedinými správnymi tvrdeniami sú:}
{B, C}{A, C}{A, B}{B}{A}{C}}
\LANGes{
\Question{
De los enunciados A, B, C, D ¿Cuáles  son correctos?
\[
\begin{array}{l}
\text{A: }\left(\frac1{x^3}\cdot\cos x\right)' =-\frac{\cos x+\sin x}{x^4},\ x\in\mathbb{R}\setminus\{0\} \\
\text{B: }\bigl(\left(1-x^3\right)\cdot\ln x \bigr)'=-3x^2\ln x+\frac1x - x^2,\ x\in\mathbb{R}^+ \\
\text{C: } \left(5^x\cdot\sqrt[5]x\right)'=5^{x-1}\sqrt[5]x\left(5\ln5+\frac1x\right),\ x\in\mathbb{R}\setminus\{0\} 
\end{array}
\]
Los enunciados correctos son solamente:}
{B, C}{A, C}{A, B}{B}{A}{C}}
\End

\Data\ID{1003230205}
\Updated{2020-07-15 03:07:21}
\Level{B}
\SubArea{Derivative}
\LANGen{
\Question{
Which of the statements A, B, C given bellow are correct?
\[
\begin{array}{l}
\text{A: } \left(\frac{2x-1}{2-x}\right)'=\frac{5-4x}{(2-x)^2},\ x\neq2 \\
\text{B: } \left(\frac{\mathrm{e}^x-1}{x}\right)'=\frac{\mathrm{e}^x(x-1)-1}{x^2},\ x\neq0 \\
\text{C: } \left(\frac{\cos x}{1-\sin x}\right)'=\frac1{1-\sin x},\ x\neq\frac{\pi}2+2k\pi,\ k\in\mathbb{Z}
\end{array}\]
The only correct statements are:}
{C}{A, C}{A, B}{B}{A}{B, C}}
\LANGpl{
\Question{
Które z wyrażeń  A, B, C podanych poniżej są poprawne?
\[
\begin{array}{l}
\text{A: } \left(\frac{2x-1}{2-x}\right)'=\frac{5-4x}{(2-x)^2},\ x\neq2 \\
\text{B: } \left(\frac{\mathrm{e}^x-1}{x}\right)'=\frac{\mathrm{e}^x(x-1)-1}{x^2},\ x\neq0 \\
\text{C: } \left(\frac{\cos x}{1-\sin x}\right)'=\frac1{1-\sin x},\ x\neq\frac{\pi}2+2k\pi,\ k\in\mathbb{Z}
\end{array}\]
Jedyne poprawne wyrażenia to:}
{C}{A, C}{A, B}{B}{A}{B, C}}
\LANGcs{
\Question{
Která z tvrzení A, B, C uvedených níže jsou správná?
\[
\begin{array}{l}
\text{A: } \left(\frac{2x-1}{2-x}\right)'=\frac{5-4x}{(2-x)^2},\ x\neq2 \\
\text{B: } \left(\frac{\mathrm{e}^x-1}{x}\right)'=\frac{\mathrm{e}^x(x-1)-1}{x^2},\ x\neq0 \\
\text{C: } \left(\frac{\cos x}{1-\sin x}\right)'=\frac1{1-\sin x},\ x\neq\frac{\pi}2+2k\pi,\ k\in\mathbb{Z}
\end{array}\]
Jedinými správnými tvrzeními jsou:}
{C}{A, C}{A, B}{B}{A}{B, C}}
\LANGsk{
\Question{
Ktoré z tvrdení A, B, C uvedených nižšie sú správne?
\[
\begin{array}{l}
\text{A: } \left(\frac{2x-1}{2-x}\right)'=\frac{5-4x}{(2-x)^2},\ x\neq2 \\
\text{B: } \left(\frac{\mathrm{e}^x-1}{x}\right)'=\frac{\mathrm{e}^x(x-1)-1}{x^2},\ x\neq0 \\
\text{C: } \left(\frac{\cos x}{1-\sin x}\right)'=\frac1{1-\sin x},\ x\neq\frac{\pi}2+2k\pi,\ k\in\mathbb{Z}
\end{array}\]
Jedinými správnymi tvrdeniami sú:}
{C}{A, C}{A, B}{B}{A}{B, C}}
\LANGes{
\Question{
De los enunciados A, B, C, D ¿Cuáles  son correctos?
\[
\begin{array}{l}
\text{A: } \left(\frac{2x-1}{2-x}\right)'=\frac{5-4x}{(2-x)^2},\ x\neq2 \\
\text{B: } \left(\frac{\mathrm{e}^x-1}{x}\right)'=\frac{\mathrm{e}^x(x-1)-1}{x^2},\ x\neq0 \\
\text{C: } \left(\frac{\cos x}{1-\sin x}\right)'=\frac1{1-\sin x},\ x\neq\frac{\pi}2+2k\pi,\ k\in\mathbb{Z}
\end{array}\]
Los enunciados correctos son solamente:}
{C}{A, C}{A, B}{B}{A}{B, C}}
\End

\Data\ID{1003230206}
\Updated{2020-07-15 03:07:21}
\Level{B}
\SubArea{Derivative}
\LANGen{
\Question{
Which of the statements A, B, C, D given bellow are incorrect?
\[
\begin{array}{l}
\text{A: }\left(\ln\frac x2\right)'=\frac1x,\ x\in\mathbb{R}^+ \\
\text{B: }\left(5\sin3x\right)'=5\cos3x \\
\text{C: }\left(\frac1{\left(x^3-1\right)^2}\right)'=\frac{-6x^2}{\left(x^3-1\right)^3},\ x\in\mathbb{R}\setminus\{1\} \\
\text{D: }\left(\ln(1+\cos x ) \right)'=\frac1{1-\sin x},\ x\neq\frac{\pi}2+2k\pi,\ k\in\mathbb{Z}
\end{array}\]
The only incorrect statements are:}
{B, D}{A, B, D}{B, C}{B}{A, C}{A, C, D}}
\LANGpl{
\Question{
Które z podanych wyrażeń A, B, C, D są niepoprawne?
\[
\begin{array}{l}
\text{A: }\left(\ln\frac x2\right)'=\frac1x,\ x\in\mathbb{R}^+ \\
\text{B: }\left(5\sin3x\right)'=5\cos3x \\
\text{C: }\left(\frac1{\left(x^3-1\right)^2}\right)'=\frac{-6x^2}{\left(x^3-1\right)^3},\ x\in\mathbb{R}\setminus\{1\} \\
\text{D: }\left(\ln(1+\cos x ) \right)'=\frac1{1-\sin x},\ x\neq\frac{\pi}2+2k\pi,\ k\in\mathbb{Z}
\end{array}\]
Jedyne niepoprawne wyrażenia to:}
{B, D}{A, B, D}{B, C}{B}{A, C}{A, C, D}}
\LANGcs{
\Question{
Která z tvrzení A, B, C, D uvedených níže nejsou správná?
\[
\begin{array}{l}
\text{A: }\left(\ln\frac x2\right)'=\frac1x,\ x\in\mathbb{R}^+ \\
\text{B: }\left(5\sin3x\right)'=5\cos3x \\
\text{C: }\left(\frac1{\left(x^3-1\right)^2}\right)'=\frac{-6x^2}{\left(x^3-1\right)^3},\ x\in\mathbb{R}\setminus\{1\} \\
\text{D: }\left(\ln(1+\cos x ) \right)'=\frac1{1-\sin x},\ x\neq\frac{\pi}2+2k\pi,\ k\in\mathbb{Z}
\end{array}
\]
Jedinými nesprávnými tvrzeními jsou:}
{B, D}{A, B, D}{B, C}{B}{A, C}{A, C, D}}
\LANGsk{
\Question{
Ktoré z tvrdení A, B, C, D uvedených nižšie nie sú správne?
\[
\begin{array}{l}
\text{A: }\left(\ln\frac x2\right)'=\frac1x,\ x\in\mathbb{R}^+ \\
\text{B: }\left(5\sin3x\right)'=5\cos3x \\
\text{C: }\left(\frac1{\left(x^3-1\right)^2}\right)'=\frac{-6x^2}{\left(x^3-1\right)^3},\ x\in\mathbb{R}\setminus\{1\} \\
\text{D: }\left(\ln(1+\cos x ) \right)'=\frac1{1-\sin x},\ x\neq\frac{\pi}2+2k\pi,\ k\in\mathbb{Z}
\end{array}
\]
Jedinými nesprávnymi tvrdeniami sú:}
{B, D}{A, B, D}{B, C}{B}{A, C}{A, C, D}}
\LANGes{
\Question{
De los enunciados A, B, C, D ¿Cuáles  son incorrectos?
\[
\begin{array}{l}
\text{A: }\left(\ln\frac x2\right)'=\frac1x,\ x\in\mathbb{R}^+ \\
\text{B: }\left(5\sin3x\right)'=5\cos3x \\
\text{C: }\left(\frac1{\left(x^3-1\right)^2}\right)'=\frac{-6x^2}{\left(x^3-1\right)^3},\ x\in\mathbb{R}\setminus\{1\} \\
\text{D: }\left(\ln(1+\cos x ) \right)'=\frac1{1-\sin x},\ x\neq\frac{\pi}2+2k\pi,\ k\in\mathbb{Z}
\end{array}\]
Los enunciados incorrectos son:}
{B, D}{A, B, D}{B, C}{B}{A, C}{A, C, D}}
\End

\Data\ID{9000063101}
\Updated{2020-07-15 03:07:34}
\Level{B}
\SubArea{Derivative}
\LANGen{
\Question{
Differentiate the following function. 
\[ 
 f(x) = \frac{x^{2} - 1} 
{x^{2} + 1}
\]}
{\(f'(x) = \frac{4x} 
{(x^{2}+1)^{2}} ,\ x\in \mathbb{R}\)}{\(f'(x) = \frac{-4x} 
{x^{2}+1},\ x\in \mathbb{R}\)}{\(f'(x) = \frac{4x^{3}} 
{(x^{2}+1)^{2}} ,\ x\in \mathbb{R}\)}{\(f'(x) = \frac{4x} 
{x^{2}+1},\ x\in \mathbb{R}\)}{}{}}
\LANGpl{
\Question{
Wyznacz pochodną funkcji. 
\[ 
 f\colon y = \frac{x^{2} - 1} 
{x^{2} + 1}
\]}
{\(f'(x) = \frac{4x} 
{(x^{2}+1)^{2}} ,\ x\in \mathbb{R}\)}{\(f'(x) = \frac{-4x} 
{x^{2}+1},\ x\in \mathbb{R}\)}{\(f'(x) = \frac{4x^{3}} 
{(x^{2}+1)^{2}} ,\ x\in \mathbb{R}\)}{\(f'(x) = \frac{4x} 
{x^{2}+1},\ x\in \mathbb{R}\)}{}{}}
\LANGcs{
\Question{
Derivace funkce \(f\colon y = \frac{x^{2}-1} 
{x^{2}+1}\) 
je rovna:}
{\(f'(x) = \frac{4x} 
{(x^{2}+1)^{2}} ,\ x\in \mathbb{R}\)}{\(f'(x) = \frac{-4x} 
{x^{2}+1},\ x\in \mathbb{R}\)}{\(f'(x) = \frac{4x^{3}} 
{(x^{2}+1)^{2}} ,\ x\in \mathbb{R}\)}{\(f'(x) = \frac{4x} 
{x^{2}+1},\ x\in \mathbb{R}\)}{}{}}
\LANGsk{
\Question{
Derivácia funkcie \(f\colon y = \frac{x^{2}-1} 
{x^{2}+1}\) 
je rovná:}
{\(f'(x) = \frac{4x} 
{(x^{2}+1)^{2}} ,\ x\in \mathbb{R}\)}{\(f'(x) = \frac{-4x} 
{x^{2}+1},\ x\in \mathbb{R}\)}{\(f'(x) = \frac{4x^{3}} 
{(x^{2}+1)^{2}} ,\ x\in \mathbb{R}\)}{\(f'(x) = \frac{4x} 
{x^{2}+1},\ x\in \mathbb{R}\)}{}{}}
\LANGes{
\Question{
Deriva la siguiente función.
\[ 
 f(x) = \frac{x^{2} - 1} 
{x^{2} + 1}
\]}
{\(f'(x) = \frac{4x} 
{(x^{2}+1)^{2}} ,\ x\in \mathbb{R}\)}{\(f'(x) = \frac{-4x} 
{x^{2}+1},\ x\in \mathbb{R}\)}{\(f'(x) = \frac{4x^{3}} 
{(x^{2}+1)^{2}} ,\ x\in \mathbb{R}\)}{\(f'(x) = \frac{4x} 
{x^{2}+1},\ x\in \mathbb{R}\)}{}{}}
\End

\Data\ID{9000063102}
\Updated{2020-07-15 03:07:34}
\Level{B}
\SubArea{Derivative}
\LANGen{
\Question{
Differentiate the following function. 
\[ 
 f(x) = \frac{x^{2} + 1} 
 {2x}
\]}
{\(f'(x) = \frac{x^{2}-1} 
 {2x^{2}} ,\ x\neq 0\)}{\(f'(x) = x,\ x\neq 0\)}{\(f'(x) = \frac{x-1} 
{2x^{2}} ,\ x\neq 0\)}{\(f'(x) = \frac{x^{2}+1} 
 {2x^{2}} ,\ x\neq 0\)}{}{}}
\LANGpl{
\Question{
Wyznacz pochodną funkcji.  
\[ 
 f\colon y = \frac{x^{2} + 1} 
 {2x}
\]}
{\(f'(x) = \frac{x^{2}-1} 
 {2x^{2}} ,\ x\neq 0\)}{\(f'(x) = x,\ x\neq 0\)}{\(f'(x) = \frac{x-1} 
{2x^{2}} ,\ x\neq 0\)}{\(f'(x) = \frac{x^{2}+1} 
 {2x^{2}} ,\ x\neq 0\)}{}{}}
\LANGcs{
\Question{
Derivace funkce \(f\colon y = \frac{x^{2}+1} 
 {2x} \) 
je rovna:}
{\(f'(x) = \frac{x^{2}-1} 
 {2x^{2}} ,\ x\neq 0\)}{\(f'(x) = x,\ x\neq 0\)}{\(f'(x) = \frac{x-1} 
{2x^{2}} ,\ x\neq 0\)}{\(f'(x) = \frac{x^{2}+1} 
 {2x^{2}} ,\ x\neq 0\)}{}{}}
\LANGsk{
\Question{
Derivácia funkcie \(f\colon y = \frac{x^{2}+1} 
 {2x} \) 
je rovná:}
{\(f'(x) = \frac{x^{2}-1} 
 {2x^{2}} ,\ x\neq 0\)}{\(f'(x) = x,\ x\neq 0\)}{\(f'(x) = \frac{x-1} 
{2x^{2}} ,\ x\neq 0\)}{\(f'(x) = \frac{x^{2}+1} 
 {2x^{2}} ,\ x\neq 0\)}{}{}}
\LANGes{
\Question{
Deriva la siguiente función.
\[ 
 f(x) = \frac{x^{2} + 1} 
 {2x}
\]}
{\(f'(x) = \frac{x^{2}-1} 
 {2x^{2}} ,\ x\neq 0\)}{\(f'(x) = x,\ x\neq 0\)}{\(f'(x) = \frac{x-1} 
{2x^{2}} ,\ x\neq 0\)}{\(f'(x) = \frac{x^{2}+1} 
 {2x^{2}} ,\ x\neq 0\)}{}{}}
\End

\Data\ID{9000063103}
\Updated{2020-07-15 03:07:34}
\Level{B}
\SubArea{Derivative}
\LANGen{
\Question{
Differentiate the following function. 
\[ 
 f(x) = \frac{x^{2} - x} 
 {x + 1}
\]}
{\(f'(x) = \frac{x^{2}+2x-1} 
 {(x+1)^{2}} ,\ x\neq - 1\)}{\(f'(x) = 2x - 1,\ x\neq - 1\)}{\(f'(x) = \frac{x^{2}+2x-1} 
 {(x+1)^{2}} ,\ x\neq 0\)}{\(f'(x) = \frac{2x} 
{(x^{2}+1)^{2}} ,\ x\neq 0\)}{}{}}
\LANGpl{
\Question{
Wyznacz pochodną funkcji. 
\[ 
 f\colon y = \frac{x^{2} - x} 
 {x + 1}
\]}
{\(f'(x) = \frac{x^{2}+2x-1} 
 {(x+1)^{2}} ,\ x\neq - 1\)}{\(f'(x) = 2x - 1,\ x\neq - 1\)}{\(f'(x) = \frac{x^{2}+2x-1} 
 {(x+1)^{2}} ,\ x\neq 0\)}{\(f'(x) = \frac{2x} 
{(x^{2}+1)^{2}} ,\ x\neq 0\)}{}{}}
\LANGcs{
\Question{
Derivace funkce \(f\colon y = \frac{x^{2}-x} 
 {x+1} \) 
je rovna:}
{\(f'(x) = \frac{x^{2}+2x-1} 
 {(x+1)^{2}} ,\ x\neq - 1\)}{\(f'(x) = 2x - 1,\ x\neq - 1\)}{\(f'(x) = \frac{x^{2}+2x-1} 
 {(x+1)^{2}} ,\ x\neq 0\)}{\(f'(x) = \frac{2x} 
{(x^{2}+1)^{2}} ,\ x\neq 0\)}{}{}}
\LANGsk{
\Question{
Derivácia funkcie \(f\colon y = \frac{x^{2}-x} 
 {x+1} \) 
je rovná:}
{\(f'(x) = \frac{x^{2}+2x-1} 
 {(x+1)^{2}} ,\ x\neq - 1\)}{\(f'(x) = 2x - 1,\ x\neq - 1\)}{\(f'(x) = \frac{x^{2}+2x-1} 
 {(x+1)^{2}} ,\ x\neq 0\)}{\(f'(x) = \frac{2x} 
{(x^{2}+1)^{2}} ,\ x\neq 0\)}{}{}}
\LANGes{
\Question{
Deriva la siguiente función.
\[ 
 f(x) = \frac{x^{2} - x} 
 {x + 1}
\]}
{\(f'(x) = \frac{x^{2}+2x-1} 
 {(x+1)^{2}} ,\ x\neq - 1\)}{\(f'(x) = 2x - 1,\ x\neq - 1\)}{\(f'(x) = \frac{x^{2}+2x-1} 
 {(x+1)^{2}} ,\ x\neq 0\)}{\(f'(x) = \frac{2x} 
{(x^{2}+1)^{2}} ,\ x\neq 0\)}{}{}}
\End

\Data\ID{9000063104}
\Updated{2020-07-15 03:07:34}
\Level{B}
\SubArea{Derivative}
\LANGen{
\Question{
Differentiate the following function. 
\[ 
 f(x)= \frac{\sin x} 
{\sin x -\cos x}
\]}
{\(f'(x) = \frac{-1} 
{(\sin x-\cos x)^{2}} ,\ x\neq \frac{\pi }{4} + k\pi ;k\in \mathbb{Z}\)}{\(f'(x) = \frac{\sin ^{2}x-\cos ^{2}x} 
{(\sin x-\cos x)^{2}} ,\ x\neq \frac{\pi }{4} + k\pi ;k\in \mathbb{Z}\)}{\(f'(x) = \frac{\sin x(\cos x+1)} 
{(\sin x-\cos x)^{2}} ,\ x\neq \frac{\pi }{4} + k\pi ;k\in \mathbb{Z}\)}{\(f'(x) = \frac{\cos ^{2}x-\sin ^{2}x} 
{(\sin x-\cos x)^{2}} ,\ x\neq \frac{\pi }{4} + k\pi ;k\in \mathbb{Z}\)}{}{}}
\LANGpl{
\Question{
Wyznacz pochodną funkcji. 
\[ 
 f\colon y = \frac{\sin x} 
{\sin x -\cos x}
\]}
{\(f'(x) = \frac{-1} 
{(\sin x-\cos x)^{2}} ,\ x\neq \frac{\pi }{4} + k\pi ;k\in \mathbb{Z}\)}{\(f'(x) = \frac{\sin ^{2}x-\cos ^{2}x} 
{(\sin x-\cos x)^{2}} ,\ x\neq \frac{\pi }{4} + k\pi ;k\in \mathbb{Z}\)}{\(f'(x) = \frac{\sin x(\cos x+1)} 
{(\sin x-\cos x)^{2}} ,\ x\neq \frac{\pi }{4} + k\pi ;k\in \mathbb{Z}\)}{\(f'(x) = \frac{\cos ^{2}x-\sin ^{2}x} 
{(\sin x-\cos x)^{2}} ,\ x\neq \frac{\pi }{4} + k\pi ;k\in \mathbb{Z}\)}{}{}}
\LANGcs{
\Question{
Derivace funkce \(f\colon y = \frac{\sin x} 
{\sin x-\cos x}\) 
je rovna:}
{\(f'(x) = \frac{-1} 
{(\sin x-\cos x)^{2}} ,\ x\neq \frac{\pi }{4} + k\pi ;k\in \mathbb{Z}\)}{\(f'(x) = \frac{\sin ^{2}x-\cos ^{2}x} 
{(\sin x-\cos x)^{2}} ,\ x\neq \frac{\pi }{4} + k\pi ;k\in \mathbb{Z}\)}{\(f'(x) = \frac{\sin x(\cos x+1)} 
{(\sin x-\cos x)^{2}} ,\ x\neq \frac{\pi }{4} + k\pi ;k\in \mathbb{Z}\)}{\(f'(x) = \frac{\cos ^{2}x-\sin ^{2}x} 
{(\sin x-\cos x)^{2}} ,\ x\neq \frac{\pi }{4} + k\pi ;k\in \mathbb{Z}\)}{}{}}
\LANGsk{
\Question{
Derivácia funkcie \(f\colon y = \frac{\sin x} 
{\sin x-\cos x}\) 
je rovná:}
{\(f'(x) = \frac{-1} 
{(\sin x-\cos x)^{2}} ,\ x\neq \frac{\pi }{4} + k\pi ;k\in \mathbb{Z}\)}{\(f'(x) = \frac{\sin ^{2}x-\cos ^{2}x} 
{(\sin x-\cos x)^{2}} ,\ x\neq \frac{\pi }{4} + k\pi ;k\in \mathbb{Z}\)}{\(f'(x) = \frac{\sin x(\cos x+1)} 
{(\sin x-\cos x)^{2}} ,\ x\neq \frac{\pi }{4} + k\pi ;k\in \mathbb{Z}\)}{\(f'(x) = \frac{\cos ^{2}x-\sin ^{2}x} 
{(\sin x-\cos x)^{2}} ,\ x\neq \frac{\pi }{4} + k\pi ;k\in \mathbb{Z}\)}{}{}}
\LANGes{
\Question{
Deriva la siguiente función.
\[ 
 f(x)= \frac{\sin x} 
{\sin x -\cos x}
\]}
{\(f'(x) = \frac{-1} 
{(\sin x-\cos x)^{2}} ,\ x\neq \frac{\pi }{4} + k\pi ;k\in \mathbb{Z}\)}{\(f'(x) = \frac{\sin ^{2}x-\cos ^{2}x} 
{(\sin x-\cos x)^{2}} ,\ x\neq \frac{\pi }{4} + k\pi ;k\in \mathbb{Z}\)}{\(f'(x) = \frac{\sin x(\cos x+1)} 
{(\sin x-\cos x)^{2}} ,\ x\neq \frac{\pi }{4} + k\pi ;k\in \mathbb{Z}\)}{\(f'(x) = \frac{\cos ^{2}x-\sin ^{2}x} 
{(\sin x-\cos x)^{2}} ,\ x\neq \frac{\pi }{4} + k\pi ;k\in \mathbb{Z}\)}{}{}}
\End

\Data\ID{9000063105}
\Updated{2020-07-15 03:07:34}
\Level{B}
\SubArea{Derivative}
\LANGen{
\Question{
Differentiate the following function. 
\[ 
 f(x) = \frac{\sqrt{x} - 1} 
{\sqrt{x} + 1}
\]}
{\(f'(x) = \frac{1} 
{\sqrt{x}(\sqrt{x}+1)^{2}} ,\ x > 0\)}{\(f'(x) = \frac{\sqrt{x}} 
{(\sqrt{x}+1)^{2}} ,\ x > 0\)}{\(f'(x) = \frac{2} 
{x(\sqrt{x}+1)^{2}} ,\ x > 0\)}{\(f'(x) = \frac{1} 
{(\sqrt{x}+1)^{2}} ,\ x > 0\)}{}{}}
\LANGpl{
\Question{
Wyznacz pochodną funkcji. 
\[ 
 f\colon y = \frac{\sqrt{x} - 1} 
{\sqrt{x} + 1}
\]}
{\(f'(x) = \frac{1} 
{\sqrt{x}(\sqrt{x}+1)^{2}} ,\ x > 0\)}{\(f'(x) = \frac{\sqrt{x}} 
{(\sqrt{x}+1)^{2}} ,\ x > 0\)}{\(f'(x) = \frac{2} 
{x(\sqrt{x}+1)^{2}} ,\ x > 0\)}{\(f'(x) = \frac{1} 
{(\sqrt{x}+1)^{2}} ,\ x > 0\)}{}{}}
\LANGcs{
\Question{
Derivace funkce \(f\colon y = \frac{\sqrt{x}-1} 
{\sqrt{x}+1}\) 
je rovna:}
{\(f'(x) = \frac{1} 
{\sqrt{x}(\sqrt{x}+1)^{2}} ,\ x > 0\)}{\(f'(x) = \frac{\sqrt{x}} 
{(\sqrt{x}+1)^{2}} ,\ x > 0\)}{\(f'(x) = \frac{2} 
{x(\sqrt{x}+1)^{2}} ,\ x > 0\)}{\(f'(x) = \frac{1} 
{(\sqrt{x}+1)^{2}} ,\ x > 0\)}{}{}}
\LANGsk{
\Question{
Derivácia funkcie \(f\colon y = \frac{\sqrt{x}-1} 
{\sqrt{x}+1}\) 
je rovná:}
{\(f'(x) = \frac{1} 
{\sqrt{x}(\sqrt{x}+1)^{2}} ,\ x > 0\)}{\(f'(x) = \frac{\sqrt{x}} 
{(\sqrt{x}+1)^{2}} ,\ x > 0\)}{\(f'(x) = \frac{2} 
{x(\sqrt{x}+1)^{2}} ,\ x > 0\)}{\(f'(x) = \frac{1} 
{(\sqrt{x}+1)^{2}} ,\ x > 0\)}{}{}}
\LANGes{
\Question{
Deriva la siguiente función.
\[ 
 f(x) = \frac{\sqrt{x} - 1} 
{\sqrt{x} + 1}
\]}
{\(f'(x) = \frac{1} 
{\sqrt{x}(\sqrt{x}+1)^{2}} ,\ x > 0\)}{\(f'(x) = \frac{\sqrt{x}} 
{(\sqrt{x}+1)^{2}} ,\ x > 0\)}{\(f'(x) = \frac{2} 
{x(\sqrt{x}+1)^{2}} ,\ x > 0\)}{\(f'(x) = \frac{1} 
{(\sqrt{x}+1)^{2}} ,\ x > 0\)}{}{}}
\End

\Data\ID{9000063106}
\Updated{2020-07-15 03:07:34}
\Level{B}
\SubArea{Derivative}
\LANGen{
\Question{
Differentiate the following function. 
\[ 
 f(x) =\sin x\cos x
\]}
{\(f'(x) =\cos 2x,\ x\in \mathbb{R}\)}{\(f'(x) = 1,\ x\in \mathbb{R}\)}{\(f'(x) = -\cos 2x,\ x\in \mathbb{R}\)}{\(f'(x) = -\sin x\cos x,\ x\in \mathbb{R}\)}{}{}}
\LANGpl{
\Question{
Wyznacz pochodną funkcji. 
\[ 
 f\colon y =\sin x\cos x
\]}
{\(f'(x) =\cos 2x,\ x\in \mathbb{R}\)}{\(f'(x) = 1,\ x\in \mathbb{R}\)}{\(f'(x) = -\cos 2x,\ x\in \mathbb{R}\)}{\(f'(x) = -\sin x\cos x,\ x\in \mathbb{R}\)}{}{}}
\LANGcs{
\Question{
Derivace funkce \(f\colon y =\sin x\cos x\) 
je rovna:}
{\(f'(x) =\cos 2x,\ x\in \mathbb{R}\)}{\(f'(x) = 1,\ x\in \mathbb{R}\)}{\(f'(x) = -\cos 2x,\ x\in \mathbb{R}\)}{\(f'(x) = -\sin x\cos x,\ x\in \mathbb{R}\)}{}{}}
\LANGsk{
\Question{
Derivácia funkcie \(f\colon y =\sin x\cos x\) 
je rovná:}
{\(f'(x) =\cos 2x,\ x\in \mathbb{R}\)}{\(f'(x) = 1,\ x\in \mathbb{R}\)}{\(f'(x) = -\cos 2x,\ x\in \mathbb{R}\)}{\(f'(x) = -\sin x\cos x,\ x\in \mathbb{R}\)}{}{}}
\LANGes{
\Question{
Deriva la siguiente función.
\[ 
 f(x) =\sin x\cos x
\]}
{\(f'(x) =\cos 2x,\ x\in \mathbb{R}\)}{\(f'(x) = 1,\ x\in \mathbb{R}\)}{\(f'(x) = -\cos 2x,\ x\in \mathbb{R}\)}{\(f'(x) = -\sin x\cos x,\ x\in \mathbb{R}\)}{}{}}
\End

\Data\ID{9000063107}
\Updated{2020-07-15 03:07:34}
\Level{B}
\SubArea{Derivative}
\LANGen{
\Question{
Differentiate the following function. 
\[ 
 f(x) =\cos x(1 +\sin x)
\]}
{\(f'(x) =\cos ^{2}x -\sin ^{2}x -\sin x,\ x\in \mathbb{R}\)}{\(f'(x) = -\sin x\cos x,\ x\in \mathbb{R}\)}{\(f'(x) =\cos x,\ x\in \mathbb{R}\)}{\(f'(x) =\sin x +\sin ^{2}x -\cos ^{2}x,\ x\in \mathbb{R}\)}{}{}}
\LANGpl{
\Question{
Wyznacz pochodną funkcji. 
\[ 
 f\colon y =\cos x(1 +\sin x)
\]}
{\(f'(x) =\cos ^{2}x -\sin ^{2}x -\sin x,\ x\in \mathbb{R}\)}{\(f'(x) = -\sin x\cos x,\ x\in \mathbb{R}\)}{\(f'(x) =\cos x,\ x\in \mathbb{R}\)}{\(f'(x) =\sin x +\sin ^{2}x -\cos ^{2}x,\ x\in \mathbb{R}\)}{}{}}
\LANGcs{
\Question{
Derivace funkce \(f\colon y =\cos x(1 +\sin x)\) 
je rovna:}
{\(f'(x) =\cos ^{2}x -\sin ^{2}x -\sin x,\ x\in \mathbb{R}\)}{\(f'(x) = -\sin x\cos x,\ x\in \mathbb{R}\)}{\(f'(x) =\cos x,\ x\in \mathbb{R}\)}{\(f'(x) =\sin x +\sin ^{2}x -\cos ^{2}x,\ x\in \mathbb{R}\)}{}{}}
\LANGsk{
\Question{
Derivácia funkcie \(f\colon y =\cos x(1 +\sin x)\) 
je rovná:}
{\(f'(x) =\cos ^{2}x -\sin ^{2}x -\sin x,\ x\in \mathbb{R}\)}{\(f'(x) = -\sin x\cos x,\ x\in \mathbb{R}\)}{\(f'(x) =\cos x,\ x\in \mathbb{R}\)}{\(f'(x) =\sin x +\sin ^{2}x -\cos ^{2}x,\ x\in \mathbb{R}\)}{}{}}
\LANGes{
\Question{
Deriva la siguiente función.
\[ 
 f(x) =\cos x(1 +\sin x)
\]}
{\(f'(x) =\cos ^{2}x -\sin ^{2}x -\sin x,\ x\in \mathbb{R}\)}{\(f'(x) = -\sin x\cos x,\ x\in \mathbb{R}\)}{\(f'(x) =\cos x,\ x\in \mathbb{R}\)}{\(f'(x) =\sin x +\sin ^{2}x -\cos ^{2}x,\ x\in \mathbb{R}\)}{}{}}
\End

\Data\ID{9000063108}
\Updated{2020-07-15 03:07:34}
\Level{B}
\SubArea{Derivative}
\LANGen{
\Question{
Differentiate the following function. 
\[ 
 f(x) = x^{5}\mathrm{e}^{x}
\]}
{\(f'(x) = x^{4}\mathrm{e}^{x}(5 + x),\ x\in \mathbb{R}\)}{\(f'(x) = 5x^{4}\mathrm{e}^{x},\ x\in \mathbb{R}\)}{\(f'(x) = x^{4}\mathrm{e}^{x}(x - 5),\ x\in \mathbb{R}\)}{\(f'(x) = x^{4}\mathrm{e}^{x}(5 + x^{2}),\ x\in \mathbb{R}\)}{}{}}
\LANGpl{
\Question{
Wyznacz pochodną funkcji. 
\[ 
 f\colon y = x^{5}\mathrm{e}^{x}
\]}
{\(f'(x) = x^{4}\mathrm{e}^{x}(5 + x),\ x\in \mathbb{R}\)}{\(f'(x) = 5x^{4}\mathrm{e}^{x},\ x\in \mathbb{R}\)}{\(f'(x) = x^{4}\mathrm{e}^{x}(x - 5),\ x\in \mathbb{R}\)}{\(f'(x) = x^{4}\mathrm{e}^{x}(5 + x^{2}),\ x\in \mathbb{R}\)}{}{}}
\LANGcs{
\Question{
Derivace funkce \(f\colon y = x^{5}\mathrm{e}^{x}\) 
je rovna:}
{\(f'(x) = x^{4}\mathrm{e}^{x}(5 + x),\ x\in \mathbb{R}\)}{\(f'(x) = 5x^{4}\mathrm{e}^{x},\ x\in \mathbb{R}\)}{\(f'(x) = x^{4}\mathrm{e}^{x}(x - 5),\ x\in \mathbb{R}\)}{\(f'(x) = x^{4}\mathrm{e}^{x}(5 + x^{2}),\ x\in \mathbb{R}\)}{}{}}
\LANGsk{
\Question{
Derivácia funkcie \(f\colon y = x^{5}\mathrm{e}^{x}\) 
je rovná:}
{\(f'(x) = x^{4}\mathrm{e}^{x}(5 + x),\ x\in \mathbb{R}\)}{\(f'(x) = 5x^{4}\mathrm{e}^{x},\ x\in \mathbb{R}\)}{\(f'(x) = x^{4}\mathrm{e}^{x}(x - 5),\ x\in \mathbb{R}\)}{\(f'(x) = x^{4}\mathrm{e}^{x}(5 + x^{2}),\ x\in \mathbb{R}\)}{}{}}
\LANGes{
\Question{
Deriva la siguiente función.
\[ 
 f(x) = x^{5}\mathrm{e}^{x}
\]}
{\(f'(x) = x^{4}\mathrm{e}^{x}(5 + x),\ x\in \mathbb{R}\)}{\(f'(x) = 5x^{4}\mathrm{e}^{x},\ x\in \mathbb{R}\)}{\(f'(x) = x^{4}\mathrm{e}^{x}(x - 5),\ x\in \mathbb{R}\)}{\(f'(x) = x^{4}\mathrm{e}^{x}(5 + x^{2}),\ x\in \mathbb{R}\)}{}{}}
\End

\Data\ID{9000063109}
\Updated{2020-07-15 03:07:34}
\Level{B}
\SubArea{Derivative}
\LANGen{
\Question{
Differentiate the following function. 
\[ 
 f(x) = 3^{x}\cdot x^{3}
\]}
{\(f'(x) = 3^{x}x^{2}(x\ln 3 + 3),\ x\in \mathbb{R}\)}{\(f'(x) = 3^{x+1}x^{2}\ln 3,\ x\in \mathbb{R}\)}{\(f'(x) = 3^{x}x^{2}(x + 3),\ x\in \mathbb{R}\)}{\(f'(x) = 3^{x}x^{2}(x\ln x + 3),\ x\in \mathbb{R}^{+}\)}{}{}}
\LANGpl{
\Question{
Wyznacz pochodną funkcji. 
\[ 
 f\colon y = 3^{x}\cdot x^{3}
\]}
{\(f'(x) = 3^{x}x^{2}(x\ln 3 + 3),\ x\in \mathbb{R}\)}{\(f'(x) = 3^{x+1}x^{2}\ln 3,\ x\in \mathbb{R}\)}{\(f'(x) = 3^{x}x^{2}(x + 3),\ x\in \mathbb{R}\)}{\(f'(x) = 3^{x}x^{2}(x\ln x + 3),\ x\in \mathbb{R}^{+}\)}{}{}}
\LANGcs{
\Question{
Derivace funkce \(f\colon y = 3^{x}\cdot x^{3}\) 
je rovna:}
{\(f'(x) = 3^{x}x^{2}(x\ln 3 + 3),\ x\in \mathbb{R}\)}{\(f'(x) = 3^{x+1}x^{2}\ln 3,\ x\in \mathbb{R}\)}{\(f'(x) = 3^{x}x^{2}(x + 3),\ x\in \mathbb{R}\)}{\(f'(x) = 3^{x}x^{2}(x\ln x + 3),\ x\in \mathbb{R}^{+}\)}{}{}}
\LANGsk{
\Question{
Derivácie funkcie \(f\colon y = 3^{x}\cdot x^{3}\) 
je rovná:}
{\(f'(x) = 3^{x}x^{2}(x\ln 3 + 3),\ x\in \mathbb{R}\)}{\(f'(x) = 3^{x+1}x^{2}\ln 3,\ x\in \mathbb{R}\)}{\(f'(x) = 3^{x}x^{2}(x + 3),\ x\in \mathbb{R}\)}{\(f'(x) = 3^{x}x^{2}(x\ln x + 3),\ x\in \mathbb{R}^{+}\)}{}{}}
\LANGes{
\Question{
Deriva la siguiente función.
\[ 
 f(x) = 3^{x}\cdot x^{3}
\]}
{\(f'(x) = 3^{x}x^{2}(x\ln 3 + 3),\ x\in \mathbb{R}\)}{\(f'(x) = 3^{x+1}x^{2}\ln 3,\ x\in \mathbb{R}\)}{\(f'(x) = 3^{x}x^{2}(x + 3),\ x\in \mathbb{R}\)}{\(f'(x) = 3^{x}x^{2}(x\ln x + 3),\ x\in \mathbb{R}^{+}\)}{}{}}
\End

\Data\ID{9000063110}
\Updated{2020-11-12 12:11:59}
\Level{B}
\SubArea{Derivative}
\LANGen{
\Question{
Differentiate the following function. 
\[ 
 f(x) =\sin x(1 +\mathop{\mathrm{tg}}\nolimits x)
\]}
{\(f'(x) =\cos x +\sin x + \frac{\sin x} 
{\cos ^{2}x},\ x\in R\setminus\{\frac{\pi}{2}+k\pi; k\in Z\}\)}{\(f'(x) =\cos x +\sin x,\ x\in R\setminus\{\frac{\pi}{2}+k\pi; k\in Z\}\)}{\(f'(x) = \frac{\sin x}
{\cos ^{2}x},\ x\in R\setminus\{\frac{\pi}{2}+k\pi; k\in Z\}\)}{\(f'(x) =\cos x + 2\sin x,\ x\in R\setminus\{\frac{\pi}{2}+k\pi; k\in Z\}\)}{}{}}
\LANGpl{
\Question{
Wyznacz pochodną funkcji. 
\[ 
 f\colon y =\sin x(1 +\mathop{\mathrm{tg}}\nolimits x)
\]}
{\(f'(x) =\cos x +\sin x + \frac{\sin x} 
{\cos ^{2}x},\ x\in R\setminus\{\frac{\pi}{2}+k\pi; k\in Z\}\)}{\(f'(x) =\cos x +\sin x,\ x\in R\setminus\{\frac{\pi}{2}+k\pi; k\in Z\}\)}{\(f'(x) = \frac{\sin x} 
{\cos ^{2}x},\ x\in R\setminus\{\frac{\pi}{2}+k\pi; k\in Z\}\)}{\(f'(x) =\cos x + 2\sin x,\ x\in R\setminus\{\frac{\pi}{2}+k\pi; k\in Z\}\)}{}{}}
\LANGcs{
\Question{
Derivace funkce \(f\colon y =\sin x(1 +\mathop{\mathrm{tg}}\nolimits x)\) 
je rovna:}
{\(f'(x) =\cos x +\sin x + \frac{\sin x} 
{\cos ^{2}x},\ x\in R\setminus\{\frac{\pi}{2}+k\pi; k\in Z\}\)}{\(f'(x) =\cos x +\sin x,\ x\in R\setminus\{\frac{\pi}{2}+k\pi; k\in Z\}\)}{\(f'(x) = \frac{\sin x} 
{\cos ^{2}x},\ x\in R\setminus\{\frac{\pi}{2}+k\pi; k\in Z\}\)}{\(f'(x) =\cos x + 2\sin x,\ x\in R\setminus\{\frac{\pi}{2}+k\pi; k\in Z\}\)}{}{}}
\LANGsk{
\Question{
Derivácia funkcie \(f\colon y =\sin x(1 +\mathop{\mathrm{tg}}\nolimits x)\) 
je rovná:}
{\(f'(x) =\cos x +\sin x + \frac{\sin x} 
{\cos ^{2}x},\ x\in R\setminus\{\frac{\pi}{2}+k\pi; k\in Z\}\)}{\(f'(x) =\cos x +\sin x,\ x\in R\setminus\{\frac{\pi}{2}+k\pi; k\in Z\}\)}{\(f'(x) = \frac{\sin x} 
{\cos ^{2}x},\ x\in R\setminus\{\frac{\pi}{2}+k\pi; k\in Z\}\)}{\(f'(x) =\cos x + 2\sin x,\ x\in R\setminus\{\frac{\pi}{2}+k\pi; k\in Z\}\)}{}{}}
\LANGes{
\Question{
Deriva la siguiente función.
\[ 
 f(x) =\sin x(1 +\mathop{\mathrm{tg}}\nolimits x)
\]}
{\(f'(x) =\cos x +\sin x + \frac{\sin x} 
{\cos ^{2}x},\ x\in R\setminus\{\frac{\pi}{2}+k\pi; k\in Z\}\)}{\(f'(x) =\cos x +\sin x,\ x\in R\setminus\{\frac{\pi}{2}+k\pi; k\in Z\}\)}{\(f'(x) = \frac{\sin x} 
{\cos ^{2}x},\ x\in R\setminus\{\frac{\pi}{2}+k\pi; k\in Z\}\)}{\(f'(x) =\cos x + 2\sin x,\ x\in R\setminus\{\frac{\pi}{2}+k\pi; k\in Z\}\)}{}{}}
\End

\Data\ID{9000063301}
\Updated{2020-07-15 03:07:37}
\Level{B}
\SubArea{Derivative}
\LANGen{
\Question{
Differentiate the following function. 
\[ 
 f(x)=\sin (2x^{2} + 1)
\]}
{\(f'(x) = 4x\cos (2x^{2} + 1),\ x\in \mathbb{R}\)}{\(f'(x) = 4x\cos x,\ x\in \mathbb{R}\)}{\(f'(x) =\cos (4x),\ x\in \mathbb{R}\)}{\(f'(x) =\sin (4x + 1),\ x\in \mathbb{R}\)}{}{}}
\LANGpl{
\Question{
Wyznacz pochodną funkcji. 
\[ 
 f\colon y =\sin (2x^{2} + 1)
\]}
{\(f'(x) = 4x\cos (2x^{2} + 1),\ x\in \mathbb{R}\)}{\(f'(x) = 4x\cos x,\ x\in \mathbb{R}\)}{\(f'(x) =\cos (4x),\ x\in \mathbb{R}\)}{\(f'(x) =\sin (4x + 1),\ x\in \mathbb{R}\)}{}{}}
\LANGcs{
\Question{
Derivace funkce \(f\colon y =\sin (2x^{2} + 1)\) 
je rovna:}
{\(f'(x) = 4x\cos (2x^{2} + 1),\ x\in \mathbb{R}\)}{\(f'(x) = 4x\cos x,\ x\in \mathbb{R}\)}{\(f'(x) =\cos (4x),\ x\in \mathbb{R}\)}{\(f'(x) =\sin (4x + 1),\ x\in \mathbb{R}\)}{}{}}
\LANGsk{
\Question{
Derivácia funkcie \(f\colon y =\sin (2x^{2} + 1)\) 
je rovná:}
{\(f'(x) = 4x\cos (2x^{2} + 1),\ x\in \mathbb{R}\)}{\(f'(x) = 4x\cos x,\ x\in \mathbb{R}\)}{\(f'(x) =\cos (4x),\ x\in \mathbb{R}\)}{\(f'(x) =\sin (4x + 1),\ x\in \mathbb{R}\)}{}{}}
\LANGes{
\Question{
Deriva la siguiente función.
\[ 
 f(x)=\sin (2x^{2} + 1)
\]}
{\(f'(x) = 4x\cos (2x^{2} + 1),\ x\in \mathbb{R}\)}{\(f'(x) = 4x\cos x,\ x\in \mathbb{R}\)}{\(f'(x) =\cos (4x),\ x\in \mathbb{R}\)}{\(f'(x) =\sin (4x + 1),\ x\in \mathbb{R}\)}{}{}}
\End

\Data\ID{9000063302}
\Updated{2020-07-15 03:07:37}
\Level{B}
\SubArea{Derivative}
\LANGen{
\Question{
Differentiate the following function. 
\[ 
 f(x)= (3x^{2} + 2)^{3}
\]}
{\(f'(x) = 18x(3x^{2} + 2)^{2},\ x\in \mathbb{R}\)}{\(f'(x) = 18x(3x^{2} + 2),\ x\in \mathbb{R}\)}{\(f'(x) = 18x^{2}(3x + 2)^{2},\ x\in \mathbb{R}\)}{\(f'(x) = 108x^{2},\ x\in \mathbb{R}\)}{}{}}
\LANGpl{
\Question{
Wyznacz pochodną funkcji. 
\[ 
 f\colon y = (3x^{2} + 2)^{3}
\]}
{\(f'(x) = 18x(3x^{2} + 2)^{2},\ x\in \mathbb{R}\)}{\(f'(x) = 18x(3x^{2} + 2),\ x\in \mathbb{R}\)}{\(f'(x) = 18x^{2}(3x + 2)^{2},\ x\in \mathbb{R}\)}{\(f'(x) = 108x^{2},\ x\in \mathbb{R}\)}{}{}}
\LANGcs{
\Question{
Derivace funkce \(f\colon y = (3x^{2} + 2)^{3}\) 
je rovna:}
{\(f'(x) = 18x(3x^{2} + 2)^{2},\ x\in \mathbb{R}\)}{\(f'(x) = 18x(3x^{2} + 2),\ x\in \mathbb{R}\)}{\(f'(x) = 18x^{2}(3x + 2)^{2},\ x\in \mathbb{R}\)}{\(f'(x) = 108x^{2},\ x\in \mathbb{R}\)}{}{}}
\LANGsk{
\Question{
Derivácia funkcie \(f\colon y = (3x^{2} + 2)^{3}\) 
je rovná:}
{\(f'(x) = 18x(3x^{2} + 2)^{2},\ x\in \mathbb{R}\)}{\(f'(x) = 18x(3x^{2} + 2),\ x\in \mathbb{R}\)}{\(f'(x) = 18x^{2}(3x + 2)^{2},\ x\in \mathbb{R}\)}{\(f'(x) = 108x^{2},\ x\in \mathbb{R}\)}{}{}}
\LANGes{
\Question{
Deriva la siguiente función.
\[ 
 f(x)= (3x^{2} + 2)^{3}
\]}
{\(f'(x) = 18x(3x^{2} + 2)^{2},\ x\in \mathbb{R}\)}{\(f'(x) = 18x(3x^{2} + 2),\ x\in \mathbb{R}\)}{\(f'(x) = 18x^{2}(3x + 2)^{2},\ x\in \mathbb{R}\)}{\(f'(x) = 108x^{2},\ x\in \mathbb{R}\)}{}{}}
\End

\Data\ID{9000070701}
\Updated{2020-07-15 03:07:37}
\Level{B}
\SubArea{Derivative}
\LANGen{
\Question{
Differentiate the following function. 
\[ 
 f(x)= (2x - 5)^{-6}
\]}
{\(f^{\prime}(x) = - \frac{12}
{(2x-5)^{7}} ;\ x\in \mathbb{R}\setminus \left \{\frac{5}
{2}\right \}\)}{\(f^{\prime}(x) = - \frac{12}
{(2x-5)^{7}} ;\ x\in \mathbb{R}\)}{\(f^{\prime}(x) = - \frac{12}
{(2x-5)^{5}} ;\ x\in \mathbb{R}\setminus \left \{\frac{5}
{2}\right \}\)}{\(f^{\prime}(x) = - \frac{12}
{(2x-5)^{5}} ;\ x\in \left (\frac{5}
{2};\infty \right )\)}{}{}}
\LANGpl{
\Question{
Wyznacz pochodną funkcji.
\[ 
 f\colon y = (2x - 5)^{-6}
\]}
{\(f^{\prime}(x) = - \frac{12}
{(2x-5)^{7}} ;\ x\in \mathbb{R}\setminus \left \{\frac{5}
{2}\right \}\)}{\(f^{\prime}(x) = - \frac{12}
{(2x-5)^{7}} ;\ x\in \mathbb{R}\)}{\(f^{\prime}(x) = - \frac{12}
{(2x-5)^{5}} ;\ x\in \mathbb{R}\setminus \left \{\frac{5}
{2}\right \}\)}{\(f^{\prime}(x) = - \frac{12}
{(2x-5)^{5}} ;\ x\in \left (\frac{5}
{2};\infty \right )\)}{}{}}
\LANGcs{
\Question{
Určete první derivaci funkce \(f\colon y = (2x - 5)^{-6}\).}
{\(f^{\prime}(x) = - \frac{12}
{(2x-5)^{7}} ;\ x\in \mathbb{R}\setminus \left \{\frac{5}
{2}\right \}\)}{\(f^{\prime}(x) = - \frac{12}
{(2x-5)^{7}} ;\ x\in \mathbb{R}\)}{\(f^{\prime}(x) = - \frac{12}
{(2x-5)^{5}} ;\ x\in \mathbb{R}\setminus \left \{\frac{5}
{2}\right \}\)}{\(f^{\prime}(x) = - \frac{12}
{(2x-5)^{5}} ;\ x\in \left (\frac{5}
{2};\infty \right )\)}{}{}}
\LANGsk{
\Question{
Určte prvú deriváciu funkcie \(f\colon y = (2x - 5)^{-6}\).}
{\(f^{\prime}(x) = - \frac{12}
{(2x-5)^{7}} ;\ x\in \mathbb{R}\setminus \left \{\frac{5}
{2}\right \}\)}{\(f^{\prime}(x) = - \frac{12}
{(2x-5)^{7}} ;\ x\in \mathbb{R}\)}{\(f^{\prime}(x) = - \frac{12}
{(2x-5)^{5}} ;\ x\in \mathbb{R}\setminus \left \{\frac{5}
{2}\right \}\)}{\(f^{\prime}(x) = - \frac{12}
{(2x-5)^{5}} ;\ x\in \left (\frac{5}
{2};\infty \right )\)}{}{}}
\LANGes{
\Question{
Deriva la siguiente función.
\[ 
 f(x)= (2x - 5)^{-6}
\]}
{\(f^{\prime}(x) = - \frac{12}
{(2x-5)^{7}} ;\ x\in \mathbb{R}\setminus \left \{\frac{5}
{2}\right \}\)}{\(f^{\prime}(x) = - \frac{12}
{(2x-5)^{7}} ;\ x\in \mathbb{R}\)}{\(f^{\prime}(x) = - \frac{12}
{(2x-5)^{5}} ;\ x\in \mathbb{R}\setminus \left \{\frac{5}
{2}\right \}\)}{\(f^{\prime}(x) = - \frac{12}
{(2x-5)^{5}} ;\ x\in \left (\frac{5}
{2};\infty \right )\)}{}{}}
\End

\Data\ID{9000070702}
\Updated{2020-07-15 03:07:37}
\Level{B}
\SubArea{Derivative}
\LANGen{
\Question{
Differentiate the following function. 
\[ 
 f(x) = (x^{2} - 3x + 2)^{\frac{1} 
{2} }
\]}
{\(f^{\prime}(x) = \frac{2x-3} 
{2\sqrt{x^{2 } -3x+2}};\ x\in \mathbb{R}\setminus \left [ 1;2\right ] \)}{\(f^{\prime}(x) = \frac{2x-3} 
{2\sqrt{x^{2 } -3x+2}};\ x\in \mathbb{R}\setminus \left (1;2\right )\)}{\(f^{\prime}(x) = (4x - 6)\sqrt{x^{2 } - 3x + 2};\ x\in \mathbb{R}\setminus \left [ 1;2\right ] \)}{\(f^{\prime}(x) = (4x - 6)\sqrt{x^{2 } - 3x + 2};\ x\in \mathbb{R}\setminus \left (1;2\right )\)}{}{}}
\LANGpl{
\Question{
Wyznacz pochodną funkcji.
\[ 
 f\colon y = (x^{2} - 3x + 2)^{\frac{1} 
{2} }
\]}
{\(f^{\prime}(x) = \frac{2x-3} 
{2\sqrt{x^{2 } -3x+2}};\ x\in \mathbb{R}\setminus \left [ 1;2\right ] \)}{\(f^{\prime}(x) = \frac{2x-3} 
{2\sqrt{x^{2 } -3x+2}};\ x\in \mathbb{R}\setminus \left (1;2\right )\)}{\(f^{\prime}(x) = (4x - 6)\sqrt{x^{2 } - 3x + 2};\ x\in \mathbb{R}\setminus \left [ 1;2\right ] \)}{\(f^{\prime}(x) = (4x - 6)\sqrt{x^{2 } - 3x + 2};\ x\in \mathbb{R}\setminus \left (1;2\right )\)}{}{}}
\LANGcs{
\Question{
Určete první derivaci funkce \(f\colon y = (x^{2} - 3x + 2)^{\frac{1} 
{2} }\).}
{\(f^{\prime}(x) = \frac{2x-3} 
{2\sqrt{x^{2 } -3x+2}};\ x\in \mathbb{R}\setminus \left \langle 1;2\right \rangle \)}{\(f^{\prime}(x) = \frac{2x-3} 
{2\sqrt{x^{2 } -3x+2}};\ x\in \mathbb{R}\setminus \left (1;2\right )\)}{\(f^{\prime}(x) = (4x - 6)\sqrt{x^{2 } - 3x + 2};\ x\in \mathbb{R}\setminus \left \langle 1;2\right \rangle \)}{\(f^{\prime}(x) = (4x - 6)\sqrt{x^{2 } - 3x + 2};\ x\in \mathbb{R}\setminus \left (1;2\right )\)}{}{}}
\LANGsk{
\Question{
Určte prvú deriváciu funkcie \(f\colon y = (x^{2} - 3x + 2)^{\frac{1} 
{2} }\).}
{\(f^{\prime}(x) = \frac{2x-3} 
{2\sqrt{x^{2 } -3x+2}};\ x\in \mathbb{R}\setminus \left \langle 1;2\right \rangle \)}{\(f^{\prime}(x) = \frac{2x-3} 
{2\sqrt{x^{2 } -3x+2}};\ x\in \mathbb{R}\setminus \left (1;2\right )\)}{\(f^{\prime}(x) = (4x - 6)\sqrt{x^{2 } - 3x + 2};\ x\in \mathbb{R}\setminus \left \langle 1;2\right \rangle \)}{\(f^{\prime}(x) = (4x - 6)\sqrt{x^{2 } - 3x + 2};\ x\in \mathbb{R}\setminus \left (1;2\right )\)}{}{}}
\LANGes{
\Question{
Deriva la siguiente función.
\[ 
 f(x) = (x^{2} - 3x + 2)^{\frac{1} 
{2} }
\]}
{\(f^{\prime}(x) = \frac{2x-3} 
{2\sqrt{x^{2 } -3x+2}};\ x\in \mathbb{R}\setminus \left [ 1;2\right ] \)}{\(f^{\prime}(x) = \frac{2x-3} 
{2\sqrt{x^{2 } -3x+2}};\ x\in \mathbb{R}\setminus \left (1;2\right )\)}{\(f^{\prime}(x) = (4x - 6)\sqrt{x^{2 } - 3x + 2};\ x\in \mathbb{R}\setminus \left [ 1;2\right ] \)}{\(f^{\prime}(x) = (4x - 6)\sqrt{x^{2 } - 3x + 2};\ x\in \mathbb{R}\setminus \left (1;2\right )\)}{}{}}
\End

\Data\ID{9000070703}
\Updated{2020-07-15 03:07:37}
\Level{B}
\SubArea{Derivative}
\LANGen{
\Question{
Differentiate the following function. 
\[ 
 f(x)= \sqrt{\sin x -\cos x}
\]}
{\(f^{\prime}(x) = \frac{\sin x+\cos x} 
{2\sqrt{\sin x-\cos x}};\ x\in \left ( \frac{\pi }{4} + 2k\pi ; \frac{5\pi } 
{4} + 2k\pi \right ),\ k\in \mathbb{Z}\)}{\(f^{\prime}(x) = \frac{\sin x+\cos x} 
{2\sqrt{\sin x-\cos x}};\ x\in \left [ \frac{\pi }{4} + 2k\pi ; \frac{5\pi } 
{4} + 2k\pi \right ] ,\ k\in \mathbb{Z}\)}{\(f^{\prime}(x) = \frac{\sin x-\cos x} 
{2\sqrt{\sin x-\cos x}};\ x\in \left [ \frac{\pi }{4} + 2k\pi ; \frac{5\pi } 
{4} + 2k\pi \right ] ,\ k\in \mathbb{Z}\)}{\(f^{\prime}(x) = \frac{\sin x-\cos x} 
{2\sqrt{\sin x-\cos x}};\ x\in \left ( \frac{\pi }{4} + 2k\pi ; \frac{5\pi } 
{4} + 2k\pi \right ),\ k\in \mathbb{Z}\)}{}{}}
\LANGpl{
\Question{
Wyznacz pochodną funkcji.
\[ 
 f\colon y = \sqrt{\sin x -\cos x}
\]}
{\(f^{\prime}(x) = \frac{\sin x+\cos x} 
{2\sqrt{\sin x-\cos x}};\ x\in \left ( \frac{\pi }{4} + 2k\pi ; \frac{5\pi } 
{4} + 2k\pi \right ),\ k\in \mathbb{Z}\)}{\(f^{\prime}(x) = \frac{\sin x+\cos x} 
{2\sqrt{\sin x-\cos x}};\ x\in \left [ \frac{\pi }{4} + 2k\pi ; \frac{5\pi } 
{4} + 2k\pi \right ] ,\ k\in \mathbb{Z}\)}{\(f^{\prime}(x) = \frac{\sin x-\cos x} 
{2\sqrt{\sin x-\cos x}};\ x\in \left [ \frac{\pi }{4} + 2k\pi ; \frac{5\pi } 
{4} + 2k\pi \right ] ,\ k\in \mathbb{Z}\)}{\(f^{\prime}(x) = \frac{\sin x-\cos x} 
{2\sqrt{\sin x-\cos x}};\ x\in \left ( \frac{\pi }{4} + 2k\pi ; \frac{5\pi } 
{4} + 2k\pi \right ),\ k\in \mathbb{Z}\)}{}{}}
\LANGcs{
\Question{
Určete první derivaci funkce \(f\colon y = \sqrt{\sin x -\cos x}\).}
{\(f^{\prime}(x) = \frac{\sin x+\cos x} 
{2\sqrt{\sin x-\cos x}};\ x\in \left ( \frac{\pi }{4} + 2k\pi ; \frac{5\pi } 
{4} + 2k\pi \right ),\ k\in \mathbb{Z}\)}{\(f^{\prime}(x) = \frac{\sin x+\cos x} 
{2\sqrt{\sin x-\cos x}};\ x\in \left \langle \frac{\pi }{4} + 2k\pi ; \frac{5\pi } 
{4} + 2k\pi \right \rangle ,\ k\in \mathbb{Z}\)}{\(f^{\prime}(x) = \frac{\sin x-\cos x} 
{2\sqrt{\sin x-\cos x}};\ x\in \left \langle \frac{\pi }{4} + 2k\pi ; \frac{5\pi } 
{4} + 2k\pi \right \rangle ,\ k\in \mathbb{Z}\)}{\(f^{\prime}(x) = \frac{\sin x-\cos x} 
{2\sqrt{\sin x-\cos x}};\ x\in \left ( \frac{\pi }{4} + 2k\pi ; \frac{5\pi } 
{4} + 2k\pi \right ),\ k\in \mathbb{Z}\)}{}{}}
\LANGsk{
\Question{
Určte prvú deriváciu funkcie \(f\colon y = \sqrt{\sin x -\cos x}\).}
{\(f^{\prime}(x) = \frac{\sin x+\cos x} 
{2\sqrt{\sin x-\cos x}};\ x\in \left ( \frac{\pi }{4} + 2k\pi ; \frac{5\pi } 
{4} + 2k\pi \right ),\ k\in \mathbb{Z}\)}{\(f^{\prime}(x) = \frac{\sin x+\cos x} 
{2\sqrt{\sin x-\cos x}};\ x\in \left \langle \frac{\pi }{4} + 2k\pi ; \frac{5\pi } 
{4} + 2k\pi \right \rangle ,\ k\in \mathbb{Z}\)}{\(f^{\prime}(x) = \frac{\sin x-\cos x} 
{2\sqrt{\sin x-\cos x}};\ x\in \left \langle \frac{\pi }{4} + 2k\pi ; \frac{5\pi } 
{4} + 2k\pi \right \rangle ,\ k\in \mathbb{Z}\)}{\(f^{\prime}(x) = \frac{\sin x-\cos x} 
{2\sqrt{\sin x-\cos x}};\ x\in \left ( \frac{\pi }{4} + 2k\pi ; \frac{5\pi } 
{4} + 2k\pi \right ),\ k\in \mathbb{Z}\)}{}{}}
\LANGes{
\Question{
Deriva la siguiente función.
\[ 
 f(x)= \sqrt{\sin x -\cos x}
\]}
{\(f^{\prime}(x) = \frac{\sin x+\cos x} 
{2\sqrt{\sin x-\cos x}};\ x\in \left ( \frac{\pi }{4} + 2k\pi ; \frac{5\pi } 
{4} + 2k\pi \right ),\ k\in \mathbb{Z}\)}{\(f^{\prime}(x) = \frac{\sin x+\cos x} 
{2\sqrt{\sin x-\cos x}};\ x\in \left [ \frac{\pi }{4} + 2k\pi ; \frac{5\pi } 
{4} + 2k\pi \right ] ,\ k\in \mathbb{Z}\)}{\(f^{\prime}(x) = \frac{\sin x-\cos x} 
{2\sqrt{\sin x-\cos x}};\ x\in \left [ \frac{\pi }{4} + 2k\pi ; \frac{5\pi } 
{4} + 2k\pi \right ] ,\ k\in \mathbb{Z}\)}{\(f^{\prime}(x) = \frac{\sin x-\cos x} 
{2\sqrt{\sin x-\cos x}};\ x\in \left ( \frac{\pi }{4} + 2k\pi ; \frac{5\pi } 
{4} + 2k\pi \right ),\ k\in \mathbb{Z}\)}{}{}}
\End

\Data\ID{9000070704}
\Updated{2020-07-15 03:07:37}
\Level{B}
\SubArea{Derivative}
\LANGen{
\Question{
Differentiate the following function. 
\[ 
 f(x) = \frac{1} 
{\cos x + 3x^{2}}
\]}
{\(f^{\prime}(x) = \frac{\sin x-6x} 
{(3x^{2}+\cos x)^{2}} ;\ x\in \mathbb{R}\)}{\(f^{\prime}(x) = \frac{6x-\sin x} 
{(3x^{2}+\cos x)^{2}} ;\ x\in \mathbb{R}\)}{\(f^{\prime}(x) = \frac{\sin x-6x} 
{3x^{2}+\cos x};\ x\in \mathbb{R}\)}{\(f^{\prime}(x) = \frac{6x-\sin x} 
{3x^{2}+\cos x};\ x\in \mathbb{R}\)}{}{}}
\LANGpl{
\Question{
Wyznacz pochodną funkcji.
\[ 
 f\colon y = \frac{1} 
{\cos x + 3x^{2}}
\]}
{\(f^{\prime}(x) = \frac{\sin x-6x} 
{(3x^{2}+\cos x)^{2}} ;\ x\in \mathbb{R}\)}{\(f^{\prime}(x) = \frac{6x-\sin x} 
{(3x^{2}+\cos x)^{2}} ;\ x\in \mathbb{R}\)}{\(f^{\prime}(x) = \frac{\sin x-6x} 
{3x^{2}+\cos x};\ x\in \mathbb{R}\)}{\(f^{\prime}(x) = \frac{6x-\sin x} 
{3x^{2}+\cos x};\ x\in \mathbb{R}\)}{}{}}
\LANGcs{
\Question{
Určete první derivaci funkce \(f\colon y = \frac{1} 
{\cos x+3x^{2}} \).}
{\(f^{\prime}(x) = \frac{\sin x-6x} 
{(3x^{2}+\cos x)^{2}} ;\ x\in \mathbb{R}\)}{\(f^{\prime}(x) = \frac{6x-\sin x} 
{(3x^{2}+\cos x)^{2}} ;\ x\in \mathbb{R}\)}{\(f^{\prime}(x) = \frac{\sin x-6x} 
{3x^{2}+\cos x};\ x\in \mathbb{R}\)}{\(f^{\prime}(x) = \frac{6x-\sin x} 
{3x^{2}+\cos x};\ x\in \mathbb{R}\)}{}{}}
\LANGsk{
\Question{
Určte prvú deriváciu funkcie \(f\colon y = \frac{1} 
{\cos x+3x^{2}} \).}
{\(f^{\prime}(x) = \frac{\sin x-6x} 
{(3x^{2}+\cos x)^{2}} ;\ x\in \mathbb{R}\)}{\(f^{\prime}(x) = \frac{6x-\sin x} 
{(3x^{2}+\cos x)^{2}} ;\ x\in \mathbb{R}\)}{\(f^{\prime}(x) = \frac{\sin x-6x} 
{3x^{2}+\cos x};\ x\in \mathbb{R}\)}{\(f^{\prime}(x) = \frac{6x-\sin x} 
{3x^{2}+\cos x};\ x\in \mathbb{R}\)}{}{}}
\LANGes{
\Question{
Deriva la siguiente función.
\[ 
 f(x) = \frac{1} 
{\cos x + 3x^{2}}
\]}
{\(f^{\prime}(x) = \frac{\sin x-6x} 
{(3x^{2}+\cos x)^{2}} ;\ x\in \mathbb{R}\)}{\(f^{\prime}(x) = \frac{6x-\sin x} 
{(3x^{2}+\cos x)^{2}} ;\ x\in \mathbb{R}\)}{\(f^{\prime}(x) = \frac{\sin x-6x} 
{3x^{2}+\cos x};\ x\in \mathbb{R}\)}{\(f^{\prime}(x) = \frac{6x-\sin x} 
{3x^{2}+\cos x};\ x\in \mathbb{R}\)}{}{}}
\End

\Data\ID{9000070705}
\Updated{2020-07-15 03:07:37}
\Level{B}
\SubArea{Derivative}
\LANGen{
\Question{
Differentiate the following function. 
\[ 
 f(x) =\ln (2x^{2} + 5x)
\]}
{\(f^{\prime}(x) = \frac{4x+5} 
{2x^{2}+5x};\ x\in \left (-\infty ;-\frac{5}
{2}\right )\cup \left (0;\infty \right )\)}{\(f^{\prime}(x) = \frac{4x+5} 
{2x^{2}+5x};\ x\in \mathbb{R}\setminus \left \{-\frac{5}
{2};0\right \}\)}{\(f^{\prime}(x) = \frac{1} 
{2x^{2}+5x};\ x\in \left (-\infty ;-\frac{5}
{2}\right )\cup \left (0;\infty \right )\)}{\(f^{\prime}(x) = \frac{1} 
{2x^{2}+5x};\ x\in \mathbb{R}\setminus \left \{-\frac{5}
{2};0\right \}\)}{}{}}
\LANGpl{
\Question{
Wyznacz pochodną funkcji.
\[ 
 f\colon y =\ln (2x^{2} + 5x)
\]}
{\(f^{\prime}(x) = \frac{4x+5} 
{2x^{2}+5x};\ x\in \left (-\infty ;-\frac{5}
{2}\right )\cup \left (0;\infty \right )\)}{\(f^{\prime}(x) = \frac{4x+5} 
{2x^{2}+5x};\ x\in \mathbb{R}\setminus \left \{-\frac{5}
{2};0\right \}\)}{\(f^{\prime}(x) = \frac{1} 
{2x^{2}+5x};\ x\in \left (-\infty ;-\frac{5}
{2}\right )\cup \left (0;\infty \right )\)}{\(f^{\prime}(x) = \frac{1} 
{2x^{2}+5x};\ x\in \mathbb{R}\setminus \left \{-\frac{5}
{2};0\right \}\)}{}{}}
\LANGcs{
\Question{
Určete první derivaci funkce \(f\colon y =\ln (2x^{2} + 5x)\).}
{\(f^{\prime}(x) = \frac{4x+5} 
{2x^{2}+5x};\ x\in \left (-\infty ;-\frac{5}
{2}\right )\cup \left (0;\infty \right )\)}{\(f^{\prime}(x) = \frac{4x+5} 
{2x^{2}+5x};\ x\in \mathbb{R}\setminus \left \{-\frac{5}
{2};0\right \}\)}{\(f^{\prime}(x) = \frac{1} 
{2x^{2}+5x};\ x\in \left (-\infty ;-\frac{5}
{2}\right )\cup \left (0;\infty \right )\)}{\(f^{\prime}(x) = \frac{1} 
{2x^{2}+5x};\ x\in \mathbb{R}\setminus \left \{-\frac{5}
{2};0\right \}\)}{}{}}
\LANGsk{
\Question{
Určte prvú deriváciu funkcie \(f\colon y =\ln (2x^{2} + 5x)\).}
{\(f^{\prime}(x) = \frac{4x+5} 
{2x^{2}+5x};\ x\in \left (-\infty ;-\frac{5}
{2}\right )\cup \left (0;\infty \right )\)}{\(f^{\prime}(x) = \frac{4x+5} 
{2x^{2}+5x};\ x\in \mathbb{R}\setminus \left \{-\frac{5}
{2};0\right \}\)}{\(f^{\prime}(x) = \frac{1} 
{2x^{2}+5x};\ x\in \left (-\infty ;-\frac{5}
{2}\right )\cup \left (0;\infty \right )\)}{\(f^{\prime}(x) = \frac{1} 
{2x^{2}+5x};\ x\in \mathbb{R}\setminus \left \{-\frac{5}
{2};0\right \}\)}{}{}}
\LANGes{
\Question{
Deriva la siguiente función.
\[ 
 f(x) =\ln (2x^{2} + 5x)
\]}
{\(f^{\prime}(x) = \frac{4x+5} 
{2x^{2}+5x};\ x\in \left (-\infty ;-\frac{5}
{2}\right )\cup \left (0;\infty \right )\)}{\(f^{\prime}(x) = \frac{4x+5} 
{2x^{2}+5x};\ x\in \mathbb{R}\setminus \left \{-\frac{5}
{2};0\right \}\)}{\(f^{\prime}(x) = \frac{1} 
{2x^{2}+5x};\ x\in \left (-\infty ;-\frac{5}
{2}\right )\cup \left (0;\infty \right )\)}{\(f^{\prime}(x) = \frac{1} 
{2x^{2}+5x};\ x\in \mathbb{R}\setminus \left \{-\frac{5}
{2};0\right \}\)}{}{}}
\End

\Data\ID{9000070706}
\Updated{2020-07-15 03:07:37}
\Level{B}
\SubArea{Derivative}
\LANGen{
\Question{
Differentiate the following function. 
\[ 
 f(x) = \sqrt{x^{2 } + 3x}
\]}
{\(f^{\prime}(x) = \frac{2x+3} 
{2\sqrt{x^{2 } +3x}};\ x\in \left (-\infty ;-3\right )\cup \left (0;\infty \right )\)}{\(f^{\prime}(x) = \frac{2x+3} 
{2\sqrt{x^{2 } +3x}};\ x\in \left (-\infty ;-3\right ] \cup \left [ 0;\infty \right )\)}{\(f^{\prime}(x) = \frac{2x+3} 
{\sqrt{x^{2 } +3x}};\ x\in \left (-\infty ;-3\right )\cup \left (0;\infty \right )\)}{\(f^{\prime}(x) = \frac{\sqrt{x^{2 } +3x}} 
 {2x+3} ;\ x\in \left (-\infty ;-3\right ] \cup \left [ 0;\infty \right )\)}{}{}}
\LANGpl{
\Question{
Wyznacz pochodną funkcji.
\[ 
 f\colon y = \sqrt{x^{2 } + 3x}
\]}
{\(f^{\prime}(x) = \frac{2x+3} 
{2\sqrt{x^{2 } +3x}};\ x\in \left (-\infty ;-3\right )\cup \left (0;\infty \right )\)}{\(f^{\prime}(x) = \frac{2x+3} 
{2\sqrt{x^{2 } +3x}};\ x\in \left (-\infty ;-3\right ] \cup \left [ 0;\infty \right )\)}{\(f^{\prime}(x) = \frac{2x+3} 
{\sqrt{x^{2 } +3x}};\ x\in \left (-\infty ;-3\right )\cup \left (0;\infty \right )\)}{\(f^{\prime}(x) = \frac{\sqrt{x^{2 } +3x}} 
 {2x+3} ;\ x\in \left (-\infty ;-3\right ] \cup \left [ 0;\infty \right )\)}{}{}}
\LANGcs{
\Question{
Určete první derivaci funkce \(f\colon y = \sqrt{x^{2 } + 3x}\).}
{\(f^{\prime}(x) = \frac{2x+3} 
{2\sqrt{x^{2 } +3x}};\ x\in \left (-\infty ;-3\right )\cup \left (0;\infty \right )\)}{\(f^{\prime}(x) = \frac{2x+3} 
{2\sqrt{x^{2 } +3x}};\ x\in \left (-\infty ;-3\right \rangle \cup \left \langle 0;\infty \right )\)}{\(f^{\prime}(x) = \frac{2x+3} 
{\sqrt{x^{2 } +3x}};\ x\in \left (-\infty ;-3\right )\cup \left (0;\infty \right )\)}{\(f^{\prime}(x) = \frac{\sqrt{x^{2 } +3x}} 
 {2x+3} ;\ x\in \left (-\infty ;-3\right \rangle \cup \left \langle 0;\infty \right )\)}{}{}}
\LANGsk{
\Question{
Určte prvú deriváciu funkcie \(f\colon y = \sqrt{x^{2 } + 3x}\).}
{\(f^{\prime}(x) = \frac{2x+3} 
{2\sqrt{x^{2 } +3x}};\ x\in \left (-\infty ;-3\right )\cup \left (0;\infty \right )\)}{\(f^{\prime}(x) = \frac{2x+3} 
{2\sqrt{x^{2 } +3x}};\ x\in \left (-\infty ;-3\right \rangle \cup \left \langle 0;\infty \right )\)}{\(f^{\prime}(x) = \frac{2x+3} 
{\sqrt{x^{2 } +3x}};\ x\in \left (-\infty ;-3\right )\cup \left (0;\infty \right )\)}{\(f^{\prime}(x) = \frac{\sqrt{x^{2 } +3x}} 
 {2x+3} ;\ x\in \left (-\infty ;-3\right \rangle \cup \left \langle 0;\infty \right )\)}{}{}}
\LANGes{
\Question{
Deriva la siguiente función.
\[ 
 f(x) = \sqrt{x^{2 } + 3x}
\]}
{\(f^{\prime}(x) = \frac{2x+3} 
{2\sqrt{x^{2 } +3x}};\ x\in \left (-\infty ;-3\right )\cup \left (0;\infty \right )\)}{\(f^{\prime}(x) = \frac{2x+3} 
{2\sqrt{x^{2 } +3x}};\ x\in \left (-\infty ;-3\right ] \cup \left [ 0;\infty \right )\)}{\(f^{\prime}(x) = \frac{2x+3} 
{\sqrt{x^{2 } +3x}};\ x\in \left (-\infty ;-3\right )\cup \left (0;\infty \right )\)}{\(f^{\prime}(x) = \frac{\sqrt{x^{2 } +3x}} 
 {2x+3} ;\ x\in \left (-\infty ;-3\right ] \cup \left [ 0;\infty \right )\)}{}{}}
\End

\Data\ID{9000070707}
\Updated{2020-07-15 03:07:37}
\Level{B}
\SubArea{Derivative}
\LANGen{
\Question{
Differentiate the following function. 
\[ 
 f(x) = \root{5}\of{x^{2} - 7x}
\]Remark: The function \(f\colon y = \root{5}\of{x}\) 
is defined for \(x\in \left < 0;\infty \right )\).}
{\(f^{\prime}(x) = \frac{2x-7} 
{5(x^{2}-7x)^{\frac{4} 
{5} }} ;\ x\in \left (-\infty ;0\right )\cup \left (7;\infty \right )\)}{\(f^{\prime}(x) = \frac{2x-7} 
{5(x^{2}-7x)^{\frac{4} 
{5} }} ;\ x\in \left (-\infty ;0\right ] \cup \left [ 7;\infty \right )\)}{\(f^{\prime}(x) = (2x - 7)\root{4}\of{x^{2} - 7x};\ x\in \left (-\infty ;0\right )\cup \left (7;\infty \right )\)}{\(f^{\prime}(x) = (2x - 7)\root{4}\of{x^{2} - 7x};\ x\in \left (-\infty ;0\right ] \cup \left [ 7;\infty \right )\)}{}{}}
\LANGpl{
\Question{
Wyznacz pochodną funkcji.
\[ 
 f\colon y = \root{5}\of{x^{2} - 7x}
\] Uwaga: Funkcja  \(f\colon y = \root{5}\of{x}\) 
jest określona dla \(x\in \left < 0;\infty \right )\).}
{\(f^{\prime}(x) = \frac{2x-7} 
{5(x^{2}-7x)^{\frac{4} 
{5} }} ;\ x\in \left (-\infty ;0\right )\cup \left (7;\infty \right )\)}{\(f^{\prime}(x) = \frac{2x-7} 
{5(x^{2}-7x)^{\frac{4} 
{5} }} ;\ x\in \left (-\infty ;0\right ] \cup \left [ 7;\infty \right )\)}{\(f^{\prime}(x) = (2x - 7)\root{4}\of{x^{2} - 7x};\ x\in \left (-\infty ;0\right )\cup \left (7;\infty \right )\)}{\(f^{\prime}(x) = (2x - 7)\root{4}\of{x^{2} - 7x};\ x\in \left (-\infty ;0\right ] \cup \left [ 7;\infty \right )\)}{}{}}
\LANGcs{
\Question{
Určete první derivaci funkce \(f\colon y = \root{5}\of{x^{2} - 7x}\). Poznámka: Funkce \(f\colon y = \root{5}\of{x}\) 
je definována pro \(x\in \left < 0;\infty \right )\).}
{\(f^{\prime}(x) = \frac{2x-7} 
{5(x^{2}-7x)^{\frac{4} 
{5} }} ;\ x\in \left (-\infty ;0\right )\cup \left (7;\infty \right )\)}{\(f^{\prime}(x) = \frac{2x-7} 
{5(x^{2}-7x)^{\frac{4} 
{5} }} ;\ x\in \left (-\infty ;0\right \rangle \cup \left \langle 7;\infty \right )\)}{\(f^{\prime}(x) = (2x - 7)\root{4}\of{x^{2} - 7x};\ x\in \left (-\infty ;0\right )\cup \left (7;\infty \right )\)}{\(f^{\prime}(x) = (2x - 7)\root{4}\of{x^{2} - 7x};\ x\in \left (-\infty ;0\right \rangle \cup \left \langle 7;\infty \right )\)}{}{}}
\LANGsk{
\Question{
Určte prvú deriváciu funkcie \(f\colon y = \root{5}\of{x^{2} - 7x}\). Poznámka: Funkcia \(f\colon y = \root{5}\of{x}\) 
je definovaná pre \(x\in \left < 0;\infty \right )\).}
{\(f^{\prime}(x) = \frac{2x-7} 
{5(x^{2}-7x)^{\frac{4} 
{5} }} ;\ x\in \left (-\infty ;0\right )\cup \left (7;\infty \right )\)}{\(f^{\prime}(x) = \frac{2x-7} 
{5(x^{2}-7x)^{\frac{4} 
{5} }} ;\ x\in \left (-\infty ;0\right \rangle \cup \left \langle 7;\infty \right )\)}{\(f^{\prime}(x) = (2x - 7)\root{4}\of{x^{2} - 7x};\ x\in \left (-\infty ;0\right )\cup \left (7;\infty \right )\)}{\(f^{\prime}(x) = (2x - 7)\root{4}\of{x^{2} - 7x};\ x\in \left (-\infty ;0\right \rangle \cup \left \langle 7;\infty \right )\)}{}{}}
\LANGes{
\Question{
Deriva la siguiente función.
\[ 
 f(x) = \root{5}\of{x^{2} - 7x}
\]Nota: La función \(f\colon y = \root{5}\of{x}\) 
está definida para \(x\in \left < 0;\infty \right )\).}
{\(f^{\prime}(x) = \frac{2x-7} 
{5(x^{2}-7x)^{\frac{4} 
{5} }} ;\ x\in \left (-\infty ;0\right )\cup \left (7;\infty \right )\)}{\(f^{\prime}(x) = \frac{2x-7} 
{5(x^{2}-7x)^{\frac{4} 
{5} }} ;\ x\in \left (-\infty ;0\right ] \cup \left [ 7;\infty \right )\)}{\(f^{\prime}(x) = (2x - 7)\root{4}\of{x^{2} - 7x};\ x\in \left (-\infty ;0\right )\cup \left (7;\infty \right )\)}{\(f^{\prime}(x) = (2x - 7)\root{4}\of{x^{2} - 7x};\ x\in \left (-\infty ;0\right ] \cup \left [ 7;\infty \right )\)}{}{}}
\End

\Data\ID{9000070708}
\Updated{2020-07-15 03:07:37}
\Level{B}
\SubArea{Derivative}
\LANGen{
\Question{
Differentiate the following function. 
\[ 
 f(x) =\ln \left (\frac{1 + x}
{1 - x}\right )
\]}
{\(f^{\prime}(x) = \frac{2} 
{1-x^{2}} ;\ x\in \left (-1;1\right )\)}{\(f^{\prime}(x) = \frac{2} 
{1-x^{2}} ;\ x\in \mathbb{R}\setminus \left \{-1;1\right \}\)}{\(f^{\prime}(x) = \frac{1-x} 
{1+x};\ x\in \left (-1;1\right )\)}{\(f^{\prime}(x) = \frac{1-x} 
{1+x};\ x\in \mathbb{R}\setminus \left \{-1;1\right \}\)}{}{}}
\LANGpl{
\Question{
Wyznacz pochodną funkcji.
\[ 
 f\colon y =\ln \left (\frac{1 + x}
{1 - x}\right )
\]}
{\(f^{\prime}(x) = \frac{2} 
{1-x^{2}} ;\ x\in \left (-1;1\right )\)}{\(f^{\prime}(x) = \frac{2} 
{1-x^{2}} ;\ x\in \mathbb{R}\setminus \left \{-1;1\right \}\)}{\(f^{\prime}(x) = \frac{1-x} 
{1+x};\ x\in \left (-1;1\right )\)}{\(f^{\prime}(x) = \frac{1-x} 
{1+x};\ x\in \mathbb{R}\setminus \left \{-1;1\right \}\)}{}{}}
\LANGcs{
\Question{
Určete první derivaci funkce \(f\colon y =\ln \left (\frac{1+x}
{1-x}\right )\).}
{\(f^{\prime}(x) = \frac{2} 
{1-x^{2}} ;\ x\in \left (-1;1\right )\)}{\(f^{\prime}(x) = \frac{2} 
{1-x^{2}} ;\ x\in \mathbb{R}\setminus \left \{-1;1\right \}\)}{\(f^{\prime}(x) = \frac{1-x} 
{1+x};\ x\in \left (-1;1\right )\)}{\(f^{\prime}(x) = \frac{1-x} 
{1+x};\ x\in \mathbb{R}\setminus \left \{-1;1\right \}\)}{}{}}
\LANGsk{
\Question{
Určte prvú deriváciu funkcie \(f\colon y =\ln \left (\frac{1+x}
{1-x}\right )\).}
{\(f^{\prime}(x) = \frac{2} 
{1-x^{2}} ;\ x\in \left (-1;1\right )\)}{\(f^{\prime}(x) = \frac{2} 
{1-x^{2}} ;\ x\in \mathbb{R}\setminus \left \{-1;1\right \}\)}{\(f^{\prime}(x) = \frac{1-x} 
{1+x};\ x\in \left (-1;1\right )\)}{\(f^{\prime}(x) = \frac{1-x} 
{1+x};\ x\in \mathbb{R}\setminus \left \{-1;1\right \}\)}{}{}}
\LANGes{
\Question{
Deriva la siguiente función.
\[ 
 f(x) =\ln \left (\frac{1 + x}
{1 - x}\right )
\]}
{\(f^{\prime}(x) = \frac{2} 
{1-x^{2}} ;\ x\in \left (-1;1\right )\)}{\(f^{\prime}(x) = \frac{2} 
{1-x^{2}} ;\ x\in \mathbb{R}\setminus \left \{-1;1\right \}\)}{\(f^{\prime}(x) = \frac{1-x} 
{1+x};\ x\in \left (-1;1\right )\)}{\(f^{\prime}(x) = \frac{1-x} 
{1+x};\ x\in \mathbb{R}\setminus \left \{-1;1\right \}\)}{}{}}
\End

\Data\ID{9000070801}
\Updated{2020-07-15 03:07:37}
\Level{B}
\SubArea{Derivative}
\LANGen{
\Question{
Differentiate the following function. 
\[ 
 f(x) = 3\sin x\cos x
\]}
{\(f'(x) = 3\cos (2x);\ x\in \mathbb{R}\)}{\(f'(x) = 3;\ x\in \mathbb{R}\)}{\(f'(x) = -3\cos x\sin x;\ x\in \mathbb{R}\)}{\(f'(x) = 3(\cos x)^{2};\ x\in \mathbb{R}\)}{}{}}
\LANGpl{
\Question{
Wyznacz pochodną funkcji.
\[ 
 f\colon y = 3\sin x\cos x
\]}
{\(f'(x) = 3\cos (2x);\ x\in \mathbb{R}\)}{\(f'(x) = 3;\ x\in \mathbb{R}\)}{\(f'(x) = -3\cos x\sin x;\ x\in \mathbb{R}\)}{\(f'(x) = 3(\cos x)^{2};\ x\in \mathbb{R}\)}{}{}}
\LANGcs{
\Question{
Určete první derivaci funkce \(f\colon y = 3\sin x\cos x\).}
{\(f'(x) = 3\cos (2x);\ x\in \mathbb{R}\)}{\(f'(x) = 3;\ x\in \mathbb{R}\)}{\(f'(x) = -3\cos x\sin x;\ x\in \mathbb{R}\)}{\(f'(x) = 3(\cos x)^{2};\ x\in \mathbb{R}\)}{}{}}
\LANGsk{
\Question{
Určte prvú deriváciu funkcie \(f\colon y = 3\sin x\cos x\).}
{\(f'(x) = 3\cos (2x);\ x\in \mathbb{R}\)}{\(f'(x) = 3;\ x\in \mathbb{R}\)}{\(f'(x) = -3\cos x\sin x;\ x\in \mathbb{R}\)}{\(f'(x) = 3(\cos x)^{2};\ x\in \mathbb{R}\)}{}{}}
\LANGes{
\Question{
Deriva la siguiente función.
\[ 
 f(x) = 3\sin x\cos x
\]}
{\(f'(x) = 3\cos (2x);\ x\in \mathbb{R}\)}{\(f'(x) = 3;\ x\in \mathbb{R}\)}{\(f'(x) = -3\cos x\sin x;\ x\in \mathbb{R}\)}{\(f'(x) = 3(\cos x)^{2};\ x\in \mathbb{R}\)}{}{}}
\End

\Data\ID{9000070807}
\Updated{2020-07-15 03:07:37}
\Level{B}
\SubArea{Derivative}
\LANGen{
\Question{
Differentiate the following function. 
\[ 
 f(x) = \frac{x^{4} + 3} 
 {x^{2}} + x^{3}
\]}
{\(f'(x) = 3x^{2} + 2x - \frac{6} 
{x^{3}} ;\ x\in \mathbb{R}\setminus \{0\}\)}{\(f'(x) = 6x^{2} - 2x - \frac{6} 
{x^{3}} ;\ x\in \mathbb{R}\setminus \{0\}\)}{\(f'(x) = 3x^{2} + 2x + \frac{6} 
{x^{3}} ;\ x\in \mathbb{R}\setminus \{0\}\)}{\(f'(x) = 6x^{2} - 2x + \frac{6} 
{x^{3}} ;\ x\in \mathbb{R}\setminus \{0\}\)}{}{}}
\LANGpl{
\Question{
Wyznacz pochodną funkcji.
\[ 
 f\colon y = \frac{x^{4} + 3} 
 {x^{2}} + x^{3}
\]}
{\(f'(x) = 3x^{2} + 2x - \frac{6} 
{x^{3}} ;\ x\in \mathbb{R}\setminus \{0\}\)}{\(f'(x) = 6x^{2} - 2x - \frac{6} 
{x^{3}} ;\ x\in \mathbb{R}\setminus \{0\}\)}{\(f'(x) = 3x^{2} + 2x + \frac{6} 
{x^{3}} ;\ x\in \mathbb{R}\setminus \{0\}\)}{\(f'(x) = 6x^{2} - 2x + \frac{6} 
{x^{3}} ;\ x\in \mathbb{R}\setminus \{0\}\)}{}{}}
\LANGcs{
\Question{
Určete první derivaci funkce \(f\colon y = \frac{x^{4}+3} 
 {x^{2}} + x^{3}\).}
{\(f'(x) = 3x^{2} + 2x - \frac{6} 
{x^{3}} ;\ x\in \mathbb{R}\setminus \{0\}\)}{\(f'(x) = 6x^{2} - 2x - \frac{6} 
{x^{3}} ;\ x\in \mathbb{R}\setminus \{0\}\)}{\(f'(x) = 3x^{2} + 2x + \frac{6} 
{x^{3}} ;\ x\in \mathbb{R}\setminus \{0\}\)}{\(f'(x) = 6x^{2} - 2x + \frac{6} 
{x^{3}} ;\ x\in \mathbb{R}\setminus \{0\}\)}{}{}}
\LANGsk{
\Question{
Určte prvú deriváciu funkcie \(f\colon y = \frac{x^{4}+3} 
 {x^{2}} + x^{3}\).}
{\(f'(x) = 3x^{2} + 2x - \frac{6} 
{x^{3}} ;\ x\in \mathbb{R}\setminus \{0\}\)}{\(f'(x) = 6x^{2} - 2x - \frac{6} 
{x^{3}} ;\ x\in \mathbb{R}\setminus \{0\}\)}{\(f'(x) = 3x^{2} + 2x + \frac{6} 
{x^{3}} ;\ x\in \mathbb{R}\setminus \{0\}\)}{\(f'(x) = 6x^{2} - 2x + \frac{6} 
{x^{3}} ;\ x\in \mathbb{R}\setminus \{0\}\)}{}{}}
\LANGes{
\Question{
Deriva la siguiente función.
\[ 
 f(x) = \frac{x^{4} + 3} 
 {x^{2}} + x^{3}
\]}
{\(f'(x) = 3x^{2} + 2x - \frac{6} 
{x^{3}} ;\ x\in \mathbb{R}\setminus \{0\}\)}{\(f'(x) = 6x^{2} - 2x - \frac{6} 
{x^{3}} ;\ x\in \mathbb{R}\setminus \{0\}\)}{\(f'(x) = 3x^{2} + 2x + \frac{6} 
{x^{3}} ;\ x\in \mathbb{R}\setminus \{0\}\)}{\(f'(x) = 6x^{2} - 2x + \frac{6} 
{x^{3}} ;\ x\in \mathbb{R}\setminus \{0\}\)}{}{}}
\End

\Data\ID{9000070808}
\Updated{2020-07-15 03:07:37}
\Level{B}
\SubArea{Derivative}
\LANGen{
\Question{
Differentiate the following function. 
\[ 
 f(x)= \frac{x} 
{x + 1}
\]}
{\(f'(x) = \frac{1} 
{(x+1)^{2}} ;\ x\in \mathbb{R}\setminus \{ - 1\}\)}{\(f'(x) = - \frac{1}
{(x+1)^{2}} ;\ x\in \mathbb{R}\setminus \{ - 1\}\)}{\(f'(x) = \frac{x} 
{(x+1)^{2}} ;\ x\in \mathbb{R}\setminus \{ - 1\}\)}{\(f'(x) = - \frac{x}
{(x+1)^{2}} ;\ x\in \mathbb{R}\setminus \{ - 1\}\)}{}{}}
\LANGpl{
\Question{
Wyznacz pochodną funkcji.
\[ 
 f\colon y = \frac{x} 
{x + 1}
\]}
{\(f'(x) = \frac{1} 
{(x+1)^{2}} ;\ x\in \mathbb{R}\setminus \{ - 1\}\)}{\(f'(x) = - \frac{1}
{(x+1)^{2}} ;\ x\in \mathbb{R}\setminus \{ - 1\}\)}{\(f'(x) = \frac{x} 
{(x+1)^{2}} ;\ x\in \mathbb{R}\setminus \{ - 1\}\)}{\(f'(x) = - \frac{x}
{(x+1)^{2}} ;\ x\in \mathbb{R}\setminus \{ - 1\}\)}{}{}}
\LANGcs{
\Question{
Určete první derivaci funkce \(f\colon y = \frac{x} 
{x+1}\).}
{\(f'(x) = \frac{1} 
{(x+1)^{2}} ;\ x\in \mathbb{R}\setminus \{ - 1\}\)}{\(f'(x) = - \frac{1}
{(x+1)^{2}} ;\ x\in \mathbb{R}\setminus \{ - 1\}\)}{\(f'(x) = \frac{x} 
{(x+1)^{2}} ;\ x\in \mathbb{R}\setminus \{ - 1\}\)}{\(f'(x) = - \frac{x}
{(x+1)^{2}} ;\ x\in \mathbb{R}\setminus \{ - 1\}\)}{}{}}
\LANGsk{
\Question{
Určte prvú deriváciu funkcie \(f\colon y = \frac{x} 
{x+1}\).}
{\(f'(x) = \frac{1} 
{(x+1)^{2}} ;\ x\in \mathbb{R}\setminus \{ - 1\}\)}{\(f'(x) = - \frac{1}
{(x+1)^{2}} ;\ x\in \mathbb{R}\setminus \{ - 1\}\)}{\(f'(x) = \frac{x} 
{(x+1)^{2}} ;\ x\in \mathbb{R}\setminus \{ - 1\}\)}{\(f'(x) = - \frac{x}
{(x+1)^{2}} ;\ x\in \mathbb{R}\setminus \{ - 1\}\)}{}{}}
\LANGes{
\Question{
Deriva la siguiente función.
\[ 
 f(x)= \frac{x} 
{x + 1}
\]}
{\(f'(x) = \frac{1} 
{(x+1)^{2}} ;\ x\in \mathbb{R}\setminus \{ - 1\}\)}{\(f'(x) = - \frac{1}
{(x+1)^{2}} ;\ x\in \mathbb{R}\setminus \{ - 1\}\)}{\(f'(x) = \frac{x} 
{(x+1)^{2}} ;\ x\in \mathbb{R}\setminus \{ - 1\}\)}{\(f'(x) = - \frac{x}
{(x+1)^{2}} ;\ x\in \mathbb{R}\setminus \{ - 1\}\)}{}{}}
\End

\Data\ID{9000070809}
\Updated{2020-07-15 03:07:37}
\Level{B}
\SubArea{Derivative}
\LANGen{
\Question{
Differentiate the following function. 
\[ 
 f(x)= 3x^{2}\sin x
\]}
{\(f'(x) = 6x\sin x + 3x^{2}\cos x;\ x\in \mathbb{R}\)}{\(f'(x) = 6x\cos x;\ x\in \mathbb{R}\)}{\(f'(x) = 3x^{2}\sin x\cos x;\ x\in \mathbb{R}\)}{\(f'(x) = -3x^{2}\sin x\cos x;\ x\in \mathbb{R}\)}{}{}}
\LANGpl{
\Question{
Wyznacz pochodną funkcji.
\[ 
 f\colon y = 3x^{2}\sin x
\]}
{\(f'(x) = 6x\sin x + 3x^{2}\cos x;\ x\in \mathbb{R}\)}{\(f'(x) = 6x\cos x;\ x\in \mathbb{R}\)}{\(f'(x) = 3x^{2}\sin x\cos x;\ x\in \mathbb{R}\)}{\(f'(x) = -3x^{2}\sin x\cos x;\ x\in \mathbb{R}\)}{}{}}
\LANGcs{
\Question{
Určete první derivaci funkce \(f\colon y = 3x^{2}\sin x\).}
{\(f'(x) = 6x\sin x + 3x^{2}\cos x;\ x\in \mathbb{R}\)}{\(f'(x) = 6x\cos x;\ x\in \mathbb{R}\)}{\(f'(x) = 3x^{2}\sin x\cos x;\ x\in \mathbb{R}\)}{\(f'(x) = -3x^{2}\sin x\cos x;\ x\in \mathbb{R}\)}{}{}}
\LANGsk{
\Question{
Určte prvú deriváciu funkcie \(f\colon y = 3x^{2}\sin x\).}
{\(f'(x) = 6x\sin x + 3x^{2}\cos x;\ x\in \mathbb{R}\)}{\(f'(x) = 6x\cos x;\ x\in \mathbb{R}\)}{\(f'(x) = 3x^{2}\sin x\cos x;\ x\in \mathbb{R}\)}{\(f'(x) = -3x^{2}\sin x\cos x;\ x\in \mathbb{R}\)}{}{}}
\LANGes{
\Question{
Deriva la siguiente función.
\[ 
 f(x)= 3x^{2}\sin x
\]}
{\(f'(x) = 6x\sin x + 3x^{2}\cos x;\ x\in \mathbb{R}\)}{\(f'(x) = 6x\cos x;\ x\in \mathbb{R}\)}{\(f'(x) = 3x^{2}\sin x\cos x;\ x\in \mathbb{R}\)}{\(f'(x) = -3x^{2}\sin x\cos x;\ x\in \mathbb{R}\)}{}{}}
\End

\Data\ID{9000063303}
\Updated{2020-07-15 03:07:37}
\Level{C}
\SubArea{Derivative}
\LANGen{
\Question{
Differentiate the following function. 
\[ 
 f(x) = \sqrt{\sin x}
\]}
{\(f'(x) = \frac{\cos x} 
{2\sqrt{\sin x}},\ x\in \mathop{\mathop{\bigcup 
 }}\nolimits _{k\in \mathbb{Z}}\left (2k\pi ;\pi + 2k\pi \right )\)}{\(f'(x) = \frac{\sin x} 
{2\sqrt{\cos x}},\ x\in \mathop{\mathop{\bigcup 
 }}\nolimits _{k\in \mathbb{Z}}\left (2k\pi ; \frac{\pi } {2} + 2k\pi \right )\)}{\(f'(x) = \frac{1} 
{2\sqrt{\sin x}},\ x\in \mathop{\mathop{\bigcup 
 }}\nolimits _{k\in \mathbb{Z}}\left (2k\pi ;\pi + 2k\pi \right )\)}{\(f'(x) = \frac{\cos x} 
{2\sqrt{\sin x}},\ x\in \mathop{\mathop{\bigcup 
 }}\nolimits _{k\in \mathbb{Z}}\left [ 2k\pi ; \frac{\pi } {2} + 2k\pi \right ] \)}{}{}}
\LANGpl{
\Question{
Wyznacz pochodną funkcji. 
\[ 
 f\colon y = \sqrt{\sin x}
\]}
{\(f'(x) = \frac{\cos x} 
{2\sqrt{\sin x}},\ x\in \mathop{\mathop{\bigcup 
 }}\nolimits _{k\in \mathbb{Z}}\left (2k\pi ;\pi + 2k\pi \right )\)}{\(f'(x) = \frac{\sin x} 
{2\sqrt{\cos x}},\ x\in \mathop{\mathop{\bigcup 
 }}\nolimits _{k\in \mathbb{Z}}\left (2k\pi ; \frac{\pi } {2} + 2k\pi \right )\)}{\(f'(x) = \frac{1} 
{2\sqrt{\sin x}},\ x\in \mathop{\mathop{\bigcup 
 }}\nolimits _{k\in \mathbb{Z}}\left (2k\pi ;\pi + 2k\pi \right )\)}{\(f'(x) = \frac{\cos x} 
{2\sqrt{\sin x}},\ x\in \mathop{\mathop{\bigcup 
 }}\nolimits _{k\in \mathbb{Z}}\left [ 2k\pi ; \frac{\pi } {2} + 2k\pi \right ] \)}{}{}}
\LANGcs{
\Question{
Derivace funkce \(f\colon y = \sqrt{\sin x}\) 
je rovna:}
{\(f'(x) = \frac{\cos x} 
{2\sqrt{\sin x}},\ x\in \mathop{\mathop{\bigcup 
 }}\nolimits _{k\in \mathbb{Z}}\left (2k\pi ;\pi + 2k\pi \right )\)}{\(f'(x) = \frac{\sin x} 
{2\sqrt{\cos x}},\ x\in \mathop{\mathop{\bigcup 
 }}\nolimits _{k\in \mathbb{Z}}\left (2k\pi ; \frac{\pi } {2} + 2k\pi \right )\)}{\(f'(x) = \frac{1} 
{2\sqrt{\sin x}},\ x\in \mathop{\mathop{\bigcup 
 }}\nolimits _{k\in \mathbb{Z}}\left (2k\pi ;\pi + 2k\pi \right )\)}{\(f'(x) = \frac{\cos x} 
{2\sqrt{\sin x}},\ x\in \mathop{\mathop{\bigcup 
 }}\nolimits _{k\in \mathbb{Z}}\left \langle 2k\pi ; \frac{\pi } {2} + 2k\pi \right \rangle \)}{}{}}
\LANGsk{
\Question{
Derivácia funkcie \(f\colon y = \sqrt{\sin x}\) 
je rovná:}
{\(f'(x) = \frac{\cos x} 
{2\sqrt{\sin x}},\ x\in \mathop{\mathop{\bigcup 
 }}\nolimits _{k\in \mathbb{Z}}\left (2k\pi ;\pi + 2k\pi \right )\)}{\(f'(x) = \frac{\sin x} 
{2\sqrt{\cos x}},\ x\in \mathop{\mathop{\bigcup 
 }}\nolimits _{k\in \mathbb{Z}}\left (2k\pi ; \frac{\pi } {2} + 2k\pi \right )\)}{\(f'(x) = \frac{1} 
{2\sqrt{\sin x}},\ x\in \mathop{\mathop{\bigcup 
 }}\nolimits _{k\in \mathbb{Z}}\left (2k\pi ;\pi + 2k\pi \right )\)}{\(f'(x) = \frac{\cos x} 
{2\sqrt{\sin x}},\ x\in \mathop{\mathop{\bigcup 
 }}\nolimits _{k\in \mathbb{Z}}\left \langle 2k\pi ; \frac{\pi } {2} + 2k\pi \right \rangle \)}{}{}}
\LANGes{
\Question{
Deriva la siguiente función.
\[ 
 f(x) = \sqrt{\sin x}
\]}
{\(f'(x) = \frac{\cos x} 
{2\sqrt{\sin x}},\ x\in \mathop{\mathop{\bigcup 
 }}\nolimits _{k\in \mathbb{Z}}\left (2k\pi ;\pi + 2k\pi \right )\)}{\(f'(x) = \frac{\sin x} 
{2\sqrt{\cos x}},\ x\in \mathop{\mathop{\bigcup 
 }}\nolimits _{k\in \mathbb{Z}}\left (2k\pi ; \frac{\pi } {2} + 2k\pi \right )\)}{\(f'(x) = \frac{1} 
{2\sqrt{\sin x}},\ x\in \mathop{\mathop{\bigcup 
 }}\nolimits _{k\in \mathbb{Z}}\left (2k\pi ;\pi + 2k\pi \right )\)}{\(f'(x) = \frac{\cos x} 
{2\sqrt{\sin x}},\ x\in \mathop{\mathop{\bigcup 
 }}\nolimits _{k\in \mathbb{Z}}\left [ 2k\pi ; \frac{\pi } {2} + 2k\pi \right ] \)}{}{}}
\End

\Data\ID{9000063304}
\Updated{2020-07-15 03:07:37}
\Level{C}
\SubArea{Derivative}
\LANGen{
\Question{
Differentiate the following function. 
\[ 
 f(x) =\ln \sqrt{x}
\]}
{\(f'(x) = \frac{1} 
{2x},\ x > 0\)}{\(f'(x) = \frac{1} 
{2x},\ x\neq 0\)}{\(f'(x) = \frac{1} 
{x},\ x > 0\)}{\(f'(x) = \frac{1} 
{x},\ x\neq 0\)}{}{}}
\LANGpl{
\Question{
Wyznacz pochodną funkcji.  
\[ 
 f\colon y =\ln \sqrt{x}
\]}
{\(f'(x) = \frac{1} 
{2x},\ x > 0\)}{\(f'(x) = \frac{1} 
{2x},\ x\neq 0\)}{\(f'(x) = \frac{1} 
{x},\ x > 0\)}{\(f'(x) = \frac{1} 
{x},\ x\neq 0\)}{}{}}
\LANGcs{
\Question{
Derivace funkce \(f\colon y =\ln \sqrt{x}\) 
je rovna:}
{\(f'(x) = \frac{1} 
{2x},\ x > 0\)}{\(f'(x) = \frac{1} 
{2x},\ x\neq 0\)}{\(f'(x) = \frac{1} 
{x},\ x > 0\)}{\(f'(x) = \frac{1} 
{x},\ x\neq 0\)}{}{}}
\LANGsk{
\Question{
Derivácia funkcie \(f\colon y =\ln \sqrt{x}\) 
je rovná:}
{\(f'(x) = \frac{1} 
{2x},\ x > 0\)}{\(f'(x) = \frac{1} 
{2x},\ x\neq 0\)}{\(f'(x) = \frac{1} 
{x},\ x > 0\)}{\(f'(x) = \frac{1} 
{x},\ x\neq 0\)}{}{}}
\LANGes{
\Question{
Deriva la siguiente función.
\[ 
 f(x) =\ln \sqrt{x}
\]}
{\(f'(x) = \frac{1} 
{2x},\ x > 0\)}{\(f'(x) = \frac{1} 
{2x},\ x\neq 0\)}{\(f'(x) = \frac{1} 
{x},\ x > 0\)}{\(f'(x) = \frac{1} 
{x},\ x\neq 0\)}{}{}}
\End

\Data\ID{9000063305}
\Updated{2020-07-15 03:07:37}
\Level{C}
\SubArea{Derivative}
\LANGen{
\Question{
Differentiate the following function. 
\[ 
f(x) = \sqrt{\frac{x - 1}
{x + 1}}
\]}
{\(f'(x) = \frac{1} 
{(x+1)^{2}} \sqrt{\frac{x+1}
{x-1}},\ x\in (-\infty ;-1)\cup (1;\infty )\)}{\(f'(x) = \frac{\sqrt{x-1}} 
{(x-1)^{2}\sqrt{x+1}},\ x\in (-\infty ;-1)\cup [ 1;\infty )\)}{\(f'(x) = \frac{x-1} 
{2\sqrt{(x+1)^{3}}} ,\ x\neq - 1\)}{\(f'(x) = \frac{x-1} 
{\sqrt{(x+1)^{3}}} ,\ x\in (-\infty ;-1)\cup [ 1;\infty )\)}{}{}}
\LANGpl{
\Question{
Wyznacz pochodną funkcji. 
\[ 
 f\colon y = \sqrt{\frac{x - 1}
{x + 1}}
\]}
{\(f'(x) = \frac{1} 
{(x+1)^{2}} \sqrt{\frac{x+1}
{x-1}},\ x\in (-\infty ;-1)\cup (1;\infty )\)}{\(f'(x) = \frac{\sqrt{x-1}} 
{(x-1)^{2}\sqrt{x+1}},\ x\in (-\infty ;-1)\cup [ 1;\infty )\)}{\(f'(x) = \frac{x-1} 
{2\sqrt{(x+1)^{3}}} ,\ x\neq - 1\)}{\(f'(x) = \frac{x-1} 
{\sqrt{(x+1)^{3}}} ,\ x\in (-\infty ;-1)\cup [ 1;\infty )\)}{}{}}
\LANGcs{
\Question{
Derivace funkce \(f\colon y = \sqrt{\frac{x-1}
{x+1}}\) 
je rovna:}
{\(f'(x) = \frac{1} 
{(x+1)^{2}} \sqrt{\frac{x+1}
{x-1}},\ x\in (-\infty ;-1)\cup (1;\infty )\)}{\(f'(x) = \frac{\sqrt{x-1}} 
{(x-1)^{2}\sqrt{x+1}},\ x\in (-\infty ;-1)\cup \langle 1;\infty )\)}{\(f'(x) = \frac{x-1} 
{2\sqrt{(x+1)^{3}}} ,\ x\neq - 1\)}{\(f'(x) = \frac{x-1} 
{\sqrt{(x+1)^{3}}} ,\ x\in (-\infty ;-1)\cup \langle 1;\infty )\)}{}{}}
\LANGsk{
\Question{
Derivácia funkcie \(f\colon y = \sqrt{\frac{x-1}
{x+1}}\) 
je rovná:}
{\(f'(x) = \frac{1} 
{(x+1)^{2}} \sqrt{\frac{x+1}
{x-1}},\ x\in (-\infty ;-1)\cup (1;\infty )\)}{\(f'(x) = \frac{\sqrt{x-1}} 
{(x-1)^{2}\sqrt{x+1}},\ x\in (-\infty ;-1)\cup \langle 1;\infty )\)}{\(f'(x) = \frac{x-1} 
{2\sqrt{(x+1)^{3}}} ,\ x\neq - 1\)}{\(f'(x) = \frac{x-1} 
{\sqrt{(x+1)^{3}}} ,\ x\in (-\infty ;-1)\cup \langle 1;\infty )\)}{}{}}
\LANGes{
\Question{
Deriva la siguiente función.
\[ 
f(x) = \sqrt{\frac{x - 1}
{x + 1}}
\]}
{\(f'(x) = \frac{1} 
{(x+1)^{2}} \sqrt{\frac{x+1}
{x-1}},\ x\in (-\infty ;-1)\cup (1;\infty )\)}{\(f'(x) = \frac{\sqrt{x-1}} 
{(x-1)^{2}\sqrt{x+1}},\ x\in (-\infty ;-1)\cup [ 1;\infty )\)}{\(f'(x) = \frac{x-1} 
{2\sqrt{(x+1)^{3}}} ,\ x\neq - 1\)}{\(f'(x) = \frac{x-1} 
{\sqrt{(x+1)^{3}}} ,\ x\in (-\infty ;-1)\cup [ 1;\infty )\)}{}{}}
\End

\Data\ID{9000063306}
\Updated{2020-07-15 03:07:37}
\Level{C}
\SubArea{Derivative}
\LANGen{
\Question{
Differentiate the following function. 
\[ 
 f(x) =\mathrm{e} ^{\sin 2x}
\]}
{\(f'(x) = 2\mathrm{e}^{\sin 2x}\cos 2x,\ x\in \mathbb{R}\)}{\(f'(x) = x\mathrm{e}^{\sin 2x}\cos 2x,\ x\in \mathbb{R}\)}{\(f'(x) =\mathrm{e} ^{\sin 2x}\sin 2x,\ x\in \mathbb{R}\)}{\(f'(x) =\mathrm{e} ^{\cos 2x},\ x\in \mathbb{R}\)}{}{}}
\LANGpl{
\Question{
Wyznacz pochodną funkcji. 
\[ 
 f\colon y =\mathrm{e} ^{\sin 2x}
\]}
{\(f'(x) = 2\mathrm{e}^{\sin 2x}\cos 2x,\ x\in \mathbb{R}\)}{\(f'(x) = x\mathrm{e}^{\sin 2x}\cos 2x,\ x\in \mathbb{R}\)}{\(f'(x) =\mathrm{e} ^{\sin 2x}\sin 2x,\ x\in \mathbb{R}\)}{\(f'(x) =\mathrm{e} ^{\cos 2x},\ x\in \mathbb{R}\)}{}{}}
\LANGcs{
\Question{
Derivace funkce \(f\colon y =\mathrm{e} ^{\sin 2x}\) 
je rovna:}
{\(f'(x) = 2\mathrm{e}^{\sin 2x}\cos 2x,\ x\in \mathbb{R}\)}{\(f'(x) = x\mathrm{e}^{\sin 2x}\cos 2x,\ x\in \mathbb{R}\)}{\(f'(x) =\mathrm{e} ^{\sin 2x}\sin 2x,\ x\in \mathbb{R}\)}{\(f'(x) =\mathrm{e} ^{\cos 2x},\ x\in \mathbb{R}\)}{}{}}
\LANGsk{
\Question{
Derivácie funkcie \(f\colon y =\mathrm{e} ^{\sin 2x}\) 
je rovná:}
{\(f'(x) = 2\mathrm{e}^{\sin 2x}\cos 2x,\ x\in \mathbb{R}\)}{\(f'(x) = x\mathrm{e}^{\sin 2x}\cos 2x,\ x\in \mathbb{R}\)}{\(f'(x) =\mathrm{e} ^{\sin 2x}\sin 2x,\ x\in \mathbb{R}\)}{\(f'(x) =\mathrm{e} ^{\cos 2x},\ x\in \mathbb{R}\)}{}{}}
\LANGes{
\Question{
Deriva la siguiente función.
\[ 
 f(x) =\mathrm{e} ^{\sin 2x}
\]}
{\(f'(x) = 2\mathrm{e}^{\sin 2x}\cos 2x,\ x\in \mathbb{R}\)}{\(f'(x) = x\mathrm{e}^{\sin 2x}\cos 2x,\ x\in \mathbb{R}\)}{\(f'(x) =\mathrm{e} ^{\sin 2x}\sin 2x,\ x\in \mathbb{R}\)}{\(f'(x) =\mathrm{e} ^{\cos 2x},\ x\in \mathbb{R}\)}{}{}}
\End

\Data\ID{9000063307}
\Updated{2020-07-15 03:07:37}
\Level{C}
\SubArea{Derivative}
\LANGen{
\Question{
Differentiate the following function. 
\[ 
 f(x) =\ln\left(\cos 2x\right)
\]}
{\(f'(x) = -2\mathop{\mathrm{tg}}\nolimits 2x,\ x\in \mathop{\mathop{\bigcup 
 }}\nolimits _{k\in \mathbb{Z}}\left (-\frac{\pi }{4} + k\pi ; \frac{\pi } {4} + k\pi \right )\)}{\(f'(x) = 2\mathop{\mathrm{tg}}\nolimits 2x,\ x\in \mathop{\mathop{\bigcup 
 }}\nolimits _{k\in \mathbb{Z}}\left (-\frac{\pi }{4} + k\pi ; \frac{\pi } {4} + k\pi \right )\)}{\(f'(x) = -2,\ x\in \mathop{\mathop{\bigcup 
 }}\nolimits _{k\in \mathbb{Z}}\left (-\frac{\pi }{4} + k\pi ; \frac{\pi } {4} + k\pi \right )\)}{\(f'(x) = 1 -\ln\left(\sin 2x\right),\ x\in \mathop{\mathop{\bigcup 
 }}\nolimits _{k\in \mathbb{Z}}\left (k\pi ; \frac{\pi } {2} + k\pi \right )\)}{}{}}
\LANGpl{
\Question{
Wyznacz pochodną funkcji. 
\[ 
 f\colon y =\ln\left(\cos 2x\right)
\]}
{\(f'(x) = -2\mathop{\mathrm{tg}}\nolimits 2x,\ x\in \mathop{\mathop{\bigcup 
 }}\nolimits _{k\in \mathbb{Z}}\left (-\frac{\pi }{4} + k\pi ; \frac{\pi } {4} + k\pi \right )\)}{\(f'(x) = 2\mathop{\mathrm{tg}}\nolimits 2x,\ x\in \mathop{\mathop{\bigcup 
 }}\nolimits _{k\in \mathbb{Z}}\left (-\frac{\pi }{4} + k\pi ; \frac{\pi } {4} + k\pi \right )\)}{\(f'(x) = -2,\ x\in \mathop{\mathop{\bigcup 
 }}\nolimits _{k\in \mathbb{Z}}\left (-\frac{\pi }{4} + k\pi ; \frac{\pi } {4} + k\pi \right )\)}{\(f'(x) = 1 -\ln\left(\sin 2x\right),\ x\in \mathop{\mathop{\bigcup 
 }}\nolimits _{k\in \mathbb{Z}}\left (k\pi ; \frac{\pi } {2} + k\pi \right )\)}{}{}}
\LANGcs{
\Question{
Derivace funkce \(f\colon y =\ln \left(\cos 2x\right)\) 
je rovna:}
{\(f'(x) = -2\mathop{\mathrm{tg}}\nolimits 2x,\ x\in \mathop{\mathop{\bigcup 
 }}\nolimits _{k\in \mathbb{Z}}\left (-\frac{\pi }{4} + k\pi ; \frac{\pi } {4} + k\pi \right )\)}{\(f'(x) = 2\mathop{\mathrm{tg}}\nolimits 2x,\ x\in \mathop{\mathop{\bigcup 
 }}\nolimits _{k\in \mathbb{Z}}\left (-\frac{\pi }{4} + k\pi ; \frac{\pi } {4} + k\pi \right )\)}{\(f'(x) = -2,\ x\in \mathop{\mathop{\bigcup 
 }}\nolimits _{k\in \mathbb{Z}}\left (-\frac{\pi }{4} + k\pi ; \frac{\pi } {4} + k\pi \right )\)}{\(f'(x) = 1 -\ln\left(\sin 2x\right),\ x\in \mathop{\mathop{\bigcup 
 }}\nolimits _{k\in \mathbb{Z}}\left (k\pi ; \frac{\pi } {2} + k\pi \right )\)}{}{}}
\LANGsk{
\Question{
Derivácia funkcie \(f\colon y =\ln \left(\cos 2x\right)\) 
je rovná:}
{\(f'(x) = -2\mathop{\mathrm{tg}}\nolimits 2x,\ x\in \mathop{\mathop{\bigcup 
 }}\nolimits _{k\in \mathbb{Z}}\left (-\frac{\pi }{4} + k\pi ; \frac{\pi } {4} + k\pi \right )\)}{\(f'(x) = 2\mathop{\mathrm{tg}}\nolimits 2x,\ x\in \mathop{\mathop{\bigcup 
 }}\nolimits _{k\in \mathbb{Z}}\left (-\frac{\pi }{4} + k\pi ; \frac{\pi } {4} + k\pi \right )\)}{\(f'(x) = -2,\ x\in \mathop{\mathop{\bigcup 
 }}\nolimits _{k\in \mathbb{Z}}\left (-\frac{\pi }{4} + k\pi ; \frac{\pi } {4} + k\pi \right )\)}{\(f'(x) = 1 -\ln\left(\sin 2x\right),\ x\in \mathop{\mathop{\bigcup 
 }}\nolimits _{k\in \mathbb{Z}}\left (k\pi ; \frac{\pi } {2} + k\pi \right )\)}{}{}}
\LANGes{
\Question{
Deriva la siguiente función.
\[ 
 f(x) =\ln\left(\cos 2x\right)
\]}
{\(f'(x) = -2\mathop{\mathrm{tg}}\nolimits 2x,\ x\in \mathop{\mathop{\bigcup 
 }}\nolimits _{k\in \mathbb{Z}}\left (-\frac{\pi }{4} + k\pi ; \frac{\pi } {4} + k\pi \right )\)}{\(f'(x) = 2\mathop{\mathrm{tg}}\nolimits 2x,\ x\in \mathop{\mathop{\bigcup 
 }}\nolimits _{k\in \mathbb{Z}}\left (-\frac{\pi }{4} + k\pi ; \frac{\pi } {4} + k\pi \right )\)}{\(f'(x) = -2,\ x\in \mathop{\mathop{\bigcup 
 }}\nolimits _{k\in \mathbb{Z}}\left (-\frac{\pi }{4} + k\pi ; \frac{\pi } {4} + k\pi \right )\)}{\(f'(x) = 1 -\ln\left(\sin 2x\right),\ x\in \mathop{\mathop{\bigcup 
 }}\nolimits _{k\in \mathbb{Z}}\left (k\pi ; \frac{\pi } {2} + k\pi \right )\)}{}{}}
\End

\Data\ID{9000063308}
\Updated{2020-07-15 03:07:37}
\Level{C}
\SubArea{Derivative}
\LANGen{
\Question{
Differentiate the following function. 
\[ 
 f(x) = \frac{1} 
{\ln x}
\]}
{\(f'(x) = \frac{-1} 
{x\ln ^{2}x},\ x > 0,\ x\neq 1\)}{\(f'(x) =\ln x,\ x > 0\)}{\(f'(x) = \frac{\ln x} 
{x},\ x\neq 0\)}{\(f'(x) = \frac{1} 
{x\ln ^{2}x},\ x\neq 0\)}{}{}}
\LANGpl{
\Question{
Wyznacz pochodną funkcji. 
\[ 
 f\colon y = \frac{1} 
{\ln x}
\]}
{\(f'(x) = \frac{-1} 
{x\ln ^{2}x},\ x > 0,\ x\neq 1\)}{\(f'(x) =\ln x,\ x > 0\)}{\(f'(x) = \frac{\ln x} 
{x},\ x\neq 0\)}{\(f'(x) = \frac{1} 
{x\ln ^{2}x},\ x\neq 0\)}{}{}}
\LANGcs{
\Question{
Derivace funkce \(f\colon y = \frac{1} 
{\ln x}\) 
je rovna:}
{\(f'(x) = \frac{-1} 
{x\ln ^{2}x},\ x > 0,\ x\neq 1\)}{\(f'(x) =\ln x,\ x > 0\)}{\(f'(x) = \frac{\ln x} 
{x},\ x\neq 0\)}{\(f'(x) = \frac{1} 
{x\ln ^{2}x},\ x\neq 0\)}{}{}}
\LANGsk{
\Question{
Derivácia funkcie \(f\colon y = \frac{1} 
{\ln x}\) 
je rovná:}
{\(f'(x) = \frac{-1} 
{x\ln ^{2}x},\ x > 0,\ x\neq 1\)}{\(f'(x) =\ln x,\ x > 0\)}{\(f'(x) = \frac{\ln x} 
{x},\ x\neq 0\)}{\(f'(x) = \frac{1} 
{x\ln ^{2}x},\ x\neq 0\)}{}{}}
\LANGes{
\Question{
Deriva la siguiente función.
\[ 
 f(x) = \frac{1} 
{\ln x}
\]}
{\(f'(x) = \frac{-1} 
{x\ln ^{2}x},\ x > 0,\ x\neq 1\)}{\(f'(x) =\ln x,\ x > 0\)}{\(f'(x) = \frac{\ln x} 
{x},\ x\neq 0\)}{\(f'(x) = \frac{1} 
{x\ln ^{2}x},\ x\neq 0\)}{}{}}
\End
